% arXiv submission template for pandoc (xelatex)
% Usage: pandoc paper.md --template=pandoc.format.arxiv.tex --pdf-engine=xelatex -o paper.tex

\documentclass[11pt,a4paper]{article}

%%%%%%%%%%%%%%%%%%%%%%%%%%%%%%%
% Core packages for xelatex
%%%%%%%%%%%%%%%%%%%%%%%%%%%%%%%
\usepackage{fontspec}
%\usepackage{unicode-math}
\usepackage[titletoc,title,header]{appendix}

%%%%%%%%%%%%%%%%%%%%%%%%%%%%%%%
% AMS math packages (arXiv standard)
% Note: amssymb is omitted because unicode-math provides all symbols
%%%%%%%%%%%%%%%%%%%%%%%%%%%%%%%
\usepackage{amsmath}
\usepackage{amssymb}
\usepackage{amsthm}
\usepackage{mathtools}
\usepackage{xparse}    % for \NewDocumentCommand

%%%%%%%%%%%%%%%%%%%%%%%%%%%%%%%
% Page layout
%%%%%%%%%%%%%%%%%%%%%%%%%%%%%%%
\usepackage[margin=1in]{geometry}

%%%%%%%%%%%%%%%%%%%%%%%%%%%%%%%
% Graphics
%%%%%%%%%%%%%%%%%%%%%%%%%%%%%%%
\usepackage{graphicx}
\usepackage{xcolor}

%%%%%%%%%%%%%%%%%%%%%%%%%%%%%%%
% Hyperlinks
%%%%%%%%%%%%%%%%%%%%%%%%%%%%%%%
\usepackage{hyperref}
\hypersetup{
  colorlinks=true,
  linkcolor=blue,
  citecolor=blue,
  urlcolor=blue,
  pdfauthor={Yoshitsugu Sekine},
  pdftitle={Interacting Lattice Bosons on the Concrete Buchholz Algebra}
}

%%%%%%%%%%%%%%%%%%%%%%%%%%%%%%%
% Bibliography
%%%%%%%%%%%%%%%%%%%%%%%%%%%%%%%
\usepackage[numbers]{natbib}
\bibliographystyle{plain}

%%%%%%%%%%%%%%%%%%%%%%%%%%%%%%%
% Theorem environments
%%%%%%%%%%%%%%%%%%%%%%%%%%%%%%%
\theoremstyle{plain}
\newtheorem{thm}{Theorem}[section]
\AtEndEnvironment{thm}{\qed}
\newtheorem{lem}[thm]{Lemma}
\AtEndEnvironment{lem}{\qed}
\newtheorem{prop}[thm]{Proposition}
\AtEndEnvironment{prop}{\qed}
\newtheorem{cor}[thm]{Corollary}
\AtEndEnvironment{cor}{\qed}
\newtheorem{conj}[thm]{Conjecture}
\AtEndEnvironment{conj}{\qed}
\newtheorem{fact}[thm]{Fact}
\AtEndEnvironment{fact}{\qed}

\theoremstyle{definition}
\newtheorem{defn}[thm]{Definition}
\AtEndEnvironment{defn}{\qed}
\newtheorem{ex}[thm]{Example}
\AtEndEnvironment{ex}{\qed}

\theoremstyle{remark}
\newtheorem{rem}[thm]{Remark}
\AtEndEnvironment{rem}{\qed}

%%%%%%%%%%%%%%%%%%%%%%%%%%%%%%%
% Custom commands file
%%%%%%%%%%%%%%%%%%%%%%%%%%%%%%%
% 分野の略称はarxivに準じる
% 以前のようなショートカットがほしければyasnippetを使おう
% 非ソート対象
\makeatletter
\providecommand*{\dashv}{\mathrel{\mathpalette\@Dashv\vDash}}
\newcommand*{\@dashv}[2]{\reflectbox{$\m@th#1#2$}}
\makeatother
\newcommand{\coloredtextbf}[1]{\textbf{\textcolor{blue}{#1}}} % 文字の強調
\newcommand{\pcoloredtextbf}[1]{\textbf{\textcolor{blue}{(#1)}}} % 強調の文字装飾
%%%%%%%%%%
% 絶対値 %
%%%%%%%%%%
\newcommand{\abscard}[1]{\abs{#1}} % 絶対値記号による濃度の記号
\newcommand{\absleft}[1]{\left| #1 \right.} % 絶対値の長い式で左だけ絶対値をつけるコマンド
\newcommand{\absright}[1]{\left. #1 \right|} % 絶対値の長い式で右だけ絶対値をつけるコマンド
\newcommand{\abssuff}[2]{\abs{#1}_{#2}} % 右下への添字つき絶対値
\newcommand{\absvol}[1]{\abs{#1}} % 体積
\newcommand{\abs}[1]{\left| #1 \right|} % 絶対値
\newcommand{\pabs}[1]{\left|#1\right|_p} % L^p系の絶対値
%%%%%%%%%%
% ノルム %
%%%%%%%%%%
\NewDocumentCommand{\weaknorm}{O{\dbk} m}{#1{#2}} % 弱(準)ノルム
\newcommand{\algnorm}{\operatorname{N}} % 分離拡大のノルム, expedition0005848
\newcommand{\boxnorm}[1]{\norm{#1}_{\txtbox}} % 箱ノルム
\newcommand{\inftynorm}[1]{\norm{#1}_{\infty}} % infty-ノルム
\newcommand{\normleft}[1]{\left\Vert #1 \right.} % 左だけあるノルム
\newcommand{\normright}[1]{\left. #1 \right\Vert} % 左だけあるノルム
\newcommand{\normsuff}[2]{\norm{#1}_{#2}} % 添字組み込みのノルム
\newcommand{\norm}[1]{\left\Vert #1 \right\Vert} % ノルム
\newcommand{\onenorm}[1]{\norm{#1}_1} % 1-ノルム
\newcommand{\opnorm}[1]{\left\Vert #1 \right\Vert_{\txtoperator}} % 作用素ノルム
\newcommand{\pnorm}[1]{\left\Vert #1 \right\Vert_p} % pノルム
\newcommand{\qnorm}[1]{\left\Vert #1 \right\Vert_q} % q-ノルム
\newcommand{\rnorm}[1]{\left\Vert #1 \right\Vert_r} % r-ノルム
\newcommand{\seminormlr}[1]{\lrbk{#1}} % lrbkによるセミノルム
\newcommand{\tnorm}[1]{\left\lvert\hspace{-1 pt}\left\lvert\hspace{-1 pt}\left\lvert {#1}\right\lvert\hspace{-1 pt}\right\lvert\hspace{-1 pt}\right\lvert} % 三本線のノルム
\newcommand{\twonorm}[1]{\norm{#1}_2}
%%%%%%%%
% 内積 %
%%%%%%%%
\newcommand{\bkt}[2]{\left\langle #1,\,#2 \right\rangle} % 内積
\newcommand{\braket}[2]{\left\langle #1 \middle| #2 \right\rangle} % 物理記法でのディラックのブラケット
\newcommand{\dbkt}[3]{\left\langle #1 \middle| #2 \middle| #3 \right\rangle} % 標準的な物理記法のディラックのブラケット, 3つの引数を取る
\newcommand{\dualbkt}[2]{\bkt{#1}{#2}_{\txtdual}} % 双対内積
\newcommand{\rbktvert}[2]{\left( #1 \vert #2 \right)} % 括弧類, 内積の記号の一つ
\newcommand{\rbkt}[2]{\left( #1,\,#2 \right)} % 内積に利用
%%%%%%%%
% 極限 %
%%%%%%%%
\newcommand{\Ltwolim}{L^2 \hyphen \lim} % L^2-極限, l.i.m
\newcommand{\slim}{\mathrm{s} \hyphen \lim} % 強極限
\newcommand{\uwlim}{\mathrm{uw} \hyphen \lim} % 超弱極限
\newcommand{\vhlim}{\txtvanhowe \hyphen \lim} % ファン・ホーフェ極限
\newcommand{\wlim}{\mathrm{w} \hyphen \lim} % 弱極限
%%%%%%%%
% 矢印 %
%%%%%%%%
\newcommand{\epito}{\twoheadrightarrow} % エピ射・全射用のエイリアス
\newcommand{\isomto}{\mathrel{\rightarrowtail\kern-1.9ex\twoheadrightarrow}} % 全単射, 同型射, ホモロジー代数・圏論でよく使う
\newcommand{\monoto}{\rightarrowtail} % モノ射・単射のエイリアス
\newcommand{\sto}{\xrightarrow{\mathrm{s}}} % 強極限
\newcommand{\uto}{\xrightarrow{\mathrm{u}}} % (広義)一様収束, expedition0001858
\newcommand{\uwto}{\xrightarrow{\mathrm{uw}}} % 超弱収束
\newcommand{\wto}{\xrightarrow{\mathrm{w}}} % 弱極限
%%%%%%%%%%
% 括弧類 %
%%%%%%%%%%
\newcommand{\basebk}[1]{\left\langle #1 \right\rangle} % 基底用の括弧
\newcommand{\cbkleft}[1]{\left\{ #1 \right.} % 片側括弧
\newcommand{\cbkright}[1]{\left. #1 \right\}} % 片側括弧
\newcommand{\cbk}[1]{\left\{ #1 \right\}} % 中括弧
\newcommand{\dbk}[1]{\left\langle #1 \right\rangle} % <f>形式の括弧, Dirac bracketの意味でdbk
\newcommand{\lrbk}[1]{\left \llbracket #1 \right \rrbracket} % 同値類の記号
\newcommand{\pairbk}[1]{\rbk{#1}} % 適当な対のための括弧
\newcommand{\rbkleft}[1]{\left( #1 \right.} % 片側括弧
\newcommand{\rbkright}[1]{\left. #1 \right)} % T片側括弧
\newcommand{\rbk}[1]{\left( #1 \right)} % 丸括弧, 共通
\newcommand{\sqbkleft}[1]{\left[ #1 \right.} % sqbkに対する片側括弧
\newcommand{\sqbkright}[1]{\left. #1 \right]} % sqbkに対する片側括弧
\newcommand{\sqbk}[1]{\left[ #1 \right]} % いわゆる大括弧
\newcommand{\sqsqbk}[1]{\sqbk{\sqbk{#1}}} % 大括弧を重ねる
\newcommand{\vecbk}[1]{\rbk{#1}} % ベクトル用の括弧
%%%%%%%%%%
% 関数類 %
%%%%%%%%%%
\newcommand{\funcond}[3]{\fun{#1}{#2 \middle| #3}} % 条件つき確率などに利用する関数
\newcommand{\funpipevert}[4]{\fun{#1}{#2 \middle| #3 \middle\Vert #4}} % 相対レニーエントロピーに利用
\newcommand{\funrbk}[2]{\fun{\rbk{#1}}{#2}} % 関数の部分に丸括弧を入れてまとめる
\newcommand{\funvert}[3]{\fun{#1}{#2 \middle\Vert #3}} % カルバック-ライブラ情報量などに利用
\newcommand{\fun}[2]{#1 \rbk{#2}} % 関数テンプレート
\newcommand{\invfun}[2]{\fun{\inv{#1}}{#2}} % 逆写像
\newcommand{\sqfuncond}[3]{\sqfun{#1}{#2 \middle| #3}} % 角括弧を使う条件つき期待値などに利用する関数
\newcommand{\sqfun}[2]{#1 \sqbk{#2}} % 角括弧を使う関数, 確率論の期待値や完全な熱力学関数に利用
\newcommand{\sqfunrbk}[2]{\sqfun{\rbk{#1}}{#2}} % 角括弧を使う関数の関数の部分に丸括弧を入れてまとめる
%%%%%%%%
% 区間 %
%%%%%%%%
\newcommand{\closedinterval}[2]{\sqbk{#1,\,#2}} % 閉区間
\newcommand{\intint}[1]{\sqbk{#1}} % 整数の区間：\sqbk{n..m} で \setone{n,n+1,\cdots,m} を表す.
\newcommand{\leftopeninterval}[2]{\left(#1, \, #2 \right]} % 左開区間
\newcommand{\nonneginterval}{\rightopeninterval{0}{\infty}} % 非負の実数の全体, 区間表記したいときのfldrealnn
\newcommand{\nonposinterval}{\leftopeninterval{-\infty}{0}} % 非正の実数の全体, 区間表記したいときのfldrealnp
\newcommand{\openinterval}[2]{\rbk{#1,\,#2}} % 開区間
\newcommand{\negativeinterval}{\openinterval{-\infty}{0}} % 厳密に負の実数の全体, 区間表記したいときのfldrealneg
\newcommand{\positiveinterval}{\openinterval{0}{\infty}} % 厳密に正の実数の全体, 区間表記したいときのfldrealplus
\newcommand{\rightopeninterval}[2]{\left[#1, \, #2 \right)} % 右開区間
\newcommand{\unitclosedinterval}{\closedinterval{0}{1}} % 単位閉区間
\newcommand{\unitopeninterval}{\openinterval{0}{1}} % 単位開区間
%%%%%%%%%%
% 交換子 %
%%%%%%%%%%
\newcommand{\anticommutator}[2]{\cbk{#1,\,#2}} % 反交換子
\newcommand{\commutator}[2]{\sqbk{#1,\,#2}} % 交換子
\newcommand{\faweakcomm}[2]{\sqbk{#1,\,#2}_{\txtweak}} % 弱交換子
\newcommand{\supercommutator}[2]{\sqbk{#1,\,#2}_{\txtsuper}} % 超交換子
%%%%%%%%%%%%%%%%%%%%%%%%%%
% ディラックのブラケット %
%%%%%%%%%%%%%%%%%%%%%%%%%%
\newcommand{\braEPR}{\bra{\mathrm{EPR}}} % EPRブラベクトル, 量子情報
\newcommand{\braplus}{\bra{+}} % スピンなどで基底を表すブラ
\newcommand{\bra}[1]{\left\langle #1 \right|} % ブラベクトル
\newcommand{\ketEPR}{\ket{\mathrm{EPR}}} % EPR状態
\newcommand{\ketGHZ}{\ket{\mathrm{GHZ}}} % GHZ状態, expedition0009308
\newcommand{\ketW}{\ket{\mathrm{W}}} % W状態, expedition0009307
\newcommand{\ketbraproj}[2]{P_{#2,#1}} % \ketbraが見づらい場合の記法
\newcommand{\ketbra}[2]{\ket{#1} \!\! \bra{#2}} % \ketと\braをくっつけて一体の作用素のように見せたいとき
\newcommand{\ketminus}{\ket{-}} % スピンなどで基底を表すケット
\newcommand{\ketone}{\ket{1}} % 二準位系での基底
\newcommand{\ketplus}{\ket{+}} % スピンなどで基底を表すブラ
\newcommand{\ketzero}{\ket{0}} % 二準位系での基底
\newcommand{\ket}[1]{\left| #1 \right\rangle} % ケットベクトル
%%%%%%%%%%%%%%%%%%
% 点列・族の添字 %
%%%%%%%%%%%%%%%%%%
\newcommand{\kinone}{k \in \semigrposint} % k in N_1
\newcommand{\kin}{k \in \monnat} % k in N
\newcommand{\lil}{\lambda \in \Lambda} % lambda in Lambda
\newcommand{\mim}{\mu \in M} % mu in M
\newcommand{\ninone}{n \in \semigrposint} % n in N_1
\newcommand{\nin}{n \in \monnat} % n in N
\newcommand{\niz}{n \in \ringratint} % n in Z
%%%%%%%%%%
% Wick積 %
%%%%%%%%%%
\newcommand{\nna}{\stackrel{\scriptscriptstyle \circ}{\scriptscriptstyle \circ}\!}
\newcommand{\nne}{\!\stackrel{\scriptscriptstyle \circ}{\scriptscriptstyle \circ}}
\newcommand{\xxa}{\stackrel{\scriptscriptstyle \times}{\scriptscriptstyle \times}\!}
\newcommand{\xxe}{\!\stackrel{\scriptscriptstyle \times}{\scriptscriptstyle \times}}
\newcommand{\bosonwick}[1]{\left. \xxa\! \hspace{0.3pt} #1 \hspace{0.3pt} \!\xxe \right.}
\newcommand{\gaugedwick}[1]{\left. \nna\! \hspace{0.3pt} #1 \hspace{0.3pt} \!\nne \right.}
\newcommand{\wick}[1]{\left. :\! \hspace{-0.5pt} #1 \hspace{-0.5pt} \!: \right.} % ウィック積
%%%%%%%%%%%%%%%%%%
% 一般・全体利用 %
%%%%%%%%%%%%%%%%%%
\DeclareMathOperator*{\argmax}{argmax} % 最大値を取る引数を与える写像
\DeclareMathOperator*{\essinf}{\mathrm{ess \, inf}} % 本質的下限
\DeclareMathOperator*{\essran}{\mathrm{ess \, ran}} % 本質的値域
\DeclareMathOperator*{\esssup}{\mathrm{ess \, sup}} % 本質的上限
\DeclareMathOperator*{\opexistsone}{\exists!}
\DeclareMathOperator*{\opexists}{\exists}
\DeclareMathOperator*{\opforall}{\forall}
\DeclareMathOperator{\arccosh}{arccosh} % 逆双曲線関数, 共通
\DeclareMathOperator{\arcsinh}{arcsinh} % 逆双曲線関数, 共通
\DeclareMathOperator{\arctanh}{arctanh} % 逆双曲線関数, 共通
\DeclareMathOperator{\bigextprod}{\bigwedge\nolimits} % 外積用の記号
\DeclareMathOperator{\cn}{cn} % 楕円関数, 共通
\DeclareMathOperator{\compl}{«} % 関数合成
\DeclareMathOperator{\compr}{»} % 関数合成
\DeclareMathOperator{\csch}{csch} % 双曲線関数, 双曲線余割, 共通
\DeclareMathOperator{\dn}{dn} % 楕円関数
\DeclareMathOperator{\len}{len} % 長さ, 共通
\DeclareMathOperator{\logpv}{Log} % 対数関数の主値
\DeclareMathOperator{\sech}{sech} % 双曲線余割, ハイパボリックセカント
\DeclareMathOperator{\sgn}{sgn} % 符号, 共通
\DeclareMathOperator{\sn}{sn} % 楕円積分
\DeclareMathOperator{\symdiff}{\raisebox{.2ex}{\scalebox{0.8}{$\triangle$}}} % 対称差
\DeclareMathOperator{\tn}{tn} % 楕円積分
\NewDocumentCommand{\imunit}{O{\mathsf{i}}}{#1} % 虚数単位, 共通
\NewDocumentCommand{\placeholder}{O{\bullet}}{#1} % 関手やホモロジーなどでハイフンや点で表す要素
\NewDocumentCommand{\trace}{O{\operatorname{Tr}}}{#1} % トレース
\newcommand{\Ad}{\operatorname{Ad}} % 群などの随伴, Adjoint, Ad_U A = U^{\ast} A UまたはAd_g A = gA\inv{g}
\newcommand{\Coim}{\operatorname{Coim}} % 余像, expedition0009942
\newcommand{\Coker}{\operatorname{Coker}} % 余核, expedition0009942
\newcommand{\Cor}{\operatorname{Cor}} % 余制限写像 expedition0008776
\newcommand{\End}{\operatorname{End}} % 自己準同型写像がなす集合
\newcommand{\Ext}{\operatorname{Ext}} % Ext群・関手
\newcommand{\Fix}{\operatorname{Fix}} % 写像の不動点
\newcommand{\Hom}{\operatorname{Hom}} % 準同型
\newcommand{\Image}{\operatorname{Im}} % 像の集合
\newcommand{\Index}{\operatorname{Index}} % 指数定理などの指数を想定
\newcommand{\Inn}{\operatorname{Inn}} % 内部自己同型群, expedition0009281
\newcommand{\Inv}{\operatorname{Inv}} % 可逆な写像
\newcommand{\Ker}{\operatorname{Ker}} % 核の集合
\newcommand{\Map}{\operatorname{Map}} % 写像
\newcommand{\Out}{\operatorname{Out}} % 外部自己同型群
\newcommand{\Ran}{\operatorname{Ran}} % 作用素論での値域を表す記号, 複素数の虚部と区別するため
\newcommand{\Supp}{\operatorname{Supp}} % 台, 加群の台
\newcommand{\Tor}{\operatorname{Tor}} % Tor群・関手
\newcommand{\Tot}{\operatorname{Tot}} % アーベル圏で全複体を与える関手, expedition0008697
\newcommand{\aaep}{\epsilon} % 代数解析で単位を分解する関数, expedition10579
\newcommand{\aaepconj}{\cmpconj{\epsilon}} % 代数解析で単位を分解する関数, expedition10579
\newcommand{\bigmiddleslash}[2]{\left. #1 \middle/ #2 \right.} % 主に商集合\setquotに利用
\newcommand{\celcius}{\operatorname{\setSymbolUpLeft{\circ}{C}}} % 摂氏温度の記号
\newcommand{\centimeter}{\operatorname{cm}} % 長さの単位
\newcommand{\closedran}{\mathop{\overline{\Ran}}} % 閉線型包
\newcommand{\clossqbk}[1]{\sqbk{#1}} % よく作用素環で閉包を表すために利用する
\newcommand{\clos}[1]{\overline{#1}} % 外測度など何らかの意味で拡張・閉包のような趣があるものの, 具体的に「---閉包」のような名前がない対象に利用
\newcommand{\cmpconjrbk}[1]{\overline{\rbk{#1}}} % 複素共役で内部に括弧をつける
\newcommand{\cmpconj}[1]{\overline{#1}} % 複素共役
\newcommand{\codim}{\opexists{codim}} % 余次元
\newcommand{\cod}{\operatorname{cod}} % 終域, codomain
\newcommand{\cof}{\operatorname{cof}} % 余因子
\newcommand{\coim}{\operatorname{coim}} % 余像, 小文字は射に対して適用する, expedition0007898
\newcommand{\coker}{\operatorname{coker}} % 余核, 小文字は射に対して適用する
\newcommand{\enumcomb}[2]{\setSymbolDownLeft{#2}{\mathrm{C}}_{#1}} % 少なくとも日本ではよく使われる組み合わせの記号
\newcommand{\enumperm}[2]{\setSymbolDownLeft{#2}{\mathrm{P}}_{#1}} % 順列
\newcommand{\const}{\operatorname{const}} % 定数
\newcommand{\dalembertian}{\Box} % ダランベール作用素
\newcommand{\dddual}[1]{#1^{\ast \ast \ast}} % 三重双対
\newcommand{\ddual}[1]{#1^{\ast \ast}} % 二重双対
\newcommand{\diag}{\operatorname{diag}} % 対角
\newcommand{\diam}{\operatorname{diam}} % 直径, 距離空間・特異単体・単体複体・リーマン幾何など各所で利用
\newcommand{\diracdelta}{\delta} % ディラックのデルタ関数
\newcommand{\diracslash}{\not \!} % ディラック方程式などで出てくるガンマ行列との縮約を取るスラッシュ
\newcommand{\dom}{\operatorname{dom}} % 定義域
\newcommand{\dualrbk}[1]{\rbk{#1}^{\ast}} % 双対作用素
\newcommand{\dualsharprbk}[1]{\rbk{#1}^{\#}} % 双対, 特に空間の双対に利用
\newcommand{\dualsharp}[1]{#1^{\#}} % 双対, 特に空間の双対に利用
\newcommand{\dual}[1]{#1^{\ast}} % 双対作用素
\newcommand{\epN}{\ep \hyphen N} % ε-N論法
\newcommand{\epdelta}{\ep \hyphen \delta} % ε-δ論法
\newcommand{\ep}{\varepsilon} % エイリアス
\newcommand{\eqclr}[1]{\lrbk{#1}} % 同値類の記号
\newcommand{\eqcol}[1]{\overline{#1}} % 同値類の記号
\newcommand{\eqcsq}[1]{\sqbk{#1}} % 同値類の記号
\newcommand{\eqdef}{\overset{\mathrm{def}}{=}} % 定義に利用する統合
\newcommand{\eqisom}{\cong} % 一般的に同型を表す状況に利用
\newcommand{\eval}{\operatorname{ev}} % 評価写像
\newcommand{\fml}[2]{\cbk{#1}_{#2}} % 族の基本系
\newcommand{\highlightcap}[3][yellow]{\tikz[baseline=(x.base)]{\node[rectangle,rounded corners,fill=#1!10](x){#2} node[below of=x, color=#1]{#3};}} % 強調のtikzサンプルで利用
\newcommand{\highlight}[2][yellow]{\tikz[baseline=(x.base)]{\node[rectangle,rounded corners,fill=#1!10](x){#2};}} % 強調のtikzサンプルで利用
\newcommand{\hypergeometricfunction}[2]{{\vphantom{F}}_{#1} F_{#2}} % 超幾何関数
\newcommand{\hyphen}{\hbox{-}} % 記号定義の補助
\newcommand{\idonebold}{\boldsymbol{1}} % 太字の1で表した恒等写像
\newcommand{\idone}{1} % 1で表した恒等写像
\newcommand{\id}{\mathrm{id}} % 恒等写像
\newcommand{\image}{\operatorname{im}} % 像に関する射
\newcommand{\inftwo}[2]{\mathop{\inf_{#1}}_{#2}} % 下限
\newcommand{\inftysmall}[1]{\mathop{d #1}} % 無限小, 超準解析用の無限小は \nainftysmall で分ける
\newcommand{\invrbk}[1]{\rbk{#1}^{-1}} % 逆写像
\newcommand{\inv}[1]{#1^{-1}} % 逆像・逆写像, \invもよく使うため独立させる
\newcommand{\kroneckerdelta}{\delta} % クロネッカーのデルタ
\newcommand{\laplacian}{\Delta} % ラプラシアン
\newcommand{\lusharp}[1]{\setSymbolUpLeft{\sharp}{#1}} % TODO expedition0006072 定義箇所で記号の可能性だけ紹介してpostcompかprecompに置き換え
\newcommand{\maxtwo}[2]{\mathop{\max_{#1}}_{#2}} % 最大値
\newcommand{\barmean}[1]{\overline{#1}} % 上つきバーによる平均
\newcommand{\meter}{\operatorname{m}} % メートル
\newcommand{\napiernum}{\mathsf{e}} % ネイピア数, 自然対数の底, 素電荷など他にもeが出てきて紛らわしい場合に利用
\newcommand{\net}[2]{\rbk{#1}_{#2}} % ネット
\newcommand{\od}[2]{\frac{d #1}{d #2}} % 一階の常微分作用素を作用させた関数
\newcommand{\onehalf}{\frac{1}{2}} % 1/2
\newcommand{\oneoverfour}{\frac{1}{4}} % 1/4
\newcommand{\overarc}[1]{\stackrel{\Large\mbox{$\frown$}}{#1}} % 円弧
\newcommand{\opchern}[1]{\operatorname{ch}} % チャーン指標
\newcommand{\opexclude}[1]{\widehat{#1}} % その項を除外する作用素
\newcommand{\ophull}{\operatorname{hull}} % 包
\newcommand{\opideal}{\operatorname{Ideal}} % イデアル
\newcommand{\opimag}{\operatorname{Im}} % 複素数の虚部
\newcommand{\opod}[1]{\frac{d}{d #1}} % 一階の常微分作用素
\newcommand{\oppd}[1]{\frac{\partial}{\partial #1}} % 偏微分作用素
\newcommand{\oppv}{\operatorname{pv}} % 主値, principal value
\newcommand{\opreal}{\operatorname{Re}} % 複素数の実部
\newcommand{\otherwise}{\mathrm{otherwise}} % 場合分けで「その他」に使う文字列
\newcommand{\pd}[2]{\frac{\partial #1}{\partial #2}} % 偏導関数
\newcommand{\pinfty}{p^{\infty}} % p進系の記号
\newcommand{\predualsharp}[1]{#1_{\#}} % 前双対
\newcommand{\rel}{\mathbin{\mathrm{rel}}} % ホモトピーでの相対的な持ち上げなどに使うrel, expedition0008480
\newcommand{\rusharprbk}[1]{\rbk{#1}^{\sharp}} % p系の処理, 置換して削除
\newcommand{\rusharp}[1]{#1^{\sharp}} % TODO expedition0006072 定義箇所で記号の可能性だけ紹介してpostcompかprecompに置き換え
\newcommand{\setSymbolDownLeft}[2]{{\vphantom{#2}}_{#1}{#2}} % 左下に添字をつけるvphantomのエイリアス
\newcommand{\setSymbolUpLeft}[2]{{\vphantom{#2}}^{#1}{#2}} % 左上に添字をつけるvphantomのエイリアス
\newcommand{\fnsineintegral}[1]{\fun{\mathrm{Si}}{#1}} % 正弦積分の記号 Si(x) の部分
\newcommand{\intsineintegral}[1]{\int_{0}^{#1} \frac{\sin t}{t} \opdmsr{t}} % 正弦積分の定義の積分
\newcommand{\supp}{\operatorname{supp}} % 台, 共通
\newcommand{\vecbf}[1]{\mathbf{#1}} % ベクトル用のコマンド
\newcommand{\vecarrow}[1]{\overrightarrow{#1}} % 矢印のベクトル
%%%%%%%%%%%%%%%%%%
% 以下ソート対象 %
%%%%%%%%%%%%%%%%%%
\DeclareMathOperator*{\algbigoplus}{\hat{\bigoplus}} % 代数的直和：フォック空間など解析で閉包と区別するため
\DeclareMathOperator*{\algbigotimes}{\hat{\bigotimes}} % 代数的テンソル：フォック空間など解析で閉包と区別するため
\DeclareMathOperator*{\algoplus}{\hat{\oplus}} % 代数的直和：フォック空間など解析で閉包と区別するため
\DeclareMathOperator*{\algotimes}{\hat{\otimes}} % 代数的テンソル積：フォック空間など解析で閉包と区別するため
\DeclareMathOperator*{\closedotimes}{\gtclos{\otimes}} % 代数的テンソル積の閉包：作用素環でよく出てくる
\DeclareMathOperator*{\freeprod}{\ast} % 群の自由積, expedition0007871
\DeclareMathOperator{\eqalgisom}{\cong} % 代数的な同型
\DeclareMathOperator{\eqapprox}{\sim} % 近似的に等しい
\DeclareMathOperator{\eqcatequiv}{\simeq} % 圏同値
\DeclareMathOperator{\eqcatisom}{\cong} % 圏の同型, expedition0006298
\DeclareMathOperator{\eqgtisom}{\approx} % 同相
\DeclareMathOperator{\eqhilbunitaryisom}{\cong} % ヒルベルト空間のユニタリ同値
\DeclareMathOperator{\eqhomotopic}{\simeq} % 道ホモトピー同値
\DeclareMathOperator{\eqoaquasiequiv}{\simeq} % $\oacstar$-環の準同値
\DeclareMathOperator{\eqpathhomotopic}{\dot{\simeq}} % 道ホモトピー同値
\DeclareMathOperator{\eqprbiddistrib}{\overset{d}{=}} % 確率論での同分布
\DeclareMathOperator{\eqprbmsrequiv}{\sim} % 確率論での測度の同値：互いに絶対連続
\DeclareMathOperator{\eqsimelem}{\sim} % 同値な二要素を表す記号
\DeclareMathOperator{\eqstarisom}{\simeq} % スター同型($\ast$-同型)
\DeclareMathOperator{\equnitaryequiv}{\simeq} % ユニタリ同値
\DeclareMathOperator{\lahodgestar}{\star} % ホッジ作用素
\DeclareMathOperator{\noteqstarisom}{\not\simeq} % スター同型($\ast$-同型)
\DeclareMathOperator{\Rep}{Rep} % 表現の空間や圏
\NewDocumentCommand{\agvariety}{O{\mathcal}}{#1} % 代数多様体, 特にイデアルが生成する代数多様体
\NewDocumentCommand{\cmdrel}{O{\omega}}{#1} % 分散関係
\NewDocumentCommand{\dfsp}{O{A}}{#1} % 微分形式の空間
\NewDocumentCommand{\eqcpointed}{O{\eqcsq} m}{#1{#2}_{\ast}} % 基点つき同値類
\NewDocumentCommand{\fnheaviside}{O{H}}{#1} % ヘビサイド関数
\NewDocumentCommand{\fthol}{O{\mathcal{O}}}{#1} % 正則関数の空間用
\NewDocumentCommand{\ftmero}{O{\mathcal{M}}}{#1} % 有理型関数の空間用
\NewDocumentCommand{\grcentralizer}{O{Z}}{#1} % 中心化群, expedition0005395
\NewDocumentCommand{\grmetform}{O{2} m m}{\grmet[#1] \! \rbkt{#2}{#3}} % (一般)相対性理論, 計量の二次形式
\NewDocumentCommand{\grmet}{O{2}}{\setSymbolDownLeft{#1}{g}} % (一般)相対性理論, 計量の文字
\NewDocumentCommand{\grnormalizer}{O{N}}{#1} % 正規化群, expedition0005395
\NewDocumentCommand{\gropasym}{O{A}}{#1} % 反対称化作用素, expedition0001365
\NewDocumentCommand{\gropsym}{O{S}}{#1} % 対称化作用素, expedition0001365
\NewDocumentCommand{\grpermorderedpair}{O{\mathcal{P}}}{#1} % 偶数個の要素の異なる順序対への分解の全体, expedition0009453
\NewDocumentCommand{\grsym}{O{\mathfrak{S}} m}{#1_{#2}} % 対称群, expedition0001319
\NewDocumentCommand{\gtbase}{O{\mathcal}}{#1} % 位相の基底, 開集合の基底, expedition0000259
\NewDocumentCommand{\gtfilter}{O{\mathcal}}{#1} % フィルター, expedition0000349
\NewDocumentCommand{\gtfmlclosed}{O{\mathcal}}{#1} % 閉集合系だけではなく, ある集合上の一般の閉集合族を表す, expedition0000227
\NewDocumentCommand{\gtfmlopen}{O{\mathcal}}{#1} % 「位相」の意味での開集合系だけではなく, ある集合上の一般の開集合族を表す, expedition0000227
\NewDocumentCommand{\gtopenball}{O{U}}{#1} % 開球
\NewDocumentCommand{\gtopencover}{O{\mathcal}}{#1} % 開被覆
\NewDocumentCommand{\gtopennbh}{O{\mathcal}}{#1} % 開近傍系
\NewDocumentCommand{\gtpreopencover}{O{\mathcal}}{#1} % 内部を取ると開被覆になる集合の族, expedition0009107
\NewDocumentCommand{\gtsubbase}{O{\mathcal}}{#1} % 位相の準基, expedition0005835
\NewDocumentCommand{\gtvicinity}{O{\mathcal}}{#1} % 一様構造での近縁
\NewDocumentCommand{\lasp}{O{\mathcal}}{#1} % 線型空間
\NewDocumentCommand{\latprightrbk}{O{\top} m}{\rbk{#2}^{#1}} % 転置, expedition0009877
\NewDocumentCommand{\latpright}{O{\top} m}{#2^{#1}} % 転置, expedition0009877
\NewDocumentCommand{\latp}{O{t} m}{\setSymbolUpLeft{#1}{#2}} % 転置, expedition0009877
\NewDocumentCommand{\lpdistribution}{O{\mu} m}{#2_{\ast,#1}} % 分布関数, expedition0010310
\NewDocumentCommand{\lpmollifier}{O{\rho}}{#1} % フリードリクスの軟化子
\NewDocumentCommand{\lpofpositive}{O{\chi}}{#1} % 正型関数
\NewDocumentCommand{\manliederiv}{O{L}}{#1} % リー微分
\NewDocumentCommand{\mansmoothnbh}{O{\mathcal}}{#1} % 多様体上の滑らかな開近傍の族
\NewDocumentCommand{\mblfmldsysgenerated}{O{d} m}{\fun{#1}{#2}} % 生成された$d$-系
\NewDocumentCommand{\mblfmlgenerated}{O{\sigma} m}{\fun{#1}{#2}} % 生成された加法族
\NewDocumentCommand{\oacorrfn}{O{\Gamma}}{#1} % 汎関数積分表示とも関わる相関関数, expedition0011645
\NewDocumentCommand{\oagnsvector}{O{\Omega}}{#1} % GNS表現のベクトル状態
\NewDocumentCommand{\oaideal}{O{\mathcal}}{#1} % 作用素環のイデアル：環や代数のイデアルとは分けておく
\NewDocumentCommand{\oanumberoperator}{O{A}}{#1} % レゾルベント環での個数作用素
\NewDocumentCommand{\oaposcone}{O{\mathcal{P}}}{#1} % 冨田-竹崎理論での正錐
\NewDocumentCommand{\oapressure}{O{P}}{#1} % トレースによる通常の圧力は大文字でP: expedition0009842, トレース状態による修正圧力は小文字でp: expedition0009797
\NewDocumentCommand{\oarepn}{O{\pi}}{#1} % 作用素環の表現
\NewDocumentCommand{\oaspnormalstate}{O{N}}{#1} % 作用素環上の正規状態の空間
\NewDocumentCommand{\oasppurestate}{O{P}}{#1} % 作用素環上の純粋状態の空間
\NewDocumentCommand{\oaspstate}{O{E}}{#1} % 作用素環の状態がなす空間
\NewDocumentCommand{\oastatevector}{O{\Omega}}{#1} % GNS表現に付随する状態ベクトル
\NewDocumentCommand{\oastate}{O{\omega}}{#1} % 作用素環の状態
\NewDocumentCommand{\opdilation}{O{\delta}}{#1} % スケーリング(伸張)の作用素, expedition0010830
\NewDocumentCommand{\opdmat}{O{\rho}}{#1} % 密度行列(密度作用素)
\NewDocumentCommand{\opfockan}{O{a}}{#1} % 場の量子論での消滅作用素
\NewDocumentCommand{\opfockcran}{O{a}}{#1^{\#}} % 場の量子論での生成・消滅作用素をまとめて表す記法
\NewDocumentCommand{\opfockcrdagger}{O{a}}{#1^{\dagger}} % 場の量子論での生成作用素のバージョン
\NewDocumentCommand{\opfockcr}{O{a}}{#1^{\ast}} % 場の量子論での生成作用素
\NewDocumentCommand{\opfocknumber}{O{N}}{#1} % 個数作用素
\NewDocumentCommand{\opfocksegalconj}{O{\pi}}{#1} % 場の量子論でのシーガルの場の作用素, expedition0009962
\NewDocumentCommand{\opfocksegal}{O{\phi}}{#1} % 場の量子論でのシーガルの場の作用素
\NewDocumentCommand{\opspecmeas}{O{E}}{#1} % スペクトル測度の標準的な記号
\NewDocumentCommand{\opspec}{O{} m}{\fun{\sigma_{#1}}{#2}} % 作用素のスペクトル, expedition0001551
\NewDocumentCommand{\optransl}{O{\tau}}{#1} % 平行移動の作用素, expedition0010829
\NewDocumentCommand{\opvarspec}{O{} m}{\operatorname{spec}_{#1} #2} % 作用素のスペクトル, スピン系の議論など \sigma ではまぎらわしい場合に利用, expedition0001551
\NewDocumentCommand{\physaction}{O{\mathcal{A}}}{#1} % 物理の作用
\NewDocumentCommand{\physcharge}{O{e}}{\mathrm{#1}} % 電荷, 素電荷などの物理量を表す. e以外にもqに利用する
\NewDocumentCommand{\physcplconst}{O{\mathsf{g}}}{#1} % 結合定数, coupling constant
\NewDocumentCommand{\physelectrostaticcapasity}{O{\mathrm{Cap}}}{#1} % 静電容量, expedition0010917
\NewDocumentCommand{\physenergy}{O{E}}{#1} % エネルギー
\NewDocumentCommand{\physgse}{O{E}}{#1_{0}} % 基底エネルギー, ground state energy
\NewDocumentCommand{\physham}{O{H}}{#1} % ハミルトニアンの雛形, たいていは下つき添字で水素原子, 自由電磁場などを表す
\NewDocumentCommand{\physlagdensity}{O{\mathcal{L}}}{#1} % ラグランジアン密度の雛形
\NewDocumentCommand{\physlag}{O{L}}{#1} % ラグランジアンの雛形
\NewDocumentCommand{\physliouvilean}{O{L}}{#1} % リウビリアンの雛形
\NewDocumentCommand{\physmass}{O{m}}{#1} % 物質の質量
\NewDocumentCommand{\prbbrownmv}{O{B}}{#1} % ブラウン運動
\NewDocumentCommand{\prbcharfun}{O{\chi}}{#1} % 特性関数
\NewDocumentCommand{\prbdist}{O{\mathcal{P}}}{#1} % 確率分布
\NewDocumentCommand{\prbgaussianmeasure}{O{\msrcal{N}}}{#1} % ガウス測度
\NewDocumentCommand{\prbnormaldist}{O{N}}{#1} % 正規分布
\NewDocumentCommand{\prbpoissonprocess}{O{N}}{#1} % ポアソン過程
\NewDocumentCommand{\prbprocess}{O{X}}{#1} % 確率過程
\NewDocumentCommand{\prbqspace}{O{\mathcal{Q}}}{#1} % $Q$-表示の土台の集合を表す
\NewDocumentCommand{\prbspsample}{O{\Omega}}{#1} % 確率空間に対する標本空間, expedition0010114
\NewDocumentCommand{\psh}{O{\mathfrak}}{#1} % 前層
\NewDocumentCommand{\qtquantumchannel}{O{\mathcal{L}}}{#1} % 量子通信路, expedition0005561, expedition0006844
\NewDocumentCommand{\repn}{O{\pi}}{#1} % 表現
\NewDocumentCommand{\schattencls}{O{\mathbb{K}}}{#1} % コンパクト作用素またはシャッテンクラス
\NewDocumentCommand{\setfmlcylinder}{O{\mathcal{C}}}{#1} % 筒集合族, expedition0000313, expedition0004973
\NewDocumentCommand{\setfml}{O{\mathcal}}{#1} % 集合族
\NewDocumentCommand{\setindex}{O{\mathcal} m}{#1{#2}} % 添字集合
\NewDocumentCommand{\setlattice}{O{\Gamma}}{#1} % 格子, expedition0010278
\NewDocumentCommand{\setspecial}{O{\mathcal} m}{#1{#2}} % 特別な集合を表すための記号, 一点集合・二点集合などを太字で表すといった用途
\NewDocumentCommand{\shdiffform}{O{\sheaf{A}}}{#1} % 微分形式がなす層
\NewDocumentCommand{\sheaf}{O{\mathfrak}}{#1} % 層の記号
\NewDocumentCommand{\smchemicalpotential}{O{\mu}}{#1} % 化学ポテンシャル
\NewDocumentCommand{\smenergydensity}{O{\varrho}}{#1} % エネルギー密度, expedition0011356
\NewDocumentCommand{\smfluctuationwithdmat}{O{\beta} m}{\smuncertaintywithdmat[#1]{#2}^2} % 密度行列に付随する正規状態でのゆらぎ, expedition0011059
\NewDocumentCommand{\sminvtemperature}{O{\beta}}{#1} % 逆温度
\NewDocumentCommand{\smlocaldensityoperator}{O{\rho}}{#1} % 大正準状態に対する局所密度作用素, expedition0011294
\NewDocumentCommand{\smmicrocanonicalstate}{O{\beta} m}{\physmean{#2}_{#1}} % 密度行列に付随する正規状態, expedition0011059
\NewDocumentCommand{\smnumberdensity}{O{\rho}}{#1} % 粒子数密度, expedition0011351
\NewDocumentCommand{\smooth}{O{\mathcal{E}}}{#1} % 滑らかな関数の空間, 特に微分形式関連でも利用, expedition0010649
\NewDocumentCommand{\smparticlenumber}{O{N}}{#1} % 平均粒子数, expedition0011351
\NewDocumentCommand{\smpressure}{O{p}}{#1} % 統計力学での圧力, expedition0011341
\NewDocumentCommand{\smspecificfreeenergy}{O{\bar{f}}}{#1} % (ヘルムホルツの)比自由エネルギー
\NewDocumentCommand{\smthermalvac}{O{\beta}}{\Omega_{#1}} % 熱的真空, expedition0011060
\NewDocumentCommand{\smuncertaintywithdmat}{O{\beta} m}{\rbk{\triangle #2}_{#1}} % 密度行列に付随する正規状態での不確かさ, expedition0011059
\NewDocumentCommand{\sphilbfrak}{O{\mathfrak}}{#1} % ヒルベルト空間
\NewDocumentCommand{\sphilb}{O{\mathcal}}{#1} % ヒルベルト空間
\NewDocumentCommand{\splowerhalf}{O{\mathbb{H}}}{#1_{\txtneg}} % 開下半空間
\NewDocumentCommand{\spupperhalf}{O{\mathbb{H}}}{#1_{\txtnonneg}} % 開上半空間
\NewDocumentCommand{\topmetric}{O{d}}{#1} % 位相空間上の距離関数
\NewDocumentCommand{\vaoutnormal}{O{\widehat}}{#1} % 外向き外法線
\newcommand{\aaantipodalmap}[1]{#1^{a}} % 対蹠写像, expedition0010618
\newcommand{\aadmodule}{\mathcal{D}} % D-加群
\newcommand{\aasingsupp}{\operatorname{sing \, supp}} % 特異台, expedition10583
\newcommand{\aasingspec}{\operatorname{SS}} % 超局所解析的な特異スペクトル, expedition0010613
\newcommand{\absconti}{\mathrm{AC}} % 絶対連続な関数がなす空間
\newcommand{\affext}[1]{\overline{#1}} % アフィン拡大実数・拡大複素数(リーマン球面)に利用
\newcommand{\aghomcoord}[1]{\sqbk{#1}} % 射影空間の斉次座標
\newcommand{\agproj}[2]{P_{#1}^{#2}} % 代数幾何での射影空間, P_{#1}^{#2}
\newcommand{\agtoideal}{\mathcal{I}} % (代数多様体の)部分集合からイデアルへの写像
\newcommand{\agtovariety}{\mathcal{V}} % イデアルから代数多様体(の部分集合)への写像
\newcommand{\algad}{\operatorname{ad}} % 環・代数の交換子を生成するad, 群のAdの微分, expedition0004725
\newcommand{\algdiffop}{\mathfrak{D}} % 微分作用素がなす環, 例 expedition0009836, expedition0007330, expedition0007329
\newcommand{\algendskew}{\operatorname{Endskew}} % expedition0004007
\newcommand{\algevenodd}[1]{#1_{\pm}} % (超代数の)偶項・奇項
\newcommand{\algeven}[1]{#1_{\txteven}} % (超代数の)偶項
\newcommand{\algideal}[1]{\mathfrak{#1}} % 環または代数のイデアル
\newcommand{\algodd}[1]{#1_{\txtodd}} % (超代数の)奇項
\newcommand{\algsetderivative}{\operatorname{Der}} % 代数上の微分作用素がなす集合, expedition0002660
\newcommand{\catabfingen}{\category{FinGenAb}} % 有限生成な可換群がなす圏
\newcommand{\catab}{\category{Ab}} % 可換群がなす圏, expedition0006433
\newcommand{\cataddfuncatoab}{\mathbf{Hom}^{\mathbf{a}}(\category{A}, \category{Ab})} % 小さなアーベル圏$\category{A}$から小さなアーベル群がなす圏$\category{Ab}$への加法関手がなす圏, expedition0007964
\newcommand{\catban}[1]{\category{Ban}_{#1}} % バナッハ空間がなす圏, expedition0009393
\newcommand{\catcatlocsmall}{\category{CAT}} % 局所小圏を対象とする圏, expedition0006288, expedition0009286
\newcommand{\catcatsmall}{\category{Cat}} % 小さな圏を対象とする圏
\newcommand{\catcluster}{\category{Cluster}} % クラスターとその細分化がなす圏
\newcommand{\catcomma}[2]{#1 \downarrow #2} % コンマ圏
\newcommand{\catcpthaus}{\category{CptHaus}} % コンパクトハウスドルフ空間がなす圏
\newcommand{\catdblcpx}[1]{\fun{\category{C}_2}{#1}} % 二重複体がなす圏
\newcommand{\catdirgraph}{\category{DirGraph}} % 有向グラフがなす圏
\newcommand{\category}[1]{\mathop{\mathsf{#1}}} % 圏を生成する
\newcommand{\cateuclidfinbased}{\cateuclidfin_{\ast}} % 基点つき有限次元ユークリッド空間がなす圏, expedition0009222
\newcommand{\cateuclidfin}{\category{FinEuclid}} % ユークリッド空間がなす圏
\newcommand{\catexactseq}[1]{\fun{\mathrm{ES}}{{#1}}} % 完全系列がなす圏
\newcommand{\catfield}{\category{Field}} % 体がなす圏
\newcommand{\catfock}{\category{Fock}} % フォック空間がなす圏
\newcommand{\catfunctor}[2]{#1^{#2}} % 関手圏
\newcommand{\catgraph}{\category{Graph}} % グラフがなす圏 expedition0009126
\newcommand{\catgroupoid}{\category{Groupoid}} % 亜群がなす圏
\newcommand{\catgrpfingen}{\category{FinGenGrp}} % 有限生成な群がなす圏
\newcommand{\catgrp}{\category{Grp}} % 群がなす圏
\newcommand{\cathaus}{\category{Haus}} % ハウスドルフ空間がなす圏
\newcommand{\cathilb}{\category{Hilb}} % ヒルベルト空間がなす圏
\newcommand{\cathom}[1]{\fun{\category{Hom}}{#1}} % 関手の圏, expedition0009417
\newcommand{\cathorcomp}{\mathop{\ast}} % 自然変換の水平合成, expedition0009442
\newcommand{\cathtpybased}{\cathtpy_{\ast}} % 基点を保つホモトピー類がなす圏, expedition0009144
\newcommand{\cathtpy}{\category{Htpy}} % ホモトピー類がなす圏
\newcommand{\catintermediatefield}[2]{\catfield_{#1}^{#2}} % 中間体がなす圏, expedition0009367
\newcommand{\catmanbased}{\catman_{\ast}} % 基点つき多様体の圏
\newcommand{\catman}{\category{Man}}
\newcommand{\catmat}[1]{\category{Mat}_{#1}} % 環係数行列がなす圏, expedition0009135
\newcommand{\catmeasure}{\category{Measure}} % 測度空間がなす圏, expedition0009145
\newcommand{\catmeas}{\category{Meas}} % 可測空間がなす圏, expedition0009131
\newcommand{\catmetricfin}{\category{FinMetric}} % 有限な距離空間がなす圏
\newcommand{\catmodcpx}[1]{#1 \mathchar`- \category{ModChCpx}}
\newcommand{\catmodel}[1]{\category{Model}_{#1}} % 数理論理学でのモデルがなす圏, expedition0009133
\newcommand{\catmodfingenleft}[1]{#1 \mathchar`- \catmodfingen} % 有限生成左加群がなす圏
\newcommand{\catmodfingenright}[1]{\catmodfingen \mathchar`- #1} % 有限生成右加群がなす圏
\newcommand{\catmodfingen}{\category{FinGenMod}} % 有限生成な加群がなす圏
\newcommand{\catmodgrleft}[1]{#1 \mathchar`- \catmodgr} % 次数つき加群がなす圏
\newcommand{\catmodgr}{\category{GrMod}} % 次数つき加群がなす圏
\newcommand{\catmodlefts}{\catmodleft{S}} % 左$S$-加群がなす圏
\newcommand{\catmodleft}[1]{#1 \mathchar`- \catmod} % 左加群がなす圏：係数環は別途指定する
\newcommand{\catmodright}[1]{\catmod \mathchar`- #1} % 右加群がなす圏：係数環は別途指定する
\newcommand{\catmodtfleft}[1]{#1 \mathchar`- \catmodtf} % 左無捻加群がなす圏
\newcommand{\catmodtf}{\category{TFMod}} % 無捻加群がなす圏, torsion-freeからTF
\newcommand{\catmod}{\category{Mod}} % 加群の圏
\newcommand{\catmonoidcommut}{\category{CMon}} % 可換なモノイドがなす圏
\newcommand{\catmon}{\category{Mon}} % モノイドの圏
\newcommand{\catmor}{\category{Mor}} % 射の集まり
\newcommand{\catobj}[1]{\operatorname{Ob} #1}
\newcommand{\catopposite}[1]{{#1}^{\txtopposite}} % 反対圏に利用.
\newcommand{\catopsets}{\category{O}} % 位相空間の開集合全体がなす圏, expedition0008797
\newcommand{\catorbitforgrp}[1]{\category{O}_{#1}} % 軌道圏, expedition0009364
\newcommand{\catordfin}{\category{FinOrd}} % 有限順序数がなす圏
\newcommand{\catpcluster}{\category{PCluster}} % パーシステントクラスター
\newcommand{\catposet}{\category{Poset}} % 半順序集合がなす圏, expedition0009132
\newcommand{\catpresheaf}[1]{\category{PSh}} % 前層がなす圏
\newcommand{\catringcommut}{\category{CRing}} % 可換環の圏
\newcommand{\catring}{\category{Ring}} % 単位的環がなす圏
\newcommand{\catrng}{\category{Rng}} % 単位的とは限らない環がなす圏
\newcommand{\catsetbased}{\catset_{\ast}} % 基点つき集合がなす圏, expedition0009121
\newcommand{\catsetfinbased}{\catsetfin_{\ast}} % 基点つき有限集合と基点を保つ写像がなす圏
\newcommand{\catsetfinisom}{\catsetfin_{iso}} % 有限集合と全単射がなす圏
\newcommand{\catsetfin}{\category{Fin}} % 有限集合と写像がなす圏
\newcommand{\catsetop}{\catopposite{\category{Set}}} % 集合がなす圏の反対圏
\newcommand{\catsetpdeffn}{\category{Set}^{\partial}} % 集合と部分関数がなす圏, expedition0009371
\newcommand{\catset}{\category{Set}} % 集合がなす圏
\newcommand{\catsheaf}{\category{Sh}} % 層がなす圏
\newcommand{\catshortexactseq}{\category{SES}} % 短完全系列がなす圏
\newcommand{\catsimplicial}{\Delta} % 単体圏, 非空の有限順序数と順序写像がなす圏
\newcommand{\catsingleobj}{\category{B} } % 一対象圏, BGのような形で使う, expedition0009146
\newcommand{\catskeleton}{\category{sk}} % 圏の骨格, expedition0009428
\newcommand{\catsliceover}[2]{\bigmiddleslash{#1}{#2}} % $c$上のスライス圏, expedition0009366
\newcommand{\catsliceunder}[2]{\bigmiddleslash{#1}{#2}} % $c$のもとでのスライス圏, expedition0009365
\newcommand{\cattopawconnbased}{\cattop_{\ast}^{\txtpathwiseconnected}} % 弧状連結な基点つき位相空間がなす圏, expedition0009427
\newcommand{\cattopbased}{\cattop_{\ast}} % 基点つき位相空間がなす圏
\newcommand{\cattopop}{\catopposite{\category{Top}}} % 位相空間がなす圏の反対圏
\newcommand{\cattop}{\category{Top}} % 位相空間がなす圏
\newcommand{\cattrgroupoid}[2]{\category{T}_{#1} #2} % 並進亜群, expedition0009430
\newcommand{\catuniverse}{\mathfrak{U}} % 圏論での宇宙
\newcommand{\catvect}{\category{Vect}} % 線型空間の圏, expedition0009125, expedition0009418
\newcommand{\clsordinals}{\operatorname{ON}} % 順序数の全体(クラス)
\newcommand{\contiseq}{c} % 数列空間用の文字, expedition0010228, expedition0010278
\newcommand{\conti}{C} % 連続または滑らかな関数の空間
\newcommand{\dgconnection}{\nabla} % 接続
\newcommand{\dgcurvatureform}{\Omega} % 接続形式の曲率2-形式, expedition0009944
\newcommand{\dggrgauge}{\mathcal{G}} % ゲージ群, ゲージ変換群, expedition0003652
\newcommand{\dgharmform}{\mathbb{H}} % 調和関数または調和形式の空間
\newcommand{\dgricci}{\operatorname{Ric}} % リッチ曲率, expedition0003770
\newcommand{\dgriemmetrbk}[3]{\dgriemmet{\rbk{#1}}{#2}{#3}} % リーマン計量
\newcommand{\dgriemmet}[3]{#1 \! \rbkt{#2}{#3}} % リーマン計量
\newcommand{\dgspconnection}{\mathcal{A}} % 接続がなす空間, expedition0003652
\newcommand{\dgvolform}{\mathrm{vol}} % 体積形式. 体積関数\fnvolと区別する
\newcommand{\distalgsed}{\mathbb{S}} % 十六元数がなす非結合的分配多元体(distributive algebra)
\newcommand{\divalgoct}{\mathbb{O}} % 八元数全体がなす可除代数(division algebra)または多元体
\newcommand{\dpopgeneralizedricci}{\mathcal} % 一般化されたワイツェンベックの公式に対するリッチ曲率, expedition0003662
\newcommand{\dstrapiddec}{\mathcal{S}} % 急減少関数の空間
\newcommand{\dstschwartz}{\dsttest^{\ast}} % シュワルツ超関数の空間
\newcommand{\dsttempered}{\dstrapiddec^{\ast}} % 緩増加超関数の空間
\newcommand{\dsttest}{\mathcal{D}} % シュワルツ超関数に対する試験関数の空間
\newcommand{\faaadj}[1]{#1^{\ast \ast}} % 共役を二回作用させる
\newcommand{\faadjrbk}[1]{\rbk{#1}^{\ast}} % 行列・作用素に対する共役・随伴. 線型代数でも内積の必要性からこれに統合. 双対空間にはdualをあてる
\newcommand{\faadjsharprbk}[1]{\rbk{#1}^{\#}} % 行列・作用素に対する共役・随伴. H^{-s}のような双対とずれるソボレフ空間などに利用する
\newcommand{\faadjsharp}[1]{#1^{\#}} % 行列・作用素に対する共役・随伴. H^{-s}のような双対とずれるソボレフ空間などに利用する
\newcommand{\faadjpresharprbk}[1]{\rbk{#1}_{\#}} % 行列・作用素に対する前共役(前双対)
\newcommand{\faadjpresharp}[1]{#1_{\#}} % 行列・作用素に対する前共役(前双対)
\newcommand{\faadj}[1]{#1^{\ast}} % 行列・作用素に対する共役・随伴. 線型代数でも内積の必要性からこれに統合. 双対空間にはdualをあてる
\newcommand{\fabanachlimit}{\mathrm{BL}} % バナッハ極限, expedition0011076
\newcommand{\faclosedconvexcone}{\mathcal{C}} % 閉凸錐, expedition0010219
\newcommand{\facomplementarysum}{\operatorname{\dot{+}}} % 相補的な閉部分空間の和, expedition0010174
\newcommand{\fafitr}[1]{\check{#1}} % フーリエ逆変換
\newcommand{\fafouriertr}{\mathcal{F}} % フーリエ変換, 逆変換は\invか\faadjで生成
\newcommand{\faftrrbk}[1]{\widehat{\rbk{#1}}} % フーリエ変換
\newcommand{\faftr}[1]{\widehat{#1}} % フーリエ変換
\newcommand{\fafunctional}[1]{\mathcal{#1}} % 汎関数
\newcommand{\falaplacetr}{\mathcal{L}} % ラプラス変換
\newcommand{\faperprbk}[1]{\rbk{#1}^{\perp}} % 直交補空間
\newcommand{\faperp}[1]{#1^{\perp}} % 直交補空間
\newcommand{\fapolarset}[1]{#1^{o}} % 極集合, expedition0010442
\newcommand{\fasetantisym}[1]{\mathcal{#1}} % 反対称集合, expedition0001331
\newcommand{\fastrongod}[2]{\txtstrong \hyphen \od{#1}{#2}} % 強微分, 強常微分
\newcommand{\fastrongopod}[1]{\txtstrong \hyphen \opod{#1}} % 強微分, 強常微分
\newcommand{\fastop}[1]{\fun{\beta}{#1}} % 関数解析での強位相の指定, expedition0010403
\newcommand{\fatop}[1]{\fun{\sigma}{#1}} % 関数解析での弱位相の指定, expedition0010402
\newcommand{\faunitclosedball}[1]{\rbk{#1}_1} % 閉単位球
\newcommand{\faweakopod}[1]{\txtweak \hyphen \opod{#1}} % 弱微分, 弱常微分
\newcommand{\fawstar}[1]{#1_{\mathrm{w}^{\ast}}} % 双対空間に対する汎弱位相
\newcommand{\fawtop}[1]{#1_{\mathrm{w}}} % 双対空間に対する汎弱位相
\newcommand{\fldalgclos}[1]{\overline{#1}} % 体の代数閉包
\newcommand{\fldchar}{\operatorname{char}} % 体の標数
\newcommand{\fldcmp}{\fld{C}} % 複素数体
\newcommand{\fldextdeg}[2]{\sqbk{#1 : #2}} % 体の拡大次数, expedition0009274
\newcommand{\fldfin}[1]{\fld{F}_{#1}} % 有限体
\newcommand{\fldfrobmap}{\fld{F}} % フロベニウス写像
\newcommand{\fldgal}[1]{\fun{\operatorname{Gal}}{#1}} % ガロア, Galを冠する関数, expedition0005798
\newcommand{\fldlaurantseries}[2]{#1 \! \cbk{\cbk{#2}}} % ローラン級数体
\newcommand{\fldmultiplicativegroup}[1]{#1^{\times}} % 体の乗法群
\newcommand{\fldratdyadic}{D_{\fldrat}} % 二進有理数, expedition0000391
\newcommand{\fldratfun}[2]{\fun{#1}{#2}} % 有理関数体, 有理式体
\newcommand{\fldrat}{\fld{Q}} % 有理数体
\newcommand{\fldrealneg}{\fld{R}_{<0}} % 厳密に負の実数全体
\newcommand{\fldrealnn}{\fld{R}_{\geq 0}} % 非負の実数全体
\newcommand{\fldrealnp}{\fld{R}_{\leq 0}} % 非正の実数全体
\newcommand{\fldrealplus}{\fld{R}_{>0}} % 厳密に正の実数全体
\newcommand{\fldreal}{\fld{R}} % 実数体
\newcommand{\fldtofield}{\mathcal{F}} % ガロア部分群から中間体への写像, expedition0005813
\newcommand{\fldtogaloisgroup}{\mathcal{G}} % 中間体からガロア部分群への写像, expedition0005813
\newcommand{\fldtransdeg}{\operatorname{trans-deg}} % 超越次数
\newcommand{\fld}[1]{\mathbb{#1}} % 実数体・複素数体・四元数体, またはこれらをまとめて表す体用のコマンド
\newcommand{\fnappell}{\phi} % リーマンの$\zeta$の変形にあたるアッペル関数
\newcommand{\fnbeta}[1]{\fun{B}{#1}} % 特殊関数のベータ関数
\newcommand{\fnceil}[1]{\left\lceil #1 \right\rceil} % 天井関数
\newcommand{\fndef}[1]{\boldsymbol{1}_{#1}} % 定義関数
\newcommand{\fnerf}[1]{\fun{\operatorname{erf}}{#1}} % 誤差関数
\newcommand{\fnexp}[1]{\fun{\exp}{#1}} % 指数関数
\newcommand{\fnfloor}[1]{\lfloor #1 \rfloor} % 床関数
\newcommand{\fngamma}[1]{\fun{\Gamma}{#1}} % 特殊関数のガンマ関数, expedition0006311
\newcommand{\fngauss}[1]{\left\lfloor #1 \right\rfloor} % ガウス記号, $x$を超えない最大の整数
\newcommand{\fnmorphismofcomplex}[1]{\mathrm{#1}} % 複体の射
\newcommand{\fnrestr}[2]{\left. #1 \right|_{#2}} % 写像の制限
\newcommand{\fnspbdd}[1]{\fun{B}{#1}} % ある集合上で有界な関数がなす空間
\newcommand{\fnspconst}[1]{\fun{\operatorname{Const}}{#1}} % 定数関数の空間
\newcommand{\fntoclose}{i_{\txtclose}} % 閉集合への包含写像
\newcommand{\fntoopen}{i_{\txtopen}} % 開集合への包含写像
\newcommand{\fnvalued}[2]{#1_{#2}} % (df)_pのようにある点での値を指定する写像
\newcommand{\fnvol}[1]{\fun{\operatorname{Vol}}{#1}} %  体積を与える関数. 実際の体積を大文字のVolとし, 体積形式は小文字のvolにする
\newcommand{\ftcauchyhilberttransform}{\operatorname{\mathbf{CH}}} % コーシー-ヒルベルト変換, expedition0010491, expedition0010511
\newcommand{\ftcauchyweiltransform}{\operatorname{\mathbf{CW}}} % コーシー-ヴェイユ変換
\newcommand{\ftcht}[1]{\check{#1}} % 解析汎関数に対するコーシー-ヒルベルト変換の作用, expedition0010483
\newcommand{\ftdbar}{\overline{\partial}} % 関数論のディーバー作用素
\newcommand{\ftdirichletsolvableregion}{\operatorname{Reg}} % ディリクレ問題が解けるような領域の全体, expedition000928
\newcommand{\ftdivisor}{\operatorname{Div}} % 因子, expedition0009889
\newcommand{\ftdsharp}{\partial^{\#}} % 複素幾何での外微分作用素
\newcommand{\ftfbt}[1]{\widehat{#1}} % 解析汎関数に対するフーリエ-ボレル変換の作用, expedition0010483
\newcommand{\ftflt}[1]{\tilde{#1}} % 指数型整関数に対するフーリエ-ラプラス変換の作用, expedition0010545
\newcommand{\ftfourierboreltransform}{\mathop{\mathbf{FB}}} % フーリエ-ボレル変換, expedition0010483
\newcommand{\ftfourierlaplacetransform}{\mathop{\mathbf{FL}}} % フーリエ-ラプラス変換
\newcommand{\ftholconv}[1]{\widehat{#1}} % 正則凸包, expedition0010460
\newcommand{\ftleraycover}[2]{\mathfrak{#1}} % 関数論でのルレイ被覆, expedition0004155やそのあとの注意
\newcommand{\ftophull}{\mathop{\mathfrak{h}}} % 包作用素, expedition0004300
\newcommand{\ftord}{\operatorname{ord}} % 零点または極の位数, expedition0003167
\newcommand{\ftperronclass}[1]{\mathfrak{P}_{#1}} % ペロン類, expedition0009838
\newcommand{\ftplurisubharm}{\operatorname{PSH}} % 多重劣調和な関数全体, expedition0004554
\newcommand{\ftpolydisc}{\mathrm{P} \triangle} % 多重円板
\newcommand{\ftpolynomialconvex}[1]{\widehat{#1}^{\txtpolynomial}} % 多項式凸包, expedition0010463
\newcommand{\ftresidue}{\operatorname{Res}} % 関数論の留数
\newcommand{\ftspentirefunctionofexpfn}{\operatorname{Exp}} % 指数型整関数の空間, expedition0010485
\newcommand{\ftspecialbd}{\partial_0} % 特殊境界, expedition0005651
\newcommand{\ftstrongplurisubharm}{\mathrm{SPSH}} % 強多重劣調和な関数全体の集合, expedition0004554
\newcommand{\ftsubharm}{\mathrm{SH}} % 劣調和関数の空間, expedition0009448
\newcommand{\grabelize}[1]{#1^{\mathrm{ab}}} % 群のアーベル化, expedition0009330
\newcommand{\graffine}{\operatorname{Aff}} % アフィン変換, expedition0005496
\newcommand{\grauto}{\operatorname{Aut}} % 自己同型群, または自己同型写像の全体
\newcommand{\grcharacter}{\operatorname{Char}} % 指標
\newcommand{\grdiff}{\operatorname{Diff}} % 微分同相群
\newcommand{\grdual}[1]{\widehat{#1}} % 双対群
\newcommand{\greuctr}[2]{\fldreal_{#1}^{#2}} % ユークリッド空間の空間並進群
\newcommand{\grind}[2]{\sqbk{#1 : #2}} % 群の指数, expedition0005326
\newcommand{\grisometry}{\operatorname{Isom}} % 等長変換群
\newcommand{\grmobius}{\mathop{\text{Möb}}} % メビウス変換
\newcommand{\grord}{\operatorname{ord}} % 群の位数, expedition0005646
\newcommand{\grprod}[1]{#1^{\txtexcludezero}} % 体の零元を除外した乗法群
\newcommand{\grprufer}[1]{\setquot{\ringpolynomial{\ringratint}{\frac{1}{#1}}}{\ringratint}} % プリューファーp群
\newcommand{\grringaug}{\mathrm{d}} % 群環の添加写像, expedition0008765
\newcommand{\grring}[2]{#1 \! \sqbk{#2}} % 群環
\newcommand{\grtransposition}[2]{(#1 \, #2)} % 群の互換
\newcommand{\grtrpos}[2]{(#1 \, #2)} % 互換, expedition0009437
\newcommand{\gtbd}{\partial} % 位相空間に対する境界作用素
\newcommand{\gtclosrbk}[1]{\overline{\rbk{#1}}} % 一般位相の閉包
\newcommand{\gtclos}[1]{\overline{#1}} % 一般位相の閉包
\newcommand{\gtfmlcompactset}[1]{\mathcal{#1}} % コンパクト集合がなす族, expedition0000329
\newcommand{\gtfmlcompactconvexset}[1]{\mathcal{#1}} % コンパクト凸集合がなす族, expedition0010532
\newcommand{\gtint}{\operatorname{Int}} % 一般位相の開核, open kernel
\newcommand{\gtoneptcpt}[1]{\dot{#1}} % 一点コンパクト化
\newcommand{\gtopclos}{\operatorname{clos}} % 一般位相の閉包, closure
\newcommand{\gttopsorgenfreyline}{\mathcal{S}} % ソルゲンフライ直線, expedition0001888
\newcommand{\habd}{\partial} % ホモロジー代数での境界作用素
\newcommand{\haboundary}{B} % 境界輪体加群
\newcommand{\hacechcohomology}{\check{H}} % チェックコホモロジー
\newcommand{\hachaincomplex}[1]{#1_{\bullet}} % 鎖複体
\newcommand{\hachain}{C} % 鎖加群
\newcommand{\hacoboundary}{B} % 余境界輪体
\newcommand{\hacochain}{C} % 余鎖
\newcommand{\hacocycle}{Z} % 余輪体
\newcommand{\hacohomology}{H} % コホモロジー
\newcommand{\hacwbd}{\underline{\partial}} % CW対に対する鎖加群上の境界作用素, expedition0007497
\newcommand{\hacycle}{Z} % 輪体加群
\newcommand{\hahomology}{H} % ホモロジー
\newcommand{\hainflation}{\operatorname{Inf}} % 膨張写像, inflation map, expedition0009331
\newcommand{\hareducedhomology}{\tilde{H}} % 被約ホモロジー, expedition0008655
\newcommand{\harelativechain}{\underline{C}} % 相対鎖加群
\newcommand{\harelativecohomology}{\underline{H}} % 相対コホモロジー
\newcommand{\harelativehomology}{\underline{H}} % 相対ホモロジー
\newcommand{\harelbd}{\overline{\partial}} % ホモロジー代数での相対境界作用素, expedition0007469
\newcommand{\harestriction}{\operatorname{res}} % 主にホモロジー代数での制限写像, 必要に応じてそれ以外にも転用
\newcommand{\htpyrevrbk}[1]{\rbk{#1}^{-}} % 道の逆, expedition0003263
\newcommand{\htpyrev}[1]{#1^{-}} % 道の逆, expedition0003263
\newcommand{\intdef}[3]{\sqbk{#1}_{#2}^{#3}} % 定積分, definite integral
\newcommand{\intleb}{\mathrm{L} \int} % ルベーグ積分: リーマン積分と区別したいときに
\newcommand{\intpv}{\oppv \int} % 主値積分
\newcommand{\intriem}{\mathrm{R} \int} % リーマン積分, ルベーグ積分と記号をわけたいときに
\newcommand{\laantilinhodgestar}{\mathop{\bar{\star}}} % 反線型ホッジ作用素, expedition0003676
\newcommand{\labilinformrank}{\operatorname{rank}} % (有限次元空間上の)双線型形式の階数
\newcommand{\labilin}{\operatorname{Bilin}} % 双線型写像の集合
\newcommand{\lacomplexication}[1]{#1_{\fldcmp}} % 複素化
\newcommand{\laextalg}{\mathcal{A}} % 外積代数
\newcommand{\laextdiff}[1]{\mathop{d #1}} % 外微分
\newcommand{\laextprod}{\wedge} % 外積
\newcommand{\laherm}{\operatorname{Herm}} % エルミート行列の全体
\newcommand{\lahodgedualind}[1]{#1^{\lahodgestar}} % ホッジ双対指数, expedition0003574
\newcommand{\lahodgeminkowskistar}{\star_M} % ホッジ-ミンコフスキー作用素
\newcommand{\lamat}[2]{\rbk{#1}_{#2}} % 行列の成分表示
\newcommand{\lapartialtp}[2]{#2^{T_{#1}}} % 部分転置, expedition0007073
\newcommand{\lapfaffian}{\operatorname{Pf}} % パフィアン
\newcommand{\larank}{\operatorname{rank}} % 線型代数での行列・線型写像の階数
\newcommand{\lasetbase}[1]{\mathcal{#1}} % 基底をなすベクトルの組 \lasetbase{E} = \basebk{e_i}_{i=1}^{N}
\newcommand{\latodf}{\operatorname{\flat}} % ベクトル場を微分形式に変換, expedition0003743
\newcommand{\latovf}{\operatorname{\sharp}} % 音楽同型, 微分形式をベクトル場に送る, expedition0003743, expedition0009967
\newcommand{\latprbk}[1]{\latp{\rbk{#1}}} % 括弧をつけた上での転置
\newcommand{\lgconstsymbol}[1]{\mathcal{#1}} % 定数記号, expedition0010147
\newcommand{\lgfnsymbol}[1]{\mathcal{#1}} % 関数記号, expedition0010147
\newcommand{\liealgheisenberg}[1]{\mathfrak{H}_{#1}} % ハイゼンベルクリー環
\newcommand{\liealgrank}{\operatorname{rank}} % (半単純)リー環の階数, expedition0005268
\newcommand{\liealguppertriangle}[1]{\mathfrak{T}_{#1}} % 上三角リー環
\newcommand{\liealg}[1]{\mathfrak{#1}} % リー環
\newcommand{\liegr}[1]{\mathrm{#1}} % リー群
\newcommand{\lorentz}{L} % ローレンツ空間, expedition0010365
\newcommand{\losetpredicate}{\mathcal{P}} % 述語記号の全体, expedition0010147
\newcommand{\lpae}{\mathop{\mathrm{a.e.}}} % ほとんどいたるところ, 「.」を含むためoperatornameが使えない
\newcommand{\lpas}{\mathop{\txtprobas}} % ほとんど確実に, 「.」を含むためoperatornameが使えない
\newcommand{\lpconjindex}[1]{#1'} % 共役指数, expedition0010223
\newcommand{\lpconv}{\ast} % たたみ込み
\newcommand{\lpfn}{\mathcal{L}} % 同値類を取らない生の関数に対するルベーグ空間
\newcommand{\lpseq}{\ell} % 数列空間, expedition0000444
\newcommand{\lpsimple}{\mathrm{SF}} % 単関数の空間
\newcommand{\lpsymrearr}[1]{#1^{\ast}} % 対称再配分, expedition0003054, expedition0003059
\newcommand{\lpvanishingatinfty}{D} % Lieb-Loss, 堀内, 無限遠で消滅的な関数, expedition0001998
\newcommand{\lp}{L} % ルベーグ空間, L_{\txtloc}^p(X), L^p(X)
\newcommand{\manalgvecfld}{\mathfrak{X}} % ベクトル場がなす代数
\newcommand{\manatlas}[1]{\mathcal{#1}} % 多様体のアトラス
\newcommand{\manclifford}{\operatorname{Cl}} % クリフォード関係, クリフォード束: expedition0005779
\newcommand{\mancospherebdle}{\dual{S}} % 余球面束
\newcommand{\mancotfunctor}{\dual{T}} % 余接束に関わる関手
\newcommand{\mancotpqfunctor}[1]{T^{\ast,#1}} % (p,q)テンソル積束に関わる関手
\newcommand{\mancrosssection}{\Gamma} % ベクトル束の切断
\newcommand{\mandgspinc}{\operatorname{Spin}^{\mathrm{c}}} % スピン$c$多様体
\newcommand{\mandgspin}{\operatorname{Spin}} % スピン多様体
\newcommand{\mandiffideal}{\mathcal{I}} % 微分形式がなす微分イデアル, expedition0007338
\newcommand{\mandisk}{\mathbb{D}} % 円板または閉球体
\newcommand{\mandistribution}{\mathcal{D}} % 多様体上の分布, expedition0003991
\newcommand{\manembedding}{\operatorname{Emb}} % 埋め込み, expedition0009279
\newcommand{\manfnspmorse}{\operatorname{Morse}} % モース関数の空間, expedition0009964
\newcommand{\mangrassmann}[3]{\fun{G_{#1,#2}}{#3}} % グラスマン多様体
\newcommand{\mangrholonomy}{\operatorname{Hol}} % ホロノミー群, expedition0009279
\newcommand{\mangrpicard}{\operatorname{Pic}} % ピカール群, expedition0004339
\newcommand{\manhirzebruchlgenus}{\mathcal{L}} % ヒルツェブルフのL種数, expedition0010182
\newcommand{\manimbedding}{\operatorname{Imb}} % はめ込み, expedition0009279
\newcommand{\manjacobi}{\operatorname{Jac}} % ヤコビ多様体
\newcommand{\manmaprank}{\operatorname{rank}} % 多様体上の写像の階数, expedition0002606
\newcommand{\mannormalfunctor}{\nu} % 法空間・法束に関わる関手
\newcommand{\manopextalmostcmpstructure}{\operatorname{\mathbf{I}}} % 定義\ref{expedition0003752}(2)の作用素
\newcommand{\manproj}{\mathbb{P}} % 微分多様体論での射影空間
\newcommand{\manriemsphere}{\overline{\fldcmp}} % リーマン球面
\newcommand{\mansetcrit}{\operatorname{Crit}} % 臨界点の集合, expedition0003231, expedition0007521, expedition0009963
\newcommand{\mansphere}{\mathbb{S}} % 球面
\newcommand{\manspherebdle}{S} % 球面束
\newcommand{\manspmoduli}{\mathcal{M}} % モジュライ空間
\newcommand{\manstiefel}[3]{\fun{V_{#1,#2}}{#3}} % シュティフェル多様体
\newcommand{\mansubmersion}{\operatorname{Subm}} % 沈め込み, expedition0002644
\newcommand{\mantanfunctor}{T} % 接束用の関手
\newcommand{\mantanvec}[2]{\fnvalued{\rbk{\oppd{#1}}}{#2}} % 接ベクトル
\newcommand{\mantensorbdle}[1]{\mathcal{#1}} % テンソル積束
\newcommand{\mantoddgenus}{\operatorname{Todd}} % トッド種数, expedition0009957
\newcommand{\mantorus}{\mathbb{T}} % トーラス
\newcommand{\mantrivbdle}[1]{\underline{#1}} % 自明束
\newcommand{\manvecfld}[1]{#1} % ベクトル場
\newcommand{\manvectbdle}{\operatorname{Vect}} % ベクトル束
\newcommand{\manyangmills}{\operatorname{YM}} % ヤン-ミルズ, expedition0006094
\newcommand{\mblclos}[1]{\overline{#1}} % 測度空間の完備化, あえて「閉包」を利用
\newcommand{\mblfmlbaire}{\mblfml{B}} % ベール集合族, expedition0006109
\newcommand{\mblfmlborelext}{\overline{\mblfml{B}}} % 拡大実数体上のボレル集合族
\newcommand{\mblfmlborel}{\mblfml{B}} % ボレル集合族
\newcommand{\mblfmldsys}{\mblfml{D}} % 可測空間論での$d$-系, Lambda系-とも言う, expedition0000538
\newcommand{\mblfmlfigure}{\mblfml{F}_0} % 区間塊の全体, expedition0000642
\newcommand{\mblfmlfinite}[1]{\mblfml{#1}} % 有限加法族
\newcommand{\mblfmlfrak}[1]{\mathfrak{#1}} % 加法族の mathfrak 表記
\newcommand{\mblfmlleb}{\mblfml{F}_{\txtleb}} % ルベーグ可測集合の全体
\newcommand{\mblfmlleftopeninterval}{\mblfml{I}} % ユークリッド空間の左開区間全体がなす集合, expedition0000887
\newcommand{\mblfmlnull}{\mblfml{N}} % 零集合の全体
\newcommand{\mblfmlpisys}[1]{\mblfml{#1}} % 可測空間論での$\pi$-系
\newcommand{\mblfmltail}[1]{\mblfml{#1}} % 末尾加法族, expedition0004740
\newcommand{\mblfml}[1]{\mathcal{#1}} % 可測集合族に関わる記号の基礎
\newcommand{\mblfn}{M} % 可測関数
\newcommand{\mblsplebwiener}{\mblfml{W}} % ルベーグ-ウィーナー空間, expedition0005039
\newcommand{\modbimod}[3]{\setSymbolDownLeft{#2}{#1}_{#3}} % 両側加群
\newcommand{\modflatdim}{\operatorname{fd}} % 加群の平坦次元
\newcommand{\modgldim}{\operatorname{\operatorname{gl \mathchar`- dim}}} % 加群の大域次元
\newcommand{\modinjdim}{\operatorname{id}} % 加群の単射次元
\newcommand{\modleft}[2]{\setSymbolDownLeft{#2}{#1}} % 左加群
\newcommand{\modlgldim}{\operatorname{l \mathchar`- gl \mathchar`- dim}} % 左大域次元
\newcommand{\modprojdim}{\operatorname{pd}} % 加群の射影次元
\newcommand{\modrank}{\operatorname{rank}} % 加群の階数
\newcommand{\modrgldim}{\operatorname{r \mathchar`- gl \mathchar`- dim}} % 加群の右大域次元
\newcommand{\modright}[2]{#1_{#2}} % 右加群
\newcommand{\modsocle}{\operatorname{Soc}} % 半単純成分, socle, expedition0008225
\newcommand{\modwgldim}{\operatorname{w \mathchar`- gl \mathchar`- dim}} % 加群の弱大域次元
\newcommand{\monfree}{\operatorname{Free}} % 集合から生成する自由モノイド, expedition0009945
\newcommand{\monnat}{\mathbb{N}} % 0を含めた自然数全体がなすモノイド
\newcommand{\msrae}[1]{#1 \hyphen \mathrm{a.e.}} % ある測度に対するほとんどいたるところ
\newcommand{\msras}[1]{#1 \hyphen \txtprobas} % ある測度に対するほとんど確実に
\newcommand{\msrbb}[1]{\mathbb{#1}} % mathbbで表す測度
\newcommand{\msrcal}[1]{\mathcal{#1}} % mathcalで表す測度
\newcommand{\msrcnt}{\mu_{\txtcnt}} % 数え上げ測度
\newcommand{\msrdirac}{\mu_{\txtdirac}} % ディラックのデルタ測度
\newcommand{\msrlaw}{\mathrm{Law}} % 法則
\newcommand{\msrlebwiener}{\mu_0} % ルベーグ-ウィーナー測度
\newcommand{\msrleb}{\mu_{\txtleb}} % ルベーグ測度
\newcommand{\msrouter}[1]{#1^{\ast}} % 外測度
\newcommand{\msrprb}{\mathrm{Pr}} % たまに役立つであろう, 立体のPrで表す確率測度
\newcommand{\msrpreleb}{m_{\txtleb}} % ユークリッド空間上のルベーグ測度の定義に利用する有限加法的測度
\newcommand{\msrsf}[1]{\mathsf{#1}} % mathsfで表す測度
\newcommand{\msrzero}{\mu_0} % 何らかの意味で基準になる測度
\newcommand{\msr}[1]{#1} % 測度の識別用
\newcommand{\nafldcmp}{\na{\fldcmp}} % 超準複素数体
\newcommand{\nafldrat}{\na{\fldrat}} % 超準有理数体
\newcommand{\nafldreal}{\na{\fldreal}} % 超準実数体
\newcommand{\nainftysmall}[1]{d #1} % 超準解析での無限小, それ以外の無限小への流用は避ける
\newcommand{\namonnat}{\na{\monnat}} % 超準自然数がなすモノイド
\newcommand{\namsrloeb}{\mu_{\mathrm{Loeb}}} % 超準解析でのローブ測度
\newcommand{\nanonstduniv}{\nauniv(\na{\fldreal})} % 超準宇宙, expedition0009306
\newcommand{\napreinternaluniv}{\tilde{\nauniv}} % 超準解析での宇宙の一種
\newcommand{\narank}{\operatorname{rank}} % 超準解析での要素の階数
\newcommand{\nasplpbase}{\mathfrak{F}} % 定義expedition0009837で定めた商空間
\newcommand{\nastdop}{\operatorname{st}} % 超準解析の標準化写像, expedition0005133
\newcommand{\nastdpart}[1]{\setSymbolUpLeft{\circ}{#1}}
\newcommand{\nastdset}[1]{\setSymbolUpLeft{\txtstandard}{#1}} % expedition0005127, 標準化された集合
\newcommand{\nastduniv}{\nauniv(\fldreal)} % 超準解析での標準宇宙, expedition0005110
\newcommand{\nauniv}{\mathcal{U}} % 超準解析での標準宇宙を表す記号
\newcommand{\na}[1]{\setSymbolUpLeft{\ast}{#1}} % 超準化作用素
\newcommand{\nfoldvar}[2]{\underline{#1}_{#2}} % #2個の変数をまとめて下線つきで表す
\newcommand{\oaadjustedentropy}{\hat{S}} % 修正エントロピー expedition0009609, 修正相対エントロピー expedition0009878
\newcommand{\oaalgentropy}{H} % 代数的エントロピー, expedition0009940
\newcommand{\oaalgweyl}{\oa{\opfockweyl}} % ワイル代数
\newcommand{\oaarakiwoodsrepn}{\oarepn^{\txtarakiwoods}} % 荒木-ウッズ表現
\newcommand{\oaarakiwoodssp}{\spfock_{\txtarakiwoods}} % 荒木-ウッズ空間
\newcommand{\oaarakiwoodsvac}{\Omega_{\txtarakiwoods}} % 荒木-ウッズ自由場の真空
\newcommand{\oaarakiwoodsweyl}[1]{W_{#1}} % 荒木-ウッズ表現でのワイル作用素
\newcommand{\oacalkin}{\operatorname{Cal}} % カルキン環
\newcommand{\oacar}{\operatorname{CAR}} % CAR環
\newcommand{\oacenter}{Z} % 作用素環の中心
\newcommand{\oacntentropy}{h} % CNTエントロピー, expedition0009938
\newcommand{\oacommutant}[1]{#1^{\prime}} % 作用素環での可換子環
\newcommand{\oacondexp}{\operatorname{E}} % 作用素環に対する条件つき期待値
\newcommand{\oacstar}{C^{\ast}} % $C^{\ast}$-環
\newcommand{\oaderiv}{\delta} % 作用素環上の微分, derivation
\newcommand{\oadoublecommutant}[1]{#1^{\prime \prime}} % 二重可換子環
\newcommand{\oaentropy}{S} % 作用素環上のエントロピーで相対エントロピー(定義\ref{expedition0009684})にも流用, オリジナルはフェルミオンの格子模型, expedition0009942
\newcommand{\oalocgibbsstate}{\phi^{c}} % 局所ギブス状態, expedition0009692
\newcommand{\oameanadjustedentropy}{\hat{s}} % 平均修正エントロピー expedition0009609
\newcommand{\oameanenergy}{e} % 平均エネルギー
\newcommand{\oameanentropy}{s} % 平均エントロピー expedition0009684
\newcommand{\oanormalizefunctional}[1]{\sqbk{#1}} % 作用素環での汎関数の正規化の記号, どれだけ一般的かは不明
\newcommand{\oaopdualcone}[1]{\check{#1}} % 双対錐の記号
\newcommand{\oapredual}[1]{#1_{\ast}} % フォン・ノイマン環の前双対
\newcommand{\oaresolventalgebra}{\oa{R}} % レゾルベント環
\newcommand{\oaresolvent}{R} % レゾルベント環の要素としてのレゾルベント
\newcommand{\oaselfdualcone}{\mathcal{P}} % 標準表現に付随する自己双対錐
\newcommand{\oasetregularrepn}{\operatorname{Reg}} % 作用素環の正則表現
\newcommand{\oasetsolforvariationalproblem}[1]{\Lambda_{#1}} % 変分原理の解の集合, expedition0009854
\newcommand{\oaspec}{\operatorname{Spec}} % 可換な$\oacstar$-環に対するスペクトル
\newcommand{\oasplocham}{\mathop{\boldsymbol{H}}} % (格子模型での)標準ハミルトニアンがなす実線型空間
\newcommand{\oaspstdpot}{\mathcal{P}} % (格子模型での)標準ポテンシャルがなす実線型空間, expedition0009673
\newcommand{\oastaralgebra}{\operatorname{\ast \hyphen \mathrm{alg}}} % 指定した集合が生成する$\ast$-環
\newcommand{\oatomitamodaut}{\sigma} % 冨田のモジュラー自己同型写像
\newcommand{\oatomitamodconj}{J} % 冨田のモジュラー共役作用素
\newcommand{\oatomitamodop}{\Delta} % 冨田のモジュラー作用素
\newcommand{\oaweyl}{\oa{W}} % 抽象的なワイル環
\newcommand{\oawstar}{W^{\ast}} % $W^{\ast}$-環
\newcommand{\oa}[1]{\mathcal{#1}} % 作用素環を表す記号用
\newcommand{\opabscontipart}[1]{#1_{\txtabsconti}} % 作用素#1の絶対連続部分, expedition0009300
\newcommand{\opalgebraicmultiplicity}[1]{\fun{m_{\txtalg}}{#1}} % 代数的多重度, expedition0010176
\newcommand{\opawsegal}{\phi} % 荒木-ウッズ表現でのシーガルの場の作用素
\newcommand{\opbottomofessspec}[1]{\Sigma_{#1}} % 真性スペクトルの下限
\newcommand{\opcharnum}[2]{\fun{\mu_{#1}}{#2}} % 作用素論でのn番目の特性数(特性レベル), expedition0009754
\newcommand{\opclos}[1]{\overline{#1}} % 作用素論での作用素に対する閉包
\newcommand{\opdifference}[1]{\mathrm{#1}} % 差分作用素, expedition0004874
\newcommand{\opdifficiencyind}{n} % 対称作用素の不足指数, expedition0009898
\newcommand{\opdiracind}{\operatorname{ind}} % ディラック作用素の指数, expedition0009280
\newcommand{\opdmsrfeynman}[1]{\mathop{\mathcal{D} #1}} % 経路積分での形式的な測度
\newcommand{\opdmsr}[1]{\mathop{d #1}} % 積分で現れる測度用のコマンド, dxのような単純な対象も含む
\newcommand{\opdpathintegralmsr}[1]{\mathop{\mathcal{D} #1}} % 物理的な経路積分で現れる「測度」
\newcommand{\opfnresolvent}[1]{\invrbk{#1}} % レゾルベント関数
\newcommand{\opfockccrrepresentationcran}{\mathrm{CCR} \hyphen \mathrm{CrAn}_{\txtfock}} % 非有界フォック正準交換関係環 expedition0011291
\newcommand{\opfockschwingerfunctional}{S} % シュウィンガー関数, expedition0009992
\newcommand{\opfocksegalradiation}{A} % 輻射場に対するシーガルの場の作用素, expedition0010034
\newcommand{\opfocksndqntdiff}{d \Gamma} % 第一種の第二量子化作用素, 微分版
\newcommand{\opfocksndqnt}{\Gamma} % 第二種の第二量子化作用素
\newcommand{\opfockvac}{\Omega} % フォック真空
\newcommand{\opfockweyl}{W} % フォック空間上のワイル作用素
\newcommand{\opfockwightmanfunctional}{W} % ワイトマン関数, expedition0009986
\newcommand{\opformclos}[1]{\overline{#1}} % 準双線型形式の閉包,
\newcommand{\opformdomain}{\mathop{Q}} % 準双線型形式の定義域, expedition0009907
\newcommand{\opforminf}[1]{\fun{\mathcal{E}_0}{#1}} % 準双線型形式の下限, expedition0009908
\newcommand{\opformsum}{\operatorname{\dot{+}}} % 作用素論での形式和
\newcommand{\opform}[1]{\mathsf{#1}} % 二次形式
\newcommand{\opfredholmind}{\operatorname{ind}} % フレドホルム指数
\newcommand{\opgeometricmultiplicity}[1]{\fun{m_{\txtgeom}}{#1}} % 幾何学的多重度, expedition0010176
\newcommand{\opgroundstateenergy}{E_0} % 基底エネルギー, expedition0001585
\newcommand{\opprojto}[1]{\sqbk{#1}} % 閉部分空間への射影作用素または閉部分空間それ自身
\newcommand{\oppr}{\mathrm{pr}} % 射影作用素
\newcommand{\oprot}{\mathcal} % 回転の作用素・行列で文字種を変えたい場合に利用
\newcommand{\opsetresolvent}[1]{\fun{\rho}{#1}} % レゾルベント集合
\newcommand{\opsingularpart}[1]{#1_{\txtsingular}} % 作用素#1の特異部分空間, expedition0009300
\newcommand{\opspabsconti}[2]{\fun{#1_\txtabsconti}{#2}} % ヒルベルト空間#1での作用素#2の絶対連続部分空間, expedition0009913
\newcommand{\opspbddlin}[1]{\fun{\mathbb{B}}{#1}} % 有界線型作用素がなす空間
\newcommand{\opspecfml}[1]{\fml{#1_{\lambda}}{\lambda \in \fldreal}} % スペクトル族
\newcommand{\opspecint}[1]{\mathcal{E}} % スペクトル積分, expedition0001621
\newcommand{\opsplin}{\mathbb{L}} % 非有界作用素まで含めた線型作用素がなす空間
\newcommand{\opsppseudodiffop}{\mathrm{PDO}} % 次数$m$以下の偏微分作用素の空間, expedition0009282
\newcommand{\opsprollnikpotential}{\mathcal{R}} % ロルニックポテンシャルの全体, expedition0003113
\newcommand{\opspsingular}[2]{\fun{#1_{\txtsingular}}{#2}} % ヒルベルト空間#1での作用素#2の特異部分空間, expedition0009913
\newcommand{\opsympform}[2]{\fun{\sigma}{#1, \, #2}} % ボソン$\oacstar$-環に付随するシンプレクティック形式
\newcommand{\optoop}[1]{\widehat{#1}} % 関数などを作用素に変換する写像
\newcommand{\opvnentropy}{S} % フォン・ノイマンエントロピー
\newcommand{\padic}[2]{#1_{#2}} % p進整数, p進数, l進数などを表す
\newcommand{\physadj}[1]{#1^{\dagger}} % 形式的共役または物理でのエルミート共役
\newcommand{\physapparatus}[1]{\mathbf{#1}} % 小澤論文でapparatusを表す記号
\newcommand{\physbfposition}[1]{\mathbf{#1}} % 小澤論文でmathbfで位置を表す記号
\newcommand{\physbkt}[2]{\left\langle #1 \vert #2 \right\rangle} % 量子力学でよく使われる内積の記法
\newcommand{\physboltzmannconst}{k_{\txtboltzmann}} % ボルツマン定数, k_B
\newcommand{\physcmpconj}[1]{{#1}^{\ast}} % 物理記法の複素共役
\newcommand{\physelecricfield}{E} % 電場
\newcommand{\physenergyfunc}{\mathcal{E}} % エネルギー汎関数
\newcommand{\physfermisurface}{\mathcal{F}} % 連結グリーン関数, expedition0009492
\newcommand{\physfnconnectedgreen}{\mathcal{G}} % 連結グリーン関数, expedition0009492
\newcommand{\physfngeneratorconnectedgreen}{\mathcal{C}} % 連結グリーン関数の生成子, expedition0009492
\newcommand{\physionization}{\mathcal{I}} % イオン化, expedition0010208
\newcommand{\physmagneticflux}{\Phi} % 物理, 磁束
\newcommand{\physmean}[1]{\dbk{#1}} % 物理, 特に統計力学によく出てくる平均の記号
\newcommand{\physmfdensity}{B} % 磁束密度, magnetic flux density
\newcommand{\physoperator}[1]{\hat{#1}} % 物理で演算子につけるハット
\newcommand{\physoptimeorder}{\operatorname{T}} % 時間順序作用素
\newcommand{\physplanckconst}{h} % プランク定数
\newcommand{\physpoissonbkt}[1]{\cbk{#1}} % 物理用のポアソン括弧
\newcommand{\physprobe}[1]{\mathbf{#1}} % 小澤論文でprobeを表す記号
\newcommand{\physqinfomatw}{\mathsf{W}} % よい名前がわからない・思いつかないため暫定
\newcommand{\physscalarpot}{\phi} % スカラーポテンシャル
\newcommand{\physsysbf}[1]{\mathbf{#1}} % 量子系などの系をmathbfで表したい場合に利用
\newcommand{\physsyscal}[1]{\mathcal{#1}} % 量子系などの系をmathcalで表したい場合に利用
\newcommand{\physsysrm}[1]{\mathrm{#1}} % 量子系などの系をmathrmで表したい場合に利用
\newcommand{\physvecpot}{A} % ベクトルポテンシャル
\newcommand{\physvec}[1]{\boldsymbol{#1}} % 物理で利用する太字のベクトル
\newcommand{\postcomp}[1]{{\vphantom{\circ}}_{#1}{\circ}} % 左から合成する, TODO 他と適切に統合, 特にlusharp, rusharp
\newcommand{\prbabsconti}{\ll} % 測度の絶対連続性
\newcommand{\prbbinshannonentropy}{h} % 二値シャノンエントロピー, 二値エントロピー, expedition0006915
\newcommand{\prbcov}{\mathrm{Cov}} % 共分散
\newcommand{\prbdbkquadvar}[1]{\dbk{#1}} % 山括弧の二次変分
\newcommand{\prbentropy}{H} % エントロピー
\newcommand{\prbexp}{\mathbb{E}} % 期待値, 条件つき期待値にも応用
\newcommand{\prbfmlmartingale}[1]{\mathcal{#1}} % マルチンゲールの族
\newcommand{\prbfmlprbmeas}[1]{\mathcal{#1}} % 確率測度の族
\newcommand{\prbgaussiansuperprocesssubsp}{\Gamma} % ガウス超過程に付随する部分空間, expedition0004961
\newcommand{\prbinvfml}[1]{\mathcal{#1}} % ある可測写像に対する不変加法族, expedition0012042
\newcommand{\prbkullbackleibler}{D} % カルバック-ライブラ情報量
\newcommand{\prbmean}[1]{\left\langle #1 \right\rangle} % 平均の別記法
\newcommand{\prbou}{\xi} % オルンシュタイン-ウーレンベック過程
\newcommand{\prbsetprbmeas}{\operatorname{Prob}} % 確率測度の空間
\newcommand{\prbsetregprbmeas}[1]{\fun{\txtprob_{\txtregular}}{#1}} % 正則な確率測度の空間
\newcommand{\prbshannonentropy}{H} % シャノンエントロピー, expedition0006895
\newcommand{\prbsingular}{\perp} % 絶対連続性に対する測度の特異性
\newcommand{\prbsqbkquadvar}[1]{\sqbk{#1}} % 角括弧の二次変分, expedition0010000
\newcommand{\prbstddev}{\Delta} % 標準偏差, expedition0006478
\newcommand{\prbvar}{\operatorname{Var}} % 分散, expedition0006478
\newcommand{\prbweakconvmeasure}{\Rightarrow} % 測度の弱収束, expedition0004769
\newcommand{\precomp}[1]{\circ_{#1}} % 右から合成する, TODO 他と適切に統合, 特にlusharp, rusharp
\newcommand{\prodthree}[3]{\mathop{\mathop{\prod_{#1}}_{#2}}_{#3}} % 三条件を持つ積
\newcommand{\prodtwo}[2]{\mathop{\prod_{#1}}_{#2}} % 二条件を持つ積
\newcommand{\pshconst}[1]{\setSymbolUpLeft{p}{#1}} % sheaf 定数前層, expedition0008790
\newcommand{\pshdirectimage}[1]{#1_{p}} % sheaf 前層の順像, expedition0008793
\newcommand{\pshinverseimage}[1]{#1^{p}} % 前層の逆像, {#1^{p}}ではなく{#1^{\vee}}で表す場合もある
\newcommand{\pshnatural}[1]{#1^{\natural}} % sheaf 層の議論で現れる
\newcommand{\pullbackrbk}[1]{\pullback{\rbk{#1}}} % 引き戻し
\newcommand{\pullback}[1]{#1^{\ast}} % 引き戻し
\newcommand{\pushoutrbk}[1]{\rbk{#1}_{\ast}} % 押し出し
\newcommand{\pushout}[1]{#1_{\ast}} % 押し出し, 共通
\newcommand{\qmebit}{\operatorname{ebit}} % エンタングルメントビット, expedition0009268
\newcommand{\qmlogneg}[1]{\fun{\mathrm{LogNeg}}{#1}} % 対数ネガティビティ関数
\newcommand{\qmpaulispin}{\sigma} % パウリ行列
\newcommand{\qmrenyientropy}{\psi} % 相対レニーエントロピー, expedition0007123
\newcommand{\qmschmidtnumber}{\operatorname{Sch}} % シュミット数, expedition0009985
\newcommand{\qmtfadmissble}{\mathcal{C}} % トーマス-フェルミ汎関数の議論で許される関数のクラス, expedition0009357周辺
\newcommand{\qtdecoder}{\mathcal} % 復号器, expedition0007131
\newcommand{\qtinstrument}{\mathcal{I}} % 量子測定理論のインストルメント
\newcommand{\qtobservable}{\mathcal} % 可観測量または自己共役作用素
\newcommand{\qtpovm}{\mathcal} % POVM, expedition0006582
\newcommand{\qtrollnik}{\mathcal} % ロルニック, expedition0003113
\newcommand{\qtsetmeasuredvalue}[1]{\mathcal{#1}} % 測定値の集合, expedition0006499
\newcommand{\qtspobservable}{\mathcal} % 可観測量または自己共役作用素がなす空間
\newcommand{\qtstatevec}[1]{#1} % 状態ベクトル, 波動関数
\newcommand{\qtstate}[1]{#1} % (作用素環上の状態の意味での)状態
\newcommand{\riemzeta}{\zeta} % リーマンのゼータ, リーマンの$\zeta$
\newcommand{\ringannideal}{\operatorname{Ann}} % 零化イデアル, expedition0007946
\newcommand{\ringassociatedprime}{\operatorname{Ass}} % 加群の随伴素イデアル全体 expedition0009267
\newcommand{\ringconvseries}[2]{\sqfun{#1}{#2}} % 収束ベキ級数環, expedition0004584
\newcommand{\ringfmlideals}[1]{\mathcal{#1}} % イデアル全体がなす集合
\newcommand{\ringfmlmaximalideals}[1]{\mathcal{#1}} % 極大イデアル全体がなす集合
\newcommand{\ringformalpowerseries}[2]{#1 \! \sqbk{\sqbk{#2}}} % 形式的ベキ級数環
\newcommand{\ringfracfld}{\operatorname{Frac}} % 分数体・剰余体, expedition0008559l
\newcommand{\ringideal}[1]{\mathsf{#1}} % 環のイデアル, 各イデアルを表す記号でイデアルの全体は \ringfmlideals
\newcommand{\ringlocalization}[2]{\inv{#1} #2} % 環の局所化, expedition0009953
\newcommand{\ringmaxspec}{\mathrm{m} \hyphen \mathrm{Spec}} % 極大スペクトル
\newcommand{\ringnilrad}{\operatorname{nilrad}} % べき零根基, expedition0008345
\newcommand{\ringopposite}[1]{{#1}^{\txtopposite}} % 反対環に利用.
\newcommand{\ringpolynomial}[2]{\sqfun{#1}{#2}} % 多項式環
\newcommand{\ringprim}{\operatorname{Prim}} % 原始イデアル
\newcommand{\ringrad}{\operatorname{rad}} % 根基, expedition0004618
\newcommand{\ringratint}{\mathbb{Z}} % 有理整数環
\newcommand{\ringregelems}[1]{#1^{\circ}} % 環の正則元全体がなす集合,expedition0008368
\newcommand{\ringspec}{\operatorname{Spec}} % 可換環のスペクトル, expedition0009966
\newcommand{\semigrposintverone}{\ringratint_{>0}} % 正の整数の全体の表記のバリエーション
\newcommand{\semigrposint}{\monnat_1} % 正の整数の全体または1以上の自然数全体
\newcommand{\seqjone}[1]{\seq{#1_{j}}{j \in \semigrposint}}
\newcommand{\seqj}[1]{\seq{#1_{j}}{j \in \monnat}} %
\newcommand{\seqkone}[1]{\seq{#1_{k}}{k \in \semigrposint}} % 添字を正の整数の$k$とする点列用のエイリアス
\newcommand{\seqk}[1]{\seq{#1_{k}}{k \in \monnat}} % 添字を自然数の$k$とする点列用のエイリアス
\newcommand{\seqna}[1]{\seq{#1_n}{n \in \namonnat}} % 超準自然数に対する点列
\newcommand{\seqnone}[1]{\seq{#1_n}{\ninone}}
\newcommand{\seqn}[1]{\seq{#1_{n}}{\nin}} % 添字を自然数の$n$とする点列用のエイリアス
\newcommand{\seq}[2]{\if\relax\detokenize{#1}\relax \rbk{#1} \else \rbk{#1}_{#2} \fi} % 点列の基本系
\newcommand{\setcardopsharp}[1]{\operatorname{\#} #1} % ある集合に対して濃度を与える作用素：#版
\newcommand{\setcardop}[1]{\operatorname{card} #1} % ある集合に対して濃度を与える作用素
\newcommand{\setcard}[1]{\mathfrak{#1}} % 濃度自体を表す
\newcommand{\setcenter}[3]{\fun{Z_{#1}^{#2}}{#3}} % 中心, 群・環など様々代数的対象に対して適用する
\newcommand{\setcmpconj}[1]{#1^{\dagger}} % 集合の複素共役, expedition0010591
\newcommand{\setcpl}[1]{#1^{c}} % 補集合, complement
\newcommand{\setcube}{\mathcal{C}} % 超立方体, expedition0009745
\newcommand{\setdecomp}{\operatorname{Decomp}} % 適当な分解の集合
\newcommand{\seteven}{\mathrm{Even}} % 偶数全体がなす集合
\newcommand{\setextremal}{\operatorname{ex}} % 端点の全体
\newcommand{\setirr}{\mathbb{I}} % 無理数全体がなす集合
\newcommand{\setisomorphism}[1]{\operatorname{Iso}} % 同型写像の集合
\newcommand{\setlowerbounds}[1]{#1^{\downarrow}} % 順序数に対する下界の全体
\newcommand{\setmultiindexpos}[1]{\semigrposint^{#1}} % 正の整数からなる多重指数の空間
\newcommand{\setmultiindex}[1]{\monnat^{#1}} % 多重指数の空間
\newcommand{\setnpt}[1]{\lrbk{#1}} % 1~nまでのn点集合, または「i..j」で「整数iからjまでの集合」を表す. 添字集合としての利用を想定
\newcommand{\setodd}{\mathrm{Odd}} % 奇数全体
\newcommand{\setone}[1]{\cbk{#1}}
\newcommand{\setorderedpair}[1]{\dbk{#1}} % 順序対
\newcommand{\setpartfin}[1]{\fun{\mathrm{P}_{\txtfin}}{#1}} % 集合#1の有限分割全体
\newcommand{\setperiodiclattice}{\operatorname{Per}} % 周期格子, expedition0009283
\newcommand{\setpowfin}[1]{\fun{\mathrm{Pow}_{\txtfin}}{#1}} % 集合#1の有限部分集合全体
\newcommand{\setpow}[1]{\fun{\mathrm{Pow}}{#1}} % 集合#1のベキ集合
\newcommand{\setprime}{\operatorname{Prime}} % 素数全体がなす集合
\newcommand{\setprojop}{\operatorname{Proj}} % 射影作用素の集合
\newcommand{\setquot}[2]{\bigmiddleslash{#1}{#2}} % 商集合
\newcommand{\set}[2]{\left\{#1 \, \middle| \, #2\right\}}
\newcommand{\shconti}{\sheaf{\conti}} % 連続関数がなす層
\newcommand{\shcrosssection}{\Gamma} % 層の切断
\newcommand{\shdirectimage}[1]{#1_{\ast}} % sheaf 層の順像
\newcommand{\shdiscontinuoussection}[1]{\sqbk{#1}} % sheaf 不連続切断の層
\newcommand{\shdstschwartz}{\dual{\shdsttest}} % sheaf シュワルツ超関数がなす層
\newcommand{\shdsttest}{\sheaf{D}} % sheaf シュワルツの試験関数がなす層
\newcommand{\shfnhol}{\sheaf{O}} % 正則関数がなす層
\newcommand{\shfnmero}{\sheaf{M}} % 恒等的に非零の有理型関数がなす層
\newcommand{\shfnsato}{\sheaf{B}} % 佐藤超関数の層
\newcommand{\shgerm}[2]{\underline{#1}_{#2}} % sheaf 層の理論での芽
\newcommand{\shhigherdirectimage}[2]{R^{#1} \shdirectimage{#2}} % 層の高次順像または高次順像関手
\newcommand{\shhom}{\sheaf{Hom}} % 層の内部ホム関手, expedition0009065
\newcommand{\shinverseimage}[1]{#1^{\ast}} % sheaf 層の逆像
\newcommand{\shproperdirectimage}[1]{#1_{!}} % sheaf 層の固有順像
\newcommand{\shrealanalytic}{\sheaf{A}} % 実解析関数がなす層
\newcommand{\shsheafify}[1]{#1^{+}} % sheaf 前層の層化, expedition0008808
\newcommand{\shsingcochain}{\sheaf{S}} % 特異余鎖がなす層
\newcommand{\shskyscraper}[2]{#1_{#2}} % 摩天楼層
\newcommand{\shsmooth}{\sheaf{\smooth[E]}} % sheaf 滑らかな関数がなす層, 上つき添字を簡潔にするためC^{\infty}は使わない
\newcommand{\shstalk}[2]{\underline{#1}_{#2}} % 層での茎
\newcommand{\shverdierdual}[1]{#1^!} % sheaf ヴェルディエ双対
\newcommand{\shzeroext}[1]{#1_!} % sheaf 層の零延長, expedition0008825
\newcommand{\skewfldquatimag}{\mathbb{I}} % 四元数がなす斜体の虚部
\newcommand{\skewfldquatpureq}{\fldreal \boldsymbol{i} + \fldreal \boldsymbol{j} + \fldreal \boldsymbol{k}} % 四元数の意味での純虚数
\newcommand{\skewfldquat}{\fld{H}} % 四元数がなす斜体
\newcommand{\smfreeenergy}{F} % (ヘルムホルツの)自由エネルギー
\newcommand{\smgrandpartitionfunc}{\Xi} % 大分配関数
\newcommand{\smmeanmagnetization}{m} % 全磁化の一サイト平均, expedition0011058
\newcommand{\smmindist}[1]{\fun{\mathfrak{D}}{#1}} % 物質の安定性の議論での距離関数, expedition0006651
\newcommand{\smpartitionfunc}{Z} % 分配関数
\newcommand{\smreducertok}[2]{#1^{(#2)}} % $k$粒子系への縮約密度作用素, expedition0010013
\newcommand{\smspinsum}[1]{\mathring{#1}} % スピン和を取る作用素, expedition0006542
\newcommand{\smtotalmagnetization}{M} % 全磁化, expedition0011057
\newcommand{\sobolev}[1]{#1} % ソボレフ空間, WとH両方
\newcommand{\spaffine}{\mathbb{A}} % アフィン空間
\newcommand{\spaffspan}{\operatorname{aff}} % アフィン包, expedition0006580
\newcommand{\spclosedconv}{\mathop{\clos{\spconv}}} % 閉凸包, expedition0002189
\newcommand{\spclosedlinspan}{\mathop{\clos{\splinspan}}} % 閉線型包
\newcommand{\spconv}{\operatorname{conv}} % 凸包, 閉凸包は\closedconv, expedition0002189
\newcommand{\speigv}{\operatorname{Eig}} % 固有値の空間, モース理論, expedition0010235
\newcommand{\speuclid}{\mathsf{Euc}} % ユークリッド空間
\newcommand{\speucsub}[1]{\mathbb{#1}} % ユークリッド空間の部分集合・部分空間
\newcommand{\spfock}{\mathcal{F}} % フォック空間
\newcommand{\spkleinbottle}[1]{\mathbb{K}^{#1}} % クラインの壷
\newcommand{\splinspan}{\operatorname{span}} % 線型包
\newcommand{\spmatsing}{S} % 特異行列(非可逆行列)の全体
\newcommand{\spmat}[2]{\fun{M_{#1}}{{#2}}} % 全行列環
\newcommand{\spminkowski}{\mathbb{M}} % ミンコフスキー空間
\newcommand{\sppt}{\mathrm{pt}} % 一点・一点集合・一点空間
\newcommand{\stataic}{\mathrm{AIC}} % 赤池情報量基準
\newcommand{\statpredictivedistribution}[1]{#1^{\ast}} % 予測分布
\newcommand{\statwaic}{\mathrm{WAIC}} % 統計学, 広く使える情報量規準
\newcommand{\sumcomb}{\sum_{\mathrm{comb}}} % 組み合わせに関わる和
\newcommand{\sumfourprimed}[4]{\mathop{\mathop{\mathop{\sumprimed_{#1}}_{#2}}_{#3}}_{#4}} % 和, 共通
\newcommand{\sumfour}[4]{\mathop{\mathop{\mathop{\sum_{#1}}_{#2}}_{#3}}_{#4}} % 四重和
\newcommand{\sumprimed}{\mathop{\sideset{}{^{\prime}}\sum}} % 右上にプライム記号をつけた\sum
\newcommand{\sumthreeprimed}[3]{\mathop{\mathop{\sumprimed_{#1}}_{#2}}_{#3}} % 和, 共通
\newcommand{\sumthree}[3]{\mathop{\mathop{\sum_{#1}}_{#2}}_{#3}} % 三重和
\newcommand{\sumtwoprimed}[2]{\mathop{\sumprimed_{#1}}_{#2}} % 和, 共通
\newcommand{\sumtwo}[2]{\mathop{\sum_{#1}}_{#2}} % 二重和
\newcommand{\supalg}[2]{\mathbb{#1}_{#2}} % 超代数
\newcommand{\supthree}[3]{\mathop{\mathop{\sup_{#1}}_{#2}}_{#3}} % 上限
\newcommand{\suptwo}[2]{\mathop{\sup_{#1}}_{#2}} % 上限
\newcommand{\tdcompthermenthalpy}[1]{\sqfun{H}{#1}} % 熱力学, エンタルピー
\newcommand{\tdcompthermentropy}[1]{\sqfun{S}{#1}} % 熱力学, エントロピー
\newcommand{\tdcompthermgibbs}[1]{\sqfun{G}{#1}} % 熱力学, ギブスの自由エネルギー
\newcommand{\tdcompthermhelmholtz}[1]{\sqfun{F}{#1}} % 熱力学, ヘルムホルツの自由エネルギー
\newcommand{\tdcompthermintenergy}[1]{\sqfun{U}{#1}} % 熱力学, 内部エネルギー
\newcommand{\tdqmax}{Q_{\txtmax}} % 最大吸熱量
\newcommand{\tdworkadiabatic}{W_{\txttdadiabatic}} % 断熱仕事
\newcommand{\tdworkcyclic}{W_{\txttdcyclic}} % サイクルがなす仕事
\newcommand{\tdworkmax}{W_{\txtmax}} % 最大仕事
\newcommand{\topaffinesimplexbarycenter}[1]{\underline{#1}} % アフィン特異単体の重心, expedition0008931
\newcommand{\topaffinesingularsimplex}[1]{\sqbk{#1}} % アフィン単体, expedition0009195
\newcommand{\topbarysubd}{\operatorname{Sd}} % 重心細分, expedition0008943
\newcommand{\topcarr}{\operatorname{carr}} % 多面体の点に対する台, expedition0009263
\newcommand{\topcellaffsimp}[1]{\dbk{#1}} % 単体複体の単体をCW複体の単体とみなしたときの記号, expedition0009199
\newcommand{\topconeop}{\operatorname{\Upsilon}} % 錐作用素, expedition0009100
\newcommand{\topdist}{\operatorname{dist}} % 距離関数, 他と記号が重複するときに利用
\newcommand{\topgrdeck}{\operatorname{Deck}} % デッキ変換群, expedition0003436
\newcommand{\tophomotopygr}[2]{\fun{\pi_{#1}}{#2}} % ホモトピー群
\newcommand{\toplink}{\operatorname{Lk}} % 絡み数
\newcommand{\topmesh}{\operatorname{mesh}} % 単体のメッシュ, expedition0009968
\newcommand{\topopenstar}[2]{\operatorname{St}_{#1} #2} % 開星型集合, expedition0009269
\newcommand{\toppolyhedron}[1]{\abs{#1}} % 単体複体の多面体
\newcommand{\topproj}[2]{#1 P^{#2}} % 位相幾何での射影空間, #1 P^{#2}
\newcommand{\toprmap}{\operatorname{R}} % 重心細分に現れる, expedition0008943
\newcommand{\topsetdirectedopencover}{\operatorname{ROCover}} % 開被覆の全体, expedition0004832
\newcommand{\topsetopencoverclass}{\operatorname{\eqc{OCover}}} % 開被覆の全体の同値類, expedition0008899
\newcommand{\topsetopencover}{\operatorname{OCover}} % 開被覆の全体, expedition0004832
\newcommand{\topsimplexdelta}[1]{\sqbk{#1}} % 特殊な特異鎖, expedition0008935
\newcommand{\topsingularchain}{S} % 特異鎖
\newcommand{\topsingularcochain}{S} % 特異余鎖
\newcommand{\topsingularconvexchain}{A} % ユークリッド空間上の特異凸鎖, expedition0008932
\newcommand{\topsingularsimplexset}{T} % 特異単体がなす加群ではなく単に特異単体の全体集合, expedition0008932
\newcommand{\topsmashproduct}{\wedge} % スマッシュ積, expedition0008842
\newcommand{\topstandardconvexsimplexset}{B} % ユークリッド空間上の標準凸単体の集合, expedition0008932
\newcommand{\topstandardsimplex}{\Delta} % 標準単体, expedition0009973
\newcommand{\topwind}{\operatorname{Wind}} % 平面での巻き付き数
\newcommand{\txtIII}{\mathrm{III}} % III型因子環などのIII
\newcommand{\txtII}{\mathrm{II}} % 第二変数などの2(II)
\newcommand{\txtI}{\mathrm{I}} % I型因子環などのI
\newcommand{\txtabsconti}{\mathrm{ac}} % 絶対連続
\newcommand{\txtalg}{\mathrm{alg}} % 代数的
\newcommand{\txtappoint}{\mathrm{ap}} % 近似点
\newcommand{\txtarakiwoods}{\mathrm{AW},\mathrm{b}} % 荒木-ウッズ
\newcommand{\txtarakiwyss}{\mathrm{AW},\mathrm{f}} % 荒木-ヴィス
\newcommand{\txtasymp}{\mathrm{asymp}} % 漸近
\newcommand{\txtasym}{\mathrm{as}} % 反対称
\newcommand{\txtatom}{\mathrm{atom}} % 原子
\newcommand{\txtaut}{\mathrm{Aut}} % 自己同型群に対する条件を表す
\newcommand{\txtbasis}{\mathrm{basis}} % 基底, 基底つき有限次元線型空間の圏などに利用
\newcommand{\txtbc}{\mathrm{BC}} % 境界条件, Boundary Condition
\newcommand{\txtbdd}{\mathrm{b}} % 有界性を表す
\newcommand{\txtbec}{\mathrm{BEC}} % BEC, ボース-アインシュタイン凝縮
\newcommand{\txtbogoliubov}{\mathrm{Bog}} % ボゴリウボフ
\newcommand{\txtboltzmann}{\mathrm{B}} % ボルツマン
\newcommand{\txtborel}{\mathrm{b}} % ボレル可測性に関わる対象を指す
\newcommand{\txtbott}{\mathrm{Bott}} % ボット
\newcommand{\txtboundedbelow}{\mathrm{bb}} % 下への有界性を表す
\newcommand{\txtbox}{\mathrm{b}} % 箱位相などの箱
\newcommand{\txtbrownbridge}{\mathrm{BB}} % ブラウン橋
\newcommand{\txtbsn}{\mathrm{b}} % ボソンのb
\newcommand{\txtcameronmartin}{\mathrm{CM}} % キャメロン-マルティン (Cameron-Martin)
\newcommand{\txtcanonical}{\mathrm{C}} % 統計力学の正準状態などの canonical
\newcommand{\txtcar}{\mathrm{CAR}} % CAR環, 正準反交換関係の環
\newcommand{\txtcir}{\mathrm{cir}} % 同心円方向
\newcommand{\txtclassical}{\mathrm{cl}} % 古典量を表すcl
\newcommand{\txtclifford}{\mathrm{cl}} % クリフォード
\newcommand{\txtclose}{\mathrm{c}} % 閉
\newcommand{\txtcnt}{\mathrm{c}} % カウント, 数え上げ
\newcommand{\txtconf}{\mathrm{conf}} % 配位, configuration
\newcommand{\txtconst}{\mathrm{const}} % 定数性を表す添字
\newcommand{\txtconti}{\mathrm{c}} % 連続性を表す添字
\newcommand{\txtcoulomb}{\mathrm{C}} % クーロンを表すC
\newcommand{\txtcpt}{\mathrm{c}} % コンパクト性を表す添字
\newcommand{\txtcritical}{\mathrm{c}} % 臨界
\newcommand{\txtcyc}{\mathrm{cyc}} % 巡回
\newcommand{\txtdel}{\mathrm{del}} % 除外近傍などに利用
\newcommand{\txtderham}{\mathrm{dR}} % ディラック関係に利用する添字
\newcommand{\txtdiag}{\mathrm{diag}} % 対角
\newcommand{\txtdirac}{\mathrm{D}} % ディラック関係に利用する添字
\newcommand{\txtdiscrete}{\mathrm{d}} % 離散性を表すdiscrete
\newcommand{\txtdivergent}{\mathrm{div}} % 発散
\newcommand{\txtdual}{\mathrm{dual}} % 双対
\newcommand{\txteff}{\mathrm{eff}} % 実効的
\newcommand{\txtelhole}{\txtel \hyphen \txthole} % 電子-正孔系
\newcommand{\txtel}{\mathrm{el}} % 電子のel
\newcommand{\txtemlorentz}{\mathrm{Lorentz}} % ローレンツ変換のローレンツ
\newcommand{\txtemmag}{\mathrm{mag}} % 電磁気の磁気
\newcommand{\txtemphoton}{\mathrm{ph}} % 光子
\newcommand{\txtem}{\mathrm{em}} % 電磁場のem
\newcommand{\txtentireextension}{\mathrm{ent}} % 整関数拡張を持つ対象, expedition0009682
\newcommand{\txtergode}{\mathrm{erg}} % エルゴード, エルゴード性
\newcommand{\txtess}{\mathrm{ess}} % 本質的スペクトルのess
\newcommand{\txteuclid}{\mathrm{E}} % ユークリッド, ユークリッド化の添字
\newcommand{\txtevenodd}{\txteven / \txtodd} % 偶奇
\newcommand{\txteven}{\mathrm{e}} % 偶
\newcommand{\txtexchange}{\mathrm{ex}} % 交換, 交換相互作用などに利用
\newcommand{\txtexcludezero}{\times} % 環や体から加法の零元を除く
\newcommand{\txtexternal}{\mathrm{ext}} % 外部
\newcommand{\txtfield}{\mathrm{fld}} % 場(field), 場のハミルトニアンなどに利用
\newcommand{\txtfin}{\mathrm{fin}} % 有限のfin
\newcommand{\txtflip}{\mathrm{flip}} % フリップ
\newcommand{\txtfkn}{\mathrm{FKN}} % ファインマン-カッツ-ネルソン, FKN
\newcommand{\txtfock}{\mathrm{F}} % フォック
\newcommand{\txtfrechet}{\mathrm{f}} % フレッシェフィルターなどに利用, Fréchet
\newcommand{\txtfrmn}{\mathrm{f}} % フェルミオンのf
\newcommand{\txtfr}{\mathrm{fr}} % 自由場などを表すfr
\newcommand{\txtfull}{\mathrm{full}} % 全空間などのfull：totを利用する場合もある
\newcommand{\txtgeom}{\mathrm{geom}} % 幾何的
\newcommand{\txtgrandcanonical}{\mathrm{GC}} % 統計力学の大正準状態など grandcanonical
\newcommand{\txtgreaterone}{>1} % 1より大きい切断, expedition0010213
\newcommand{\txtgs}{\mathrm{gs}} % 基底状態
\newcommand{\txthilbertschmidt}{\mathrm{HS}} % ヒルベルト-シュミット作用素などに利用
\newcommand{\txthole}{\mathrm{hole}} % 正孔(hole)
\newcommand{\txtholomorphic}{\mathrm{hol}} % 関数論の正則性
\newcommand{\txthorizontal}{\mathrm{h}} % 水平成分
\newcommand{\txthydrogen}{\mathrm{hyd}} % 水素原子のハミルトニアンに利用
\newcommand{\txtimag}{\mathrm{im}} % 虚
\newcommand{\txtinflaredred}{\mathrm{IR}} % 赤外
\newcommand{\txtinfomld}{\mathrm{ML}} % 最尤復号
\newcommand{\txtinteraction}{\mathrm{I}} % 相互作用
\newcommand{\txtircapitalized}{\mathrm{IR}} % 赤外または infrared
\newcommand{\txtirsingular}{\mathrm{irs}} % 赤外特異
\newcommand{\txtkms}{\mathrm{KMS}} % KMS状態の集合を表す
\newcommand{\txtleb}{\mathrm{Leb}} % ルベーグ
\newcommand{\txtleft}{\mathrm{l}} % 左
\newcommand{\txtleqone}{\leq 1} % 1以下に対する切断, expedition0010213
\newcommand{\txtlipschitz}{\mathrm{Lip}} % リプシッツ連続などのLip
\newcommand{\txtloc}{\mathrm{loc}} % 局所
\newcommand{\txtloop}{\mathrm{loop}} % ループ空間などのループに利用
\newcommand{\txtlow}{\mathrm{low}} % 低
\newcommand{\txtmax}{\mathrm{max}} % 最大
\newcommand{\txtmeanfield}{\mathrm{mf}} % 平均場
\newcommand{\txtmeromorphic}{\mathrm{mero}} % 関数論の有理型性
\newcommand{\txtmin}{\mathrm{min}} % 最小
\newcommand{\txtmix}{\mathrm{mix}} % 混合
\newcommand{\txtnconv}{\mathrm{nconv}} % non conventional
\newcommand{\txtneg}{\mathrm{-}} % 負の部分
\newcommand{\txtnelson}{\mathrm{N}} % ネルソンの頭文字
\newcommand{\txtnonent}{\mathrm{nonent}} % 非エンタングルメント
\newcommand{\txtnonneg}{\mathrm{+}} % 正または非負の部分
\newcommand{\txtnonrel}{\mathrm{NR}} % 非相対論的
\newcommand{\txtnonzero}{\mathrm{nz}} % 非零, 非零モードを表す
\newcommand{\txtoaaw}{\mathrm{AW}} % 荒木-ウッズ・荒木-ヴィスの添字
\newcommand{\txtodd}{\mathrm{o}} % 奇
\newcommand{\txtopen}{\mathrm{o}} % 開
\newcommand{\txtoperator}{\mathrm{op}} % 反対環・反対圏
\newcommand{\txtopposite}{\mathrm{op}} % 反対環・反対圏
\newcommand{\txtosterwalderschrader}{\mathrm{OS}} % オスターヴァルダー-シュラーダー
\newcommand{\txtpartial}{\mathrm{p}} % 部分トレースなどの「部分」
\newcommand{\txtparticle}{\mathrm{p}} % 粒子のp
\newcommand{\txtparticular}{\mathrm{p}} % 特解(particular solution)のp
\newcommand{\txtpathwiseconnected}{\mathrm{pc}} % 弧状連結
\newcommand{\txtpaulifierz}{\mathrm{PF}} % パウリ-フィールツ
\newcommand{\txtpauli}{\mathrm{Pauli}} % パウリ
\newcommand{\txtperiodicbdc}{\mathrm{P}} % 周期境界条件
\newcommand{\txtperiod}{\mathrm{per}} % 周期, 周期境界条件
\newcommand{\txtphysgrav}{\mathrm{grav}} % 重力
\newcommand{\txtphysho}{\mathrm{ho}} % 調和振動を表すho
\newcommand{\txtphysnho}{\mathrm{nho}} % 非調和振動を表すnho
\newcommand{\txtphysrel}{\mathrm{rel}} % 相対論的
\newcommand{\txtphysultrarel}{\mathrm{urel}} % 超相対論的
\newcommand{\txtphys}{\mathrm{phys}} % 物理
\newcommand{\txtpoint}{\mathrm{pt}} % 点, 一点集合は\sppt
\newcommand{\txtpoisson}{\mathrm{poisson}} % ポアソン
\newcommand{\txtpolynomial}{\mathrm{poly}} % 多項式係数の微分作用素環などに利用
\newcommand{\txtprobas}{\mathrm{a.s.}} % ほとんど確実に
\newcommand{\txtprobev}{\mathrm{ev}} % eventually(確率論)
\newcommand{\txtprobio}{\mathrm{io}} % infinitely often(確率論)
\newcommand{\txtprob}{\mathrm{Prob}} % 確率
\newcommand{\txtps}{\mathrm{ps}} % piecewise smooth, 区分的に滑らか
\newcommand{\txtpwsmooth}{\mathrm{ps}} % 区分的に滑らか
\newcommand{\txtqmepr}{\mathrm{EPR}} % 量子力学でのEPR
\newcommand{\txtqmev}{\mathop{\mathrm{eV}}} % 電子ボルト
\newcommand{\txtqmtfd}{\mathrm{TFD}} % トーマスフェルミディラック, Thomas-Fermi-Dirac, expedition0006730
\newcommand{\txtqmtf}{\mathrm{TF}} % トーマスフェルミ, Thomas-Fermi, expedition0006730
\newcommand{\txtquadvanhove}{\mathrm{qvH}} % 二次ファン・ホーフェ
\newcommand{\txtradiation}{\mathrm{rad}} % 輻射
\newcommand{\txtrad}{\mathrm{rad}} % 動径方向
\newcommand{\txtreal}{\mathrm{real}} % 実
\newcommand{\txtregular}{\mathrm{reg}} % 正則な測度などの正則
\newcommand{\txtrelative}{\mathrm{rel}} % 相対
\newcommand{\txtrel}{\mathrm{R}} % 相対論的
\newcommand{\txtreminder}{\mathrm{R}} % 剰余項
\newcommand{\txtrenormalization}{\mathrm{ren}} % くり込み
\newcommand{\txtresidual}{\mathrm{r}} % 剰余
\newcommand{\txtright}{\mathrm{r}} % 右
\newcommand{\txtsa}{\mathrm{sa}} % 自己共役
\newcommand{\txtschrodinger}{\mathrm{S}} % シュレディンガー表現などに利用
\newcommand{\txtsegal}{\mathrm{S}} % シーガルの場の作用素に利用
\newcommand{\txtselfconsistent}{\mathrm{sc}} % 自己無撞着
\newcommand{\txtself}{\mathrm{self}} % 自己エネルギーなどの自己(self)
\newcommand{\txtsingular}{\mathrm{s}} % 特異
\newcommand{\txtspinboson}{\mathrm{SB}} % スピン-ボソン模型
\newcommand{\txtspin}{\mathrm{spin}} % スピン
\newcommand{\txtspph}{\mathrm{ph}} % 相空間
\newcommand{\txtstable}{\mathrm{s}} % 安定, 安定多様体などに利用
\newcommand{\txtstandard}{\mathrm{st}} % 超準解析での標準
\newcommand{\txtstark}{\mathrm{stark}} % シュタルクハミルトニアンに利用
\newcommand{\txtstateap}{\mathrm{EAP}} % 事後推定値, expected a posterior
\newcommand{\txtstrongresolvent}{\mathrm{s} \hyphen \text{res}} % 強レゾルベント収束
\newcommand{\txtstrong}{\mathrm{s}} % 強位相などの強
\newcommand{\txtsuper}{\mathrm{s}} % 超代数の超
\newcommand{\txtsym}{\mathrm{s}} % 対称, テンソル積に使う
\newcommand{\txttdadiabatic}{\mathrm{ad}} % 断熱
\newcommand{\txttdaq}{\mathrm{aq}} % 断熱準静, adiabatic quasistatic
\newcommand{\txttda}{\mathrm{a}} % 断熱操作, adiabatic operation
\newcommand{\txttdcyclic}{\mathrm{cyc}} % サイクルがなす仕事に利用
\newcommand{\txttdiq}{\mathrm{iq}} % 等温準静, isothermal quasistatic
\newcommand{\txttdiw}{\mathrm{i}_0} % 広義等温操作, isothermal operation in the wider sense
\newcommand{\txttdi}{\mathrm{i}} % 等温操作, isothermal operation
\newcommand{\txttot}{\mathrm{tot}} % 全電荷などの「全」
\newcommand{\txttruncatedfn}{\mathrm{T}} % 連結相関関数を表す記号
\newcommand{\txtultraviolet}{\mathrm{UV}} % 紫外
\newcommand{\txtuninorm}{\mathrm{unif}} % 一様
\newcommand{\txtunstable}{\mathrm{u}} % 不安定, 不安定多様体などに利用
\newcommand{\txtvanderwaals}{\mathrm{vdW}} % ファンデルワールス
\newcommand{\txtvanhowe}{\mathrm{vH}} % ファン・ホーフェ
\newcommand{\txtvanish}{\mathrm{van}} % 消滅, 原点0で消える関数の関数空間などに利用
\newcommand{\txtvertical}{\mathrm{v}} % 垂直成分
\newcommand{\txtweak}{\mathrm{w}} % 弱位相などの弱
\newcommand{\txtwick}{\mathrm{W}} % Wick積による部分空間, expedition0009962
\newcommand{\txtwstar}{w^{\ast}} % 汎弱位相用の記号
\newcommand{\txtzeeman}{\mathrm{zeeman}} % ゼーマンハミルトニアンに利用
\newcommand{\vadiv}{\operatorname{div}} % ベクトル場の発散, リーマン多様体上の議論にも援用, expedition0007443, expedition0002435
\newcommand{\vagrad}{\operatorname{grad}} % 勾配作用素
\newcommand{\vainnprod}{\cdot} % ベクトル解析枠, 高校以来の内積の記号
\newcommand{\vanishsobolev}{D} % LiebLoss1や物質の安定性に現れるソボレフ空間
\newcommand{\varot}{\operatorname{rot}} % 回転, expedition0002223
\newcommand{\vavecprod}{\times} % ベクトル解析枠, 三次元での外積・ベクトル積
\newcommand{\vecone}{\boldsymbol{1}} % 太文字の1, 恒等写像ではなく適当な意味でのベクトルに対して利用
\newcommand{\vhlnt}[2]{\fun{\boldsymbol{n}_{#1}^{-}}{#2}} % ファン・ホーフェ極限に現れる, 互いに素なC_aの平行移動でIに完全に含まれる最大個数
\newcommand{\vhlsets}[2]{\fun{\Gamma_{#1}^{-}}{#2}} % Bratteli-Robinson 2 P287の和集合
\newcommand{\vhslnt}[2]{\fun{\boldsymbol{n}_{#1}^{\pm}}{#2}} % vhsntとvhlntを両方まとめて書きたい場合に使う
\newcommand{\vhsnt}[2]{\fun{\boldsymbol{n}_{#1}^{+}}{#2}} % ファン・ホーフェ極限に現れる, C_aの平行移動で領域Iを覆うために必要なC_aの最小個数
\newcommand{\vhssets}[2]{\fun{\Gamma_{#1}^{+}}{#2}} % Bratteli-Robinson 2 P287の和集合
\newcommand{\vhsurf}[2]{\fun{\mathrm{surf}_{#1}}{#2}} % 厚さrの表面集合


%%%%%%%%%%%%%%%%%%%%%%%%%%%%%%%
% Additional header includes
%%%%%%%%%%%%%%%%%%%%%%%%%%%%%%%

%%%%%%%%%%%%%%%%%%%%%%%%%%%%%%%
% Title information
%%%%%%%%%%%%%%%%%%%%%%%%%%%%%%%
\title{Interacting Lattice Bosons on the Concrete Buchholz Algebra\\\vspace{0.5em}{\large All-Temperature Equilibrium States from Functional-Integral Local-Number Bounds}}

\author{%
Yoshitsugu Sekine\\{\small\texttt{4429sekine@gmail.com}}%
}

\date{\today}

%%%%%%%%%%%%%%%%%%%%%%%%%%%%%%%
% Document body
%%%%%%%%%%%%%%%%%%%%%%%%%%%%%%%
\begin{document}

\maketitle

\begin{abstract}
Interacting Bose--Hubbard dynamics is formulated on the concrete Buchholz algebra in Fock space. Under repulsive finite-range interactions and a low-activity gap, a discrete functional integral gives volume-uniform exponential local-number bounds at all temperatures. These bounds remove the local-moment hypothesis and yield finite-volume approximation and KMS accumulation states in \cite{DeuchertLampartLemm001}.

\noindent\textbf{Keywords:} Bose--Hubbard model, bosonic Fock space, Buchholz algebra, thermodynamic limit, KMS state, local particle number, Feynman--Kac representation
\end{abstract}


\setcounter{tocdepth}{3}
\tableofcontents

\section{Introduction}\label{introduction}

The scarcity of physically relevant automorphism groups on the Weyl algebra has long limited its use for interacting bosonic systems. The natural automorphisms of the Weyl algebra are induced by symplectic transformations and describe quadratic Hamiltonians, whereas even simple nonquadratic Hamiltonians may fail to preserve the algebra. The resolvent algebra was introduced to overcome this restriction and to incorporate resolvents of Hamiltonians and other unbounded observables into a representation-independent \(\oacstar\)-algebraic framework \cite[Sections 1 and 5]{BuchholzGrundling002}. Its rich ideal structure is not an incidental complication but a source of physical information \cite{DetlevBuchholz001}. This structure has been used in general criteria for Bose--Einstein condensation and in the analysis of trapped and interacting Bose systems \cite{BahnsBuchholz001,DetlevBuchholz002,DetlevBuchholz004,DetlevBuchholz005,DetlevBuchholz012}. These applications also show why the resolvent algebra must be evaluated through concrete models, of which only a limited range has been treated to date. One principal aim of this paper is to extend that range to the equilibrium theory of the Bose--Hubbard model.

Like the Weyl algebra, the resolvent algebra is defined before a representation is selected. This representation independence is essential to its role as an algebra of observables. The contrast with a von Neumann algebra makes the point precise. After an automorphism group and a ground, equilibrium, or vacuum state have been chosen, the corresponding von Neumann algebra is obtained as the weak closure of the \(\oacstar\)-algebra in the GNS representation of that state. It retains the folium and the thermodynamic sector of the reference state. A nonfactorial representation may have a nontrivial center, and its central decomposition records classical order parameters or the central-spectrum variables that distinguish sectors. This is the micro--macro interpretation of the center in sector theory \cite{IzumiOjima002}. In symmetry-breaking situations, the center of the mixed representation records the classical distribution of phases, while its factor components describe the individual quantum phases. The abstract resolvent algebra does not build such a classical component into the kinematics. Its center is trivial, even though its proper ideals remain physically informative \cite[Theorem 4.10]{BuchholzGrundling002}. The restriction to one reference folium is also visible in thermal representations. For example, factorial KMS representations of type III at distinct inverse temperatures are disjoint \cite[Theorem 5.3.35]{BratteliRobinson004}. The normal states of a von Neumann algebra selected by one state remain in that folium. Such an algebra cannot serve as a representation-independent container for all thermodynamic sectors or for unrelated dynamics. An abstract \(\oacstar\)-algebra is instead intended to retain the kinematics before any such choice. Developing as much of the dynamics and equilibrium theory as possible at this abstract level is a basic requirement rather than a matter of formal generality.

The first operator-algebraic problem for a concrete Hamiltonian is the construction of its automorphism group. This is the algebraic counterpart of self-adjointness in the operator formulation, but it involves two additional obstructions. The first is continuity. Even free bosonic dynamics need not be point-norm continuous in the natural \(\oacstar\)-norm \cite[Section 7.1]{BuchholzGrundling002}. A continuous automorphism group may emerge only after a suitable representation has been selected and the weak closure has been taken, provided that the dynamics is continuously implementable there. Such a passage to a von Neumann algebra resolves continuity within the chosen folium but reintroduces the representation dependence that the abstract resolvent algebra was designed to avoid.

The second obstruction is invariance of the algebra. Free bosonic dynamics acts by automorphisms of the resolvent algebra, although the action need not be point-norm continuous, whereas interacting dynamics can send elements of the gauge-invariant resolvent algebra into a larger algebra. An interacting evolution that leaves the resolvent algebra requires a larger ambient \(\oacstar\)-algebra on which it acts by automorphisms. For continuous Bose systems, particle-sector consistency in the Fock representation produces such an extension of the gauge-invariant resolvent algebra and accommodates broad classes of interacting dynamics \cite{DetlevBuchholz002,DetlevBuchholz003}. The same structure supports the analysis of proper condensates, for which infinite particle accumulation is detected by particle-number resolvents \cite{DetlevBuchholz004,DetlevBuchholz005}. The corresponding lattice extension is called the Buchholz algebra in \cite[Definition 2.2]{DeuchertLampartLemm001}. Within the Fock representation, it contains the gauge-invariant resolvent algebra and is preserved by the Bose--Hubbard dynamics \cite[Proposition 3.1]{DeuchertLampartLemm001}. In this sense, the Buchholz algebra is an ambient algebra for the gauge-invariant resolvent algebra. The continuous-space construction also contains nonzero invariant subalgebras on which the automorphism group is point-norm continuous, obtained by taking the continuous elements of the action \cite[Theorem 4.6]{DetlevBuchholz002}. Both the continuity and invariance problems have partial solutions once the Fock representation is fixed.

The Bose--Hubbard model is a stringent test of this framework because every lattice site has an infinite-dimensional local Hilbert space. An interaction-picture cluster expansion gives volume-uniform local particle-number estimates for its finite-volume Gibbs states at sufficiently high temperature and for arbitrary fixed chemical potential \cite[Theorem 1]{GongKuwaharaTong001}. Combined with the Deuchert--Lampart--Lemm automorphism and KMS construction, these estimates yield equilibrium states in a high-temperature region without a separate low-activity restriction on the chemical potential. The infinite-volume Bose--Hubbard dynamics on the concrete Buchholz algebra, its finite-volume approximation, and its KMS accumulation states are established in \cite[Proposition 3.1 and Theorems 4.1 and 5.2]{DeuchertLampartLemm001}. The finite-volume approximation and equilibrium-state conclusions, however, assume uniform polynomial moments of every required order for the local particle number. These are the hypotheses in \cite[equations (37) and (90)]{DeuchertLampartLemm001}. This state condition is technically severe. It can fail in condensed regimes with divergent local occupation, including proper-condensate limits. It also leaves the existence theorem conditional on an estimate that is not derived from the Hamiltonian. Superstability and a functional-integral representation are proposed as a route to the missing estimate in \cite[Remark 5.3]{DeuchertLampartLemm001}.

The present paper implements that proposal. For repulsive finite-range density interactions, we assume that the chemical potential makes the grand-canonical one-particle Hamiltonian bounded below by a strictly positive constant. A discrete functional integral yields a volume-uniform exponential bound on the local particle number at every positive inverse temperature. Theorem \ref{thm:fock-all-temperature-local-number} states the precise result. The exponential estimate is stronger than the polynomial condition used in the earlier approximation argument. It permits the occupation cutoff to be removed directly and gives the all-temperature finite-volume approximation in Theorem \ref{thm:fock-all-temperature-local-approximation} without imposing a local-number hypothesis on the state. The same argument produces weak-\(\ast\) accumulation states that are invariant on the gauge-invariant resolvent subalgebra and satisfy the KMS boundary relation there for the infinite-volume automorphism group, as stated in Corollary \ref{cor:fock-all-temperature-kms-states}. The high-temperature result and the present theorem cover complementary regimes. The high-temperature result allows arbitrary fixed chemical potential at sufficiently high temperature, whereas the present theorem allows every positive temperature under a low-activity gap condition. The gap assumption does not assert the absence of Bose--Einstein condensation in the Bose--Hubbard model. It excludes the threshold regime in which the chemical potential closes the one-particle gap, and that regime lies outside the present theorem.

The present construction nevertheless remains tied to the bosonic Fock representation. From the viewpoint of algebraic quantum statistical mechanics, the Buchholz algebra should ultimately be defined without a distinguished reference representation and should serve as an abstract ambient algebra for the gauge-invariant resolvent algebra. The Bose--Hubbard dynamics should then be defined as an automorphism group of this abstract algebra, while the ordinary \(\oacstar\)-KMS condition should be imposed on an appropriate invariant subalgebra on which the action is point-norm continuous. This would preserve the ability of the abstract algebra to accommodate distinct folia, temperatures, and phases before a GNS representation is selected.

The functional integral used here should also admit an intrinsic formulation at that level. The correspondence between stochastically positive KMS systems, periodic stochastic processes, and positive semigroups was developed in a formally \(\oacstar\)-algebraic setting by Klein--Landau \cite[Sections 1 and 6--8]{KleinLandau001}. Its reconstruction nevertheless passes through a GNS Hilbert space and the associated von Neumann algebra. Standard-form perturbation theory likewise constructs perturbed Liouvilleans and KMS states after a representation has been selected \cite{DerezinskiJaksicPillet001}. For bosonic fields, the Araki--Woods representation can be identified with a Gaussian periodic path space, and this equivalence yields interacting dynamics and equilibrium states through functional integration \cite[Sections 17.1.5 and 21.4--21.5]{DerezinskiGerard001}. These results are sharp at the von Neumann-algebraic level, but a deeper relation between probability, stochastic processes, and abstract \(\oacstar\)-algebras is needed for the resolvent and Buchholz algebras. Establishing the present all-temperature result in the fixed Fock representation is a first step toward that representation-independent theory.

The paper first fixes the Fock-space Bose--Hubbard model and the concrete Buchholz algebra. It then isolates the missing local-number input in the existing dynamical and equilibrium arguments, constructs the discrete functional integral, proves the exponential bound, and derives the infinite-volume dynamics and KMS states in the low-activity regime.

\section{Fock-Space Bose--Hubbard Model}\label{fock-space-bosehubbard-model}

The concrete model fixes the spatial and particle-number structures before either dynamics or states are taken to an infinite-volume limit.

\subsection{Graph, Fock spaces, and local observables}\label{graph-fock-spaces-and-local-observables}

Let \(\Gamma\) be a connected countable graph of uniformly bounded degree. The graph distance is denoted by \(\fun{\topdist}{x,y}\). For a finite \(X \Subset \Gamma\) and \(R \geq 0\), the closed graph neighborhood is \[\begin{aligned}
X[R]
& =
\set{y \in \Gamma}{\fun{\topdist}{y,X} \leq R}.
\end{aligned}\] The interior boundary is \[\begin{aligned}
\partial X
& =
\set{x \in X}{\fun{\topdist}{x,\Gamma \setminus X} = 1}.
\end{aligned}\] Assume that \(\Gamma\) is \(d\)-dimensional with surface parameter \(\sigma > 0\) in the sense that \[\begin{aligned}
\sup_{x \in \Gamma}
\sup_{R \geq 1}
\frac{\abscard{\partial\rbk{x[R]}}}{R^{d - 1}}
& \leq
\sigma.
\end{aligned}\] This is Definition 2.1 of \cite{DeuchertLampartLemm001} and implies the polynomial ball-volume bound \begin{equation}
\begin{aligned}
\sup_{x\in\Gamma}\abscard{x[R]}
& \leq
\sigma R^d,
\quad
R\geq1.
\end{aligned}
\label{eq:fock-polynomial-ball-volume-growth}
\end{equation} The phrase polynomial ball-volume growth refers to the estimate \eqref{eq:fock-polynomial-ball-volume-growth}.

For \(X \Subset \Gamma\) or \(X = \Gamma\), the one-particle and bosonic Fock spaces are \[\begin{aligned}
\sphilbfrak{h}_{X}
& =
\fun{\lpseq^{2}}{X},
&
\sphilb{F}_{X}
& =
\bigoplus_{n \in \monnat}\sphilbfrak{h}_{X}^{\otimes_{\txtsym}^{n}}.
\end{aligned}\] For \(n \in \monnat\), the orthogonal projection \(P_{X,n}\) onto the \(n\)-particle summand is defined by \begin{equation}
\begin{aligned}
\Ran P_{X,n}
& =
\sphilbfrak{h}_{X}^{\otimes_{\txtsym}^{n}}.
\end{aligned}
\label{eq:fock-particle-sector-projection}
\end{equation} The total number operator is \(\opfocknumber_X = \sum_{n \in \monnat}n P_{X,n}\). For \(m\in\monnat\), define \begin{equation}
\begin{aligned}
P_{X,\leq m}
& =
\sum_{n=0}^{m}P_{X,n},
&
\sphilb{D}_{X,\txtfin}
& =
\bigcup_{m\in\monnat}P_{X,\leq m}\sphilb{F}_X.
\end{aligned}
\label{eq:fock-finite-particle-subspace}
\end{equation} Thus \(\sphilb{D}_{X,\txtfin}\) consists precisely of the vectors having nonzero components in only finitely many particle-number summands. For \(x \in X\), the creation, annihilation, and local occupation operators on \(\sphilb{D}_{X,\txtfin}\) are \(\faadj{a_x}\), \(a_x\), and \(N_x = \faadj{a_x}a_x\). For \(x,y\in X\), their canonical commutation relations and the resulting local-number commutators on \(\sphilb{D}_{X,\txtfin}\) are \begin{equation}
\begin{aligned}
\commutator{a_x}{\faadj{a_y}}
& =
\kroneckerdelta_{x,y},
&
\commutator{a_x}{a_y}
& =
0,
\\
\commutator{\faadj{a_x}}{a_y}
& =
-\kroneckerdelta_{x,y},
&
\commutator{\faadj{a_x}}{\faadj{a_y}}
& =
0,
\\
\commutator{N_x}{\faadj{a_y}}
& =
\faadj{a_x}\commutator{a_x}{\faadj{a_y}}
+
\commutator{\faadj{a_x}}{\faadj{a_y}}a_x
=
\kroneckerdelta_{x,y}\faadj{a_x}
=
\kroneckerdelta_{x,y}\faadj{a_y},
\\
\commutator{N_x}{a_y}
& =
\faadj{a_x}\commutator{a_x}{a_y}
+
\commutator{\faadj{a_x}}{a_y}a_x
=
-\kroneckerdelta_{x,y}a_x
=
-\kroneckerdelta_{x,y}a_y.
\end{aligned}
\label{eq:fock-creation-annihilation-number-commutators}
\end{equation} For finite \(X \Subset \Gamma\), let \(\Omega_X\) be the Fock vacuum. For each multi-index \(\boldsymbol{\nu} = \seq{\nu_x}{x \in X} \in \monnat^X\), define \begin{equation}
\begin{aligned}
\Psi_{X,\boldsymbol{\nu}}
& =
\prod_{x \in X}
\frac{\rbk{\faadj{a_x}}^{\nu_x}}{\sqrt{\nu_x!}}
\Omega_X,
\\
N_x\Psi_{X,\boldsymbol{\nu}}
& =
\nu_x\Psi_{X,\boldsymbol{\nu}},
&
x
& \in
X,
\\
\opfocknumber_X\Psi_{X,\boldsymbol{\nu}}
& =
\rbk{\sum_{x \in X}\nu_x}\Psi_{X,\boldsymbol{\nu}}.
\end{aligned}
\label{eq:fock-occupation-number-basis}
\end{equation} The creation operators at distinct sites commute, and hence the product does not depend on the ordering of \(X\). The vectors \(\Psi_{X,\boldsymbol{\nu}}\), \(\boldsymbol{\nu} \in \monnat^X\), form the orthonormal occupation-number basis of \(\sphilb{F}_X\).

Let \(\physham[h]_{X}\) be the bounded hopping operator on \(\sphilbfrak{h}_{X}\) with matrix elements \[\begin{aligned}
\rbk{\physham[h]_{X}}_{x,y}
& =
\begin{cases}
-1,
& x \sim y,
\\
0,
& x \not\sim y.
\end{cases}
\end{aligned}\] Assume that the uniform degree of the graph is strictly positive, and set \begin{equation}
\begin{aligned}
0<C_h
& =
\sup_{x \in \Gamma}\abscard{\set{y \in \Gamma}{y \sim x}}.
\end{aligned}
\label{eq:fock-uniform-degree-bound}
\end{equation} Then \(\norm{\physham[h]_{\Lambda}} \leq C_h\) for every finite \(\Lambda\), and \(-\physham[h]_{\Lambda} - C_h\) has nonnegative off-diagonal entries and nonpositive row sums. Fix \(U > 0\) and \(r\in\semigrposint\), and let \(w \colon \Gamma \times \Gamma \to \fldreal\) be a bounded symmetric function satisfying \begin{equation}
\begin{aligned}
\fun{w}{x,y}
& \geq
0,
&
\fun{\topdist}{x,y}>r
& \implies
\fun{w}{x,y}=0.
\end{aligned}
\label{eq:fock-interaction-range}
\end{equation} The density interaction is defined by \[\begin{aligned}
\fun{v}{x,y}
& =
U\fun{\fndef{\setone{x}}}{y} + \fun{w}{x,y}.
\end{aligned}\] Thus \(v\) is pointwise nonnegative, and its interaction range is at most the integer \(r\) fixed in \eqref{eq:fock-interaction-range}. For finite \(\Lambda \Subset \Gamma\), the number-conserving Hamiltonian and its grand-canonical shift are defined on \(\sphilb{D}_{\Lambda,\txtfin}\) from \eqref{eq:fock-finite-particle-subspace} by \begin{equation}
\begin{aligned}
\physham_{\txtfr,\Lambda}
& =
\sum_{x,y \in \Lambda}
\rbk{\physham[h]_{\Lambda}}_{x,y}\faadj{a_x} a_y,
\\
V_{\Lambda}
& =
\sum_{x,y \in \Lambda}
\fun{v}{x,y}\faadj{a_x}\faadj{a_y} a_x a_y,
\\
H_{\Lambda}
& =
\physham_{\txtfr,\Lambda} + V_{\Lambda},
\\
K_{\Lambda,\smchemicalpotential_{\Lambda}}
& =
H_{\Lambda} - \smchemicalpotential_{\Lambda}\opfocknumber_{\Lambda}.
\end{aligned}
\label{eq:fock-bose-hubbard-hamiltonian}
\end{equation} The ordered-pair convention in \eqref{eq:fock-bose-hubbard-hamiltonian} agrees with \cite[equation (10)]{DeuchertLampartLemm001} and introduces no additional factor \(1/2\).

The Hamiltonian commutes with \(\opfocknumber_{\Lambda}\). Its restriction \(H_{\Lambda,n}\) to \(P_{\Lambda,n}\sphilb{F}_{\Lambda}\) is bounded for every \(n \in \monnat\). For \(X = \Gamma\), the same formula on the \(n\)-particle space defines a bounded operator \(H_{\Gamma,n}\) because the hopping has bounded degree and the interaction is bounded and finite-range. The symbol \(H_{\Gamma}\) denotes the self-adjoint direct sum \(\bigoplus_{n\in\monnat}H_{\Gamma,n}\) on \(\sphilb{F}_{\Gamma}\).

\subsection{Concrete Buchholz algebra}\label{concrete-buchholz-algebra}

The gauge-invariant resolvent algebra on \(\sphilbfrak{h}_{X}\) is denoted by \(\fun{\oaresolventalgebra}{X}^{\gamma}\). The concrete Buchholz algebra of \cite[Definition 2.2]{DeuchertLampartLemm001} is \[\begin{aligned}
\oa{B}_{X}
& =
\set{A \in \opspbddlin{\sphilb{F}_{X}}}
{\text{for every $n \in \monnat$ there is $R \in
\fun{\oaresolventalgebra}{X}^{\gamma}$ with $A P_{X,\leq n} = R P_{X,\leq n}$}}.
\end{aligned}\] For \(X = \Gamma\), write \(\oa{B} = \oa{B}_{\Gamma}\). Proposition 2.4 of the cited paper identifies \(\oa{B}\) with the bounded inverse limit of its particle-sector coordinate algebras. This characterization explains why the algebraic dynamics can exist without a uniform continuity estimate in the particle-number index \(n\).

\section{Existing Dynamics on the Concrete Buchholz Algebra and the Missing Estimate}\label{existing-dynamics-on-the-concrete-buchholz-algebra-and-the-missing-estimate}

The full automorphism group on the concrete Buchholz algebra exists before any state or moment condition is selected. The moment condition enters only when this group is approximated by spatially finite Hamiltonians and implemented in an infinite-particle representation.

\subsection{Automorphisms on the full Buchholz algebra}\label{automorphisms-on-the-full-buchholz-algebra}

For \(X \Subset \Gamma\) or \(X = \Gamma\), let \(\alpha_X\) denote the one-parameter family whose time-\(t\) map is defined by \[\begin{aligned}
\fun{\alpha_{X,t}}{A}
& =
\napiernum^{\imunit tH_X}A\napiernum^{-\imunit tH_X},
\quad
A \in \oa{B}_X,
\quad
t \in \fldreal.
\end{aligned}\]

\begin{fact}[Deuchert--Lampart--Lemm, Proposition 3.1]\label{fact:dll-buchholz-automorphism}
If $\fun{v}{x,y}$ tends to $0$ as $\fun{\topdist}{x,y} \to \infty$, then
for every finite $X\Subset\Gamma$ and for $X=\Gamma$, $\alpha_X$ is a group of $\ast$-automorphisms of
$\oa{B}_X$, whose time-$t$ member is $\alpha_{X,t}$ for every $t \in \fldreal$.
No continuity in $t$ is asserted.
\end{fact}

The particle-sector Dyson expansion and the consistency maps of the inverse limit prove Fact \ref{fact:dll-buchholz-automorphism}; see \cite[Proposition 3.1 and equations (21)--(29)]{DeuchertLampartLemm001}. Example 3.2 of \cite{DeuchertLampartLemm001} proves that the group need not be weakly continuous. The problem addressed here is not existence of \(\alpha_{\Gamma}\) on \(\oa{B}\) but approximation of its time-\(t\) map \(\alpha_{\Gamma,t}\) by \(\alpha_{\Lambda,t}\) on local gauge-invariant observables.

\subsection{The obstruction in the DLL approximation argument}\label{the-obstruction-in-the-dll-approximation-argument}

Theorem 4.1 of \cite{DeuchertLampartLemm001} derives finite-volume approximation only after assuming a volume-uniform local particle-number estimate. Its proof propagates the assumed estimate, uses it to remove a local occupation cutoff, and then applies a state-independent finite-range propagation bound to the cutoff dynamics. The equilibrium argument in Theorem 5.2 of the same paper inherits this dependence.

The present proof does not use either conditional theorem. The functional-integral estimate proved in this paper controls the cutoff tails directly for the Gibbs family. We then combine that control with the finite-range shell expansion and pass the KMS boundary functions to the accumulation state.

\section{Finite-Volume and Infinite-Volume Equilibrium States}\label{finite-volume-and-infinite-volume-equilibrium-states}

The equilibrium analysis begins with the trace-class Gibbs operator on a finite Fock space and passes its KMS boundary relation to a state of the infinite Buchholz algebra.

\subsection{Superstability and finite-volume Gibbs states}\label{superstability-and-finite-volume-gibbs-states}

Fix \(\beta > 0\) and a bounded family \(\seq{\smchemicalpotential_{\Lambda}}{\Lambda \Subset \Gamma}\) of chemical potentials. The finite-volume partition function and Gibbs state are \begin{equation}
\begin{aligned}
Z_{\Lambda,\beta}
& =
\sqfun{\trace_{\sphilb{F}_{\Lambda}}}{
\fnexp{-\beta K_{\Lambda,\smchemicalpotential_{\Lambda}}}},
\\
\fun{\oastate[\psi_{\Lambda}]}{A}
& =
\frac{
\sqfun{\trace_{\sphilb{F}_{\Lambda}}}{
A\fnexp{-\beta K_{\Lambda,\smchemicalpotential_{\Lambda}}}}}
{Z_{\Lambda,\beta}}.
\end{aligned}
\label{eq:fock-finite-volume-gibbs-state}
\end{equation}

\begin{prop}[Finite-volume Gibbs existence]\label{prop:fock-finite-volume-gibbs-existence}
For every finite $\Lambda \Subset \Gamma$ and $\beta > 0$, the operator
$\fnexp{-\beta K_{\Lambda,\smchemicalpotential_{\Lambda}}}$ is trace class and
\eqref{eq:fock-finite-volume-gibbs-state} defines a normal state on $\opspbddlin{\sphilb{F}_{\Lambda}}$.
Its restriction to $\fun{\oaresolventalgebra}{\Lambda}^{\gamma}$ satisfies the finite-volume $\beta$-KMS
boundary relation for the number-conserving dynamics.
\end{prop}

\begin{proof}
Let
$c_U = U + C_h + \sup_{\Lambda \Subset \Gamma}\abs{\smchemicalpotential_{\Lambda}}$, where
$C_h$ is the uniform degree bound in \eqref{eq:fock-uniform-degree-bound}.
Since $w \geq 0$ and the second quantization of $\physham[h]_{\Lambda}$ is bounded below by
$-C_h \opfocknumber_{\Lambda}$, one has the quadratic-form bound
$$\begin{aligned}
K_{\Lambda,\smchemicalpotential_{\Lambda}}
& \geq
U \sum_{x \in \Lambda}N_x^2 - c_U \sum_{x \in \Lambda}N_x
\end{aligned}$$
on $\sphilb{D}_{\Lambda,\txtfin}$ defined in
\eqref{eq:fock-finite-particle-subspace}.
For every occupation configuration
$\seq{\nu_x}{x \in \Lambda} \in \monnat^{\Lambda}$, the scalar estimate is
$$\begin{aligned}
U \sum_{x \in \Lambda}\nu_x^2 - c_U \sum_{x \in \Lambda}\nu_x
& \geq
\frac{U}{2}\sum_{x \in \Lambda}\nu_x^2
-
\frac{c_U^2}{2U}\abscard{\Lambda}.
\end{aligned}$$
The min--max principle and the occupation-number basis
\eqref{eq:fock-occupation-number-basis} give the trace majorant
$$\begin{aligned}
Z_{\Lambda,\beta}
& \leq
\fnexp{\frac{\beta c_U^2}{2U}\abscard{\Lambda}}
\prod_{x \in \Lambda}
\sum_{n \in \monnat}\fnexp{-\frac{\beta U}{2}n^2}
<
\infty.
\end{aligned}$$
The Gibbs density matrix is consequently well-defined.
Trace cyclicity for analytic elements proves the finite-volume KMS boundary relation.
On gauge-invariant observables, adding $-\smchemicalpotential_{\Lambda}\opfocknumber_{\Lambda}$ to the generator
does not change the real-time action because $\opfocknumber_{\Lambda}$ commutes with those observables.
\end{proof}

\subsection{Infinite-volume equilibrium problem}\label{infinite-volume-equilibrium-problem}

The Fock factorization \(\sphilb{F}_{\Gamma} \cong \sphilb{F}_{\Lambda} \otimes \sphilb{F}_{\Lambda^{\mathrm{c}}}\) and the vacuum vector \(\Omega_{\Lambda^{\mathrm{c}}}\) define the state from equation (89) of \cite{DeuchertLampartLemm001} by \begin{equation}
\begin{aligned}
\overline{\rho}_{\Lambda}
& =
\frac{\fnexp{-\beta K_{\Lambda,\smchemicalpotential_{\Lambda}}}}
{Z_{\Lambda,\beta}}
\otimes
\ketbra{\Omega_{\Lambda^{\mathrm{c}}}}{\Omega_{\Lambda^{\mathrm{c}}}},
\\
\fun{\oastate[\overline{\psi}_{\Lambda}]}{A}
& =
\sqfun{\trace_{\sphilb{F}_{\Gamma}}}{A\overline{\rho}_{\Lambda}},
\quad
A \in \oa{B}.
\end{aligned}
\label{eq:fock-vacuum-extended-gibbs-state}
\end{equation} The state space of \(\oa{B}_{\Gamma}\) is weak-\(\ast\) compact, and the net \(\seq{\oastate[\overline{\psi}_{\Lambda}]}{\Lambda \Subset \Gamma}\) has weak-\(\ast\) accumulation points. Weak-\(\ast\) compactness alone does not pass the KMS boundary condition because the dynamics is controlled through unbounded local occupations before the cutoffs are removed.

No continuity of \(\alpha_{\Gamma}\) on \(\oa{B}\) is assumed. The equilibrium conclusion to be proved is a state \(\oastate[\psi]\) on \(\oa{B}\) that is invariant under \(\alpha_{\Gamma}\) on \(\fun{\oaresolventalgebra}{\Gamma}^{\gamma}\) and satisfies the corresponding KMS boundary relation there. For \(A,B \in \fun{\oaresolventalgebra}{\Gamma}^{\gamma}\), this means the existence of a bounded scalar function \(F_{A,B}\) that is holomorphic on \(-\beta < \opimag z < 0\), continuous on the closed strip, and satisfies \begin{equation}
\begin{aligned}
\fun{F_{A,B}}{t}
& =
\fun{\oastate[\psi]}{\fun{\alpha_{\Gamma,t}}{A}B},
&
\fun{F_{A,B}}{t-\imunit\beta}
& =
\fun{\oastate[\psi]}{B\fun{\alpha_{\Gamma,t}}{A}}.
\end{aligned}
\label{eq:fock-infinite-volume-kms-boundary}
\end{equation} Continuity in this statement concerns only the scalar function \(F_{A,B}\) and is part of the conclusion. It does not assert point-norm, weak, or strong continuity of the automorphism group on \(\oa{B}\).

\section{All-Temperature Particle-Number Bounds at Low Activity}\label{all-temperature-particle-number-bounds-at-low-activity}

The discrete functional integral compares marked open paths with an interacting gas of unmarked loops. The argument uses nonnegativity of \(v\) and a uniform gap below the one-particle band.

\subsection{Sub-Markov kernels and the chemical-potential gap}\label{sub-markov-kernels-and-the-chemical-potential-gap}

Let \(C_h\) be the uniform degree bound defined by \eqref{eq:fock-uniform-degree-bound}. Assume that there is \(\varepsilon > 0\) such that \begin{equation}
\begin{aligned}
\smchemicalpotential_{\Lambda}
& \leq
-C_h - \varepsilon
\end{aligned}
\label{eq:fock-path-chemical-potential-gap}
\end{equation} for every finite \(\Lambda \Subset \Gamma\). The elementary probabilistic definitions are included because the bridge measure and its normalization enter the particle-number estimate directly.

\begin{defn}[Finite state space and sub-Markovian kernels]\label{def:fock-finite-state-space-submarkov}
A finite state space is a finite set $E$ whose elements are called states.
A matrix kernel $R$ on the finite state space $E$ is substochastic when
$$\begin{aligned}
\fun{R}{x,y}
& \geq
0,
&
\sum_{y\in E}\fun{R}{x,y}
& \leq
1
\end{aligned}$$
for every $x,y\in E$.
It is stochastic when every row sum is equal to $1$.
A stochastic kernel is a transition kernel that conserves mass.

A family of matrix kernels $P_t$, $t\geq0$, on $E$ is sub-Markovian when
$$\begin{aligned}
\fun{P_0}{x,y}
& =
\fun{\fndef{\setone{x}}}{y},
\\
\sum_{z\in E}
\fun{P_s}{x,z}\fun{P_t}{z,y}
& =
\fun{P_{s+t}}{x,y},
\\
\fun{P_t}{x,y}
& \geq
0,
\\
\sum_{y\in E}\fun{P_t}{x,y}
&\leq
1
\end{aligned}$$
for every $s,t\geq0$ and $x,y\in E$.
\end{defn}

\begin{defn}[Cemetery state and absorbing state]\label{def:fock-cemetery-absorbing-state}
Let $E$ be a finite state space in the sense of
Definition \ref{def:fock-finite-state-space-submarkov}, and let $\partial\notin E$.
For a substochastic kernel $R$ on $E$, define a stochastic kernel $\widehat{R}$ on
$E\cup\setone{\partial}$ by
$$\begin{aligned}
\fun{\widehat{R}}{x,y}
& =
\fun{R}{x,y},
&
x,y
& \in E,
\\
\fun{\widehat{R}}{x,\partial}
& =
1-\sum_{y\in E}\fun{R}{x,y},
&
x
& \in E,
\\
\fun{\widehat{R}}{\partial,\partial}
& =
1,
\\ %%%%%%%%%%%%%%%%
\fun{\widehat{R}}{\partial,y}
& =
0,
\quad y\in E.
\end{aligned}$$
The extra point $\partial$ is called the cemetery state: the missing row mass of $R$ is the probability of a
transition from $x\in E$ to $\partial$.
The two identities in the last two lines mean that $\partial$ is an absorbing state for $\widehat{R}$; after a
transition to $\partial$, the kernel assigns probability $1$ to remaining at $\partial$ and probability $0$ to
returning to $E$.

For a process $Z=\seq{Z_t}{t\geq0}$ with values in $E\cup\setone{\partial}$, the state $\partial$ is absorbing
when, for every sample point and all $0\leq s\leq t$,
$$\begin{aligned}
Z_s=\partial
& \implies
Z_t=\partial.
\end{aligned}$$
\end{defn}

\begin{defn}[Discrete-time Markov chain]
Let $E$ be a finite state space in the sense of Definition \ref{def:fock-finite-state-space-submarkov}, let $R$
be a stochastic transition kernel on $E$, and let $\msrprb$ be a probability measure on a measurable space.
An $E$-valued sequence $\seq{Y_n}{n\in\monnat}$ is a discrete-time Markov chain with transition kernel $R$ when
$$\begin{aligned}
\funcond{\msrprb}{Y_{n+1}=y}{Y_0,\ldots,Y_n}
& =
\fun{R}{Y_n,y}
\end{aligned}$$
almost surely for every $n\in\monnat$ and $y\in E$.
The initial law is specified separately by the distribution of $Y_0$.
\end{defn}

\begin{defn}[Time-homogeneous Markov process]\label{def:fock-time-homogeneous-markov-process}
Let $E$ be a finite state space, let
$Z=\seq{Z_t}{t\geq0}$ be an $E$-valued process on a probability space
$\rbk{\Sigma,\mblfml{F},\msrprb}$, and let $P_t$, $t\geq0$, be stochastic kernels on $E$.
For $s\geq0$, let $\mblfml{F}_{Z,s}$ be the sigma-algebra generated by $Z_u$, $0\leq u\leq s$.
The process $Z$ is a time-homogeneous Markov process with transition kernels $P_t$ when, for every
$0\leq s\leq t$ and $y\in E$,
$$\begin{aligned}
\sqfuncond{\prbexp_{\msrprb}}{
\fndef{\setone{Z_t=y}}
}{\mblfml{F}_{Z,s}}
& =
\fun{P_{t-s}}{Z_s,y}
\end{aligned}$$
almost surely.
\end{defn}

\begin{defn}[Poisson process]
Let $\msrprb$ be a probability measure on a measurable space.
A process $N_t$, $t\geq0$, is a Poisson process of rate $c>0$ when $N_0=0$ almost surely, its increments over
pairwise disjoint time intervals are independent, and
$$\begin{aligned}
\fun{\msrprb}{N_t-N_s=k}
& =
\napiernum^{-c\rbk{t-s}}
\frac{\rbk{c\rbk{t-s}}^{k}}{k!}
\end{aligned}$$
for every $0\leq s<t$ and $k\in\monnat$.
Its paths are right-continuous, nondecreasing, integer-valued, and have jumps of size $1$.
\end{defn}

\begin{defn}[Exponential law]\label{def:fock-exponential-law}
Let $\msrprb$ be a probability measure on a measurable space.
A nonnegative random variable $T$ has the exponential law of rate $r>0$ when
$$\begin{aligned}
\fun{\msrprb}{T>t}
& =
\napiernum^{-rt}
\end{aligned}$$
for every $t\geq0$.
\end{defn}

\begin{defn}[Càdlàg path]\label{def:fock-cadlag-path}
Let $S$ be a metric space and let $I=\closedinterval{0}{a}$ for some $a>0$, or
$I=\rightopeninterval{0}{\infty}$.
A path $z\colon I\to S$ is càdlàg when it is right-continuous and has a left limit at every positive time in
$I$.
A process $Z=\seq{Z_t}{t\in I}$ is càdlàg when the path
$t\mapsto\fun{Z_t}{\omega}$ is càdlàg for every sample point $\omega$.
\end{defn}

\begin{defn}[Process killed at a nonnegative random variable]\label{def:fock-killed-process}
Let $E$ be a finite state space in the sense of
Definition \ref{def:fock-finite-state-space-submarkov}, and let $\partial$ be an absorbing state in the sense of
Definition \ref{def:fock-cemetery-absorbing-state}.
Let $X=\seq{X_t}{t\geq0}$ be a process with values in $E\cup\setone{\partial}$, and let $T$ be a nonnegative
random variable on the same probability space.
A process $\xi=\seq{\xi_t}{t\geq0}$ is called $X$ killed at $T$, or a killed process, when, for every sample point
$\omega$ and every $t\geq0$,
$$\begin{aligned}
\fun{\xi_t}{\omega}
& =
\begin{cases}
\fun{X_t}{\omega},
&
0\leq t<\fun{T}{\omega},
\\
\partial,
&
t\geq\fun{T}{\omega}.
\end{cases}
\end{aligned}$$
\end{defn}

\begin{lem}[Uniformization of the hopping kernel]\label{lem:fock-submarkov-uniformization}
For a finite $\Lambda\Subset\Gamma$, define
\begin{equation}
\begin{aligned}
Q_{\Lambda}
& =
-\physham[h]_{\Lambda}-C_h.
\end{aligned}
\label{eq:fock-killed-generator}
\end{equation}
The family $\napiernum^{tQ_{\Lambda}}$, $t\geq0$, is sub-Markovian.
For $R_{\Lambda}=-C_h^{-1}\physham[h]_{\Lambda}$, one has
\begin{equation}
\begin{aligned}
\napiernum^{tQ_{\Lambda}}
& =
\napiernum^{-tC_h}
\sum_{n\in\monnat}
\frac{\rbk{tC_h}^{n}}{n!}R_{\Lambda}^{n}.
\end{aligned}
\label{eq:fock-submarkov-uniformization}
\end{equation}
\end{lem}

\begin{proof}
The off-diagonal entries of $Q_{\Lambda}$ are nonnegative, and its row sums satisfy
$$\begin{aligned}
\sum_{y\in\Lambda}\rbk{Q_{\Lambda}}_{x,y}
& =
\abscard{\set{y\in\Lambda}{y\sim x}}-C_h
\leq
0.
\end{aligned}$$
The matrix $R_{\Lambda}$ has nonnegative entries and row sums at most $1$.
Every power $R_{\Lambda}^{n}$ has the same two properties.
Expanding $\napiernum^{tC_hR_{\Lambda}}$ proves
\eqref{eq:fock-submarkov-uniformization}.
The Poisson coefficients in that formula are nonnegative and sum to $1$.
Their convex combination has nonnegative entries and row sums at most $1$ for every $t\geq0$.
The matrix-exponential semigroup law completes the sub-Markovianity assertion.
\end{proof}

\begin{lem}[Probabilistic representation of $p_{\Lambda,t}$]\label{lem:fock-killed-walk-kernel}
Assume \eqref{eq:fock-path-chemical-potential-gap}, and let $\Lambda\Subset\Gamma$ be finite.
Let $Q_{\Lambda}$ be the operator defined by \eqref{eq:fock-killed-generator}.
Define $r_{\Lambda}$ and $p_{\Lambda,t}$, $t\geq0$, by
\begin{equation}
\begin{aligned}
r_{\Lambda}
& =
-\rbk{\smchemicalpotential_{\Lambda}+C_h},
\\
p_{\Lambda,t}
& =
\napiernum^{-t\rbk{\physham[h]_{\Lambda}-\smchemicalpotential_{\Lambda}}}
=
\napiernum^{-tr_{\Lambda}}\napiernum^{tQ_{\Lambda}}.
\end{aligned}
\label{eq:fock-killed-kernel-definition}
\end{equation}
Then, for every $x\in\Lambda$, there are a probability space
$\rbk{\Sigma_{x},\mblfml{F}_{x},\msrprb_x}$, a nonnegative random variable
$\tau_{\Lambda}$ with exponential law of rate $r_{\Lambda}$, and a càdlàg Markov process
$\xi=\seq{\xi_t}{t\geq0}$ on $\Lambda\cup\setone{\partial}$ in the sense of
Definitions \ref{def:fock-cadlag-path} and \ref{def:fock-time-homogeneous-markov-process} such that
$$\begin{aligned}
\fun{\msrprb_x}{\xi_0=x}
& =
1,
&
\fun{\msrprb_x}{\tau_{\Lambda}>t}
& =
\napiernum^{-r_{\Lambda}t},
\quad
t\geq0,
\end{aligned}$$
and $\partial$ is an absorbing state for $\xi$ in the sense of
Definition \ref{def:fock-cemetery-absorbing-state}.
The subscript $x$ in $\msrprb_x$ records the initial state $x$.
Its transition probabilities satisfy
\begin{equation}
\begin{aligned}
\fun{\msrprb_x}{\xi_t=y}
& =
\fun{p_{\Lambda,t}}{x,y},
&
x,y
& \in
\Lambda,
\\
\fun{\msrprb_x}{\xi_t=\partial}
& =
1-
\sum_{y\in\Lambda}\fun{p_{\Lambda,t}}{x,y}.
\end{aligned}
\label{eq:fock-killed-walk-transition}
\end{equation}
For every $t\geq0$, the transition probabilities satisfy the uniform row estimate
\begin{equation}
\begin{aligned}
\sum_{y \in \Lambda}\fun{p_{\Lambda,t}}{x,y}
& \leq
\napiernum^{-t\varepsilon}.
\end{aligned}
\label{eq:fock-killed-walk-row-bound}
\end{equation}
\end{lem}

\begin{proof}
Let $R_{\Lambda}$ be the substochastic kernel in
Lemma \ref{lem:fock-submarkov-uniformization}.
Its stochastic extension to $\Lambda\cup\setone{\partial}$ is
$$\begin{aligned}
\fun{\widehat{R}_{\Lambda}}{x,y}
& =
\fun{R_{\Lambda}}{x,y},
&
x,y
& \in
\Lambda,
\\
\fun{\widehat{R}_{\Lambda}}{x,\partial}
& =
1-
\sum_{y\in\Lambda}\fun{R_{\Lambda}}{x,y},
&
x
& \in
\Lambda,
\\
\fun{\widehat{R}_{\Lambda}}{\partial,\partial}
& =
1,
\\
\fun{\widehat{R}_{\Lambda}}{\partial,y}
& =
0,
&
y
& \in
\Lambda.
\end{aligned}$$
Thus $\partial$ is the cemetery state of $R_{\Lambda}$ and is absorbing for
$\widehat{R}_{\Lambda}$ in the sense of Definition \ref{def:fock-cemetery-absorbing-state}.
Let $\msrprb_x^{Y}$ be the law of the discrete-time Markov chain with transition kernel
$\widehat{R}_{\Lambda}$ and initial state $Y_0=x$.
Let $\msrprb^{N}$ be the law of a Poisson process of rate $C_h$.
On the product space with law $\msrprb_x^{X}=\msrprb_x^{Y}\otimes\msrprb^{N}$, define
$X_t=Y_{N_t}$.
For $x,y\in\Lambda$, the restriction of every $n$-step kernel of $\widehat{R}_{\Lambda}$ to
$\Lambda\times\Lambda$ is $R_{\Lambda}^{n}$.
The assertion holds for $n=0$ because both kernels are the identity on $\Lambda$.
If it holds for $n$, the absorbing-state identity
$\fun{\widehat{R}_{\Lambda}}{\partial,y}=0$ from
Definition \ref{def:fock-cemetery-absorbing-state} gives the induction step
$$\begin{aligned}
&
\fun{\widehat{R}_{\Lambda}^{n+1}}{x,y}
=
\sum_{z\in\Lambda}
\fun{\widehat{R}_{\Lambda}^{n}}{x,z}
\fun{\widehat{R}_{\Lambda}}{z,y}
+
\fun{\widehat{R}_{\Lambda}^{n}}{x,\partial}
\fun{\widehat{R}_{\Lambda}}{\partial,y}
\\ %%%%%%%%%%%%%%%%
&=
\sum_{z\in\Lambda}
\fun{R_{\Lambda}^{n}}{x,z}
\fun{R_{\Lambda}}{z,y}
=
\fun{R_{\Lambda}^{n+1}}{x,y}.
\end{aligned}$$
The Markov property of $Y$ yields
$$\begin{aligned}
\fun{\msrprb_x^{Y}}{Y_n=y}
=
\fun{\widehat{R}_{\Lambda}^{n}}{x,y}
=
\fun{R_{\Lambda}^{n}}{x,y}
\end{aligned}$$
for every $n\in\monnat$ and $x,y\in\Lambda$.
The events $\setone{N_t=n}$, $n\in\monnat$, form a disjoint partition, and the product measure makes the
entire chain $Y$ independent of the Poisson process $N$.
The law of total probability and the Poisson distribution give
$$\begin{aligned}
&\fun{\msrprb_x^{X}}{X_t=y}
=
\sum_{n\in\monnat}
\fun{\msrprb_x^{X}}{X_t=y,\ N_t=n}
\\
& =
\sum_{n\in\monnat}
\fun{\msrprb_x^{X}}{Y_n=y,\ N_t=n}
=
\sum_{n\in\monnat}
\fun{\msrprb_x^{Y}}{Y_n=y}
\fun{\msrprb^{N}}{N_t=n}
\\
& =
\sum_{n\in\monnat}
\fun{R_{\Lambda}^{n}}{x,y}
\napiernum^{-tC_h}
\frac{\rbk{tC_h}^{n}}{n!}
=
\fun{
\napiernum^{-tC_h}
\sum_{n\in\monnat}
\frac{\rbk{tC_h}^{n}}{n!}R_{\Lambda}^{n}
}{x,y}
=
\fun{\napiernum^{tQ_{\Lambda}}}{x,y}.
\end{aligned}$$
The last equality is the uniformization formula
\eqref{eq:fock-submarkov-uniformization}.

\begin{samepage}
Let $\rbk{\Sigma_x^{X},\mblfml{F}_x^{X},\msrprb_x^{X}}$ denote the probability space carrying $X$ constructed
here, and let
$\rbk{\Sigma_{\tau_{\Lambda}},\mblfml{F}_{\tau_{\Lambda}},
\msrprb_{r_{\Lambda}}^{\tau_{\Lambda}}}$ carry a nonnegative random variable
$\tau_{\Lambda}$ with exponential law of rate $r_{\Lambda}$ in the sense of
Definition \ref{def:fock-exponential-law}.
The inequality $r_{\Lambda}\geq\varepsilon>0$ guarantees that this probability law is defined.
The probability space in the assertion is the product triple
$$\begin{aligned}
\pairbk{\Sigma_x,
\mblfml{F}_x,
\msrprb_x}
=
\pairbk{\Sigma_x^{X}\times\Sigma_{\tau_{\Lambda}},
\mblfml{F}_x^{X}\otimes\mblfml{F}_{\tau_{\Lambda}},
\msrprb_x^{X}\otimes\msrprb_{r_{\Lambda}}^{\tau_{\Lambda}}},
\end{aligned}$$
\end{samepage}
where $X$ and $\tau_{\Lambda}$ are lifted to $\Sigma_x$ by the coordinate projections.
For $\omega_X\in\Sigma_x^{X}$ and
$\omega_{\tau_{\Lambda}}\in\Sigma_{\tau_{\Lambda}}$, define $\xi = \net{\xi_t}{t \geq 0}$ as $X$ killed at $\tau_{\Lambda}$ in the sense
of Definition \ref{def:fock-killed-process} by
\begin{equation}
\begin{aligned}
\fun{\xi_t}{\omega_X,\omega_{\tau_{\Lambda}}}
& =
\begin{cases}
\fun{X_t}{\omega_X},
&
0\leq t<\fun{\tau_{\Lambda}}{\omega_{\tau_{\Lambda}}},
\\
\partial,
&
t\geq\fun{\tau_{\Lambda}}{\omega_{\tau_{\Lambda}}}.
\end{cases}
\end{aligned}
\label{eq:fock-killed-process-definition}
\end{equation}
Because $\tau_{\Lambda}$ is a random variable on $\Sigma_x$, it is not a deterministic parameter indexing the
process $\xi$.
For $y\in\Lambda$, the endpoint event is
$$\begin{aligned}
\setone{\xi_t=y}
& =
\setone{\tau_{\Lambda}>t}\cap\setone{X_t=y}.
\end{aligned}$$
The product law makes $\tau_{\Lambda}$ independent of the entire process $X$.
The definition of the exponential law and the transition identity for $X$ give
$$\begin{aligned}
\fun{\msrprb_x}{\xi_t=y}
& =
\fun{\msrprb_x}{\tau_{\Lambda}>t}
\fun{\msrprb_x}{X_t=y}
=
\napiernum^{-r_{\Lambda}t}
\fun{\napiernum^{tQ_{\Lambda}}}{x,y}
=
\fun{p_{\Lambda,t}}{x,y}.
\end{aligned}$$
The sub-Markov row sum and $r_{\Lambda}\geq\varepsilon$ give the claimed uniform estimate through
$$\begin{aligned}
\sum_{y\in\Lambda}\fun{p_{\Lambda,t}}{x,y}
=
\napiernum^{-r_{\Lambda}t}
\sum_{y\in\Lambda}
\fun{\napiernum^{tQ_{\Lambda}}}{x,y}
\leq
\napiernum^{-r_{\Lambda}t}
\leq
\napiernum^{-\varepsilon t}.
\end{aligned}$$
Subtracting this row sum from $1$ gives the cemetery-state formula in
\eqref{eq:fock-killed-walk-transition}.
The Poisson time change preserves the Markov property of the chain $Y$.
For $s,t\geq0$, Definition \ref{def:fock-exponential-law} gives
$$\begin{aligned}
\funcond{\msrprb_x}{\tau_{\Lambda}>s+t}{\tau_{\Lambda}>s}
& =
\frac{\napiernum^{-r_{\Lambda}\rbk{s+t}}}{\napiernum^{-r_{\Lambda}s}}
=
\napiernum^{-r_{\Lambda}t}.
\end{aligned}$$
Together with the Markov property of $X$ and the independence of $X$ and $\tau_{\Lambda}$, this identity proves
that $\xi$ is Markov.
The identities $\fun{\widehat{R}_{\Lambda}}{\partial,\partial}=1$ and
\eqref{eq:fock-killed-process-definition} show directly that $\xi_s=\partial$ implies $\xi_t=\partial$ for
every $t\geq s$; hence $\partial$ is absorbing for $\xi$ according to
Definition \ref{def:fock-cemetery-absorbing-state}.
\end{proof}

The continuous-time bridge path space is defined by \begin{equation}
\begin{aligned}
\Omega_{\Lambda;x,y}^{t}
& =
\set{\omega\colon\closedinterval{0}{t}\to\Lambda}
{\text{$\omega$ is càdlàg, has finitely many jumps,
$\omega_0=x $ and $\omega_t=y$}}.
\end{aligned}
\label{eq:fock-continuous-bridge-path-space}
\end{equation} For \(\omega\in\Omega_{\Lambda;x,y}^{t}\) and \(0\leq s\leq t\), the notation \(\omega_s\) denotes the value of \(\omega\) at time \(s\). This path space carries the sigma-algebra generated by its coordinate evaluation maps.

\begin{lem}[Unnormalized bridge measure]\label{lem:fock-unnormalized-bridge-measure}
Assume \eqref{eq:fock-path-chemical-potential-gap}, and let $\Lambda\Subset\Gamma$ be finite.
Fix $x,y\in\Lambda$ and $t>0$.
Let $p_{\Lambda,t}$ be defined by \eqref{eq:fock-killed-kernel-definition}, and let
$\rbk{\Sigma_x,\mblfml{F}_x,\msrprb_x}$ and $\xi$ be the probability space and process whose existence is
asserted in Lemma \ref{lem:fock-killed-walk-kernel} for the initial state $x$.
Let $\Omega_{\Lambda;x,y}^{t}$ be the measurable path space defined by
\eqref{eq:fock-continuous-bridge-path-space}.
There is a finite positive measure $\msrcal{W}_{\Lambda;x,y}^{t}$ on this space such that, for every
$n\in\semigrposint$, $0=s_0<s_1<\cdots<s_n=t$, and
$z_0=x,z_1,\ldots,z_{n-1},z_n=y$ in $\Lambda$,
\begin{equation}
\begin{aligned}
\fun{\msrcal{W}_{\Lambda;x,y}^{t}}{
\set{\omega\in\Omega_{\Lambda;x,y}^{t}}
{\omega_{s_j}=z_j,\ 1\leq j\leq n-1}}
& =
\prod_{j=0}^{n-1}
\fun{p_{\Lambda,s_{j+1}-s_j}}{z_j,z_{j+1}}.
\end{aligned}
\label{eq:fock-bridge-cylinder-distribution}
\end{equation}
Its total mass is
\begin{equation}
\begin{aligned}
\fun{\msrcal{W}_{\Lambda;x,y}^{t}}{\Omega_{\Lambda;x,y}^{t}}
& =
\fun{p_{\Lambda,t}}{x,y}.
\end{aligned}
\label{eq:fock-bridge-total-mass}
\end{equation}
For $s,t>0$, $x,y,z\in\Lambda$,
$\omega\in\Omega_{\Lambda;x,z}^{s}$, and
$\eta\in\Omega_{\Lambda;z,y}^{t}$, define their concatenation by
$$\begin{aligned}
\rbk{\omega\diamond_s\eta}_u
& =
\begin{cases}
\omega_u, & 0\leq u\leq s,\\
\eta_{u-s}, & s<u\leq s+t.
\end{cases}
\end{aligned}$$
For every nonnegative measurable function $F$ on
$\bigsqcup_{x,y\in\Lambda}\Omega_{\Lambda;x,y}^{s+t}$, the bridge measures satisfy
\begin{equation}
\begin{aligned}
&
\sum_{z\in\Lambda}
\int_{\Omega_{\Lambda;x,z}^{s}}
\int_{\Omega_{\Lambda;z,y}^{t}}
\fun{F}{\omega\diamond_s\eta}
\opdmsr{\msrcal{W}_{\Lambda;z,y}^{t}}(\eta)
\opdmsr{\msrcal{W}_{\Lambda;x,z}^{s}}(\omega)
=
\int_{\Omega_{\Lambda;x,y}^{s+t}}
\fun{F}{\zeta}
\opdmsr{\msrcal{W}_{\Lambda;x,y}^{s+t}}(\zeta).
\end{aligned}
\label{eq:fock-bridge-convolution-law}
\end{equation}
\end{lem}

\begin{proof}
For a measurable set $A\subset\Omega_{\Lambda;x,y}^{t}$, define
$$\begin{aligned}
\fun{\msrcal{W}_{\Lambda;x,y}^{t}}{A}
& =
\fun{\msrprb_x}{
\setone{\fnrestr{\xi}{\closedinterval{0}{t}}\in A}}.
\end{aligned}$$
This restriction of the path law is a finite positive measure.
Fix $0=s_0<s_1<\cdots<s_n=t$ and
$z_0=x,z_1,\ldots,z_{n-1},z_n=y$ in $\Lambda$.
For $0\leq u\leq t$, let $\mblfml{F}_{\xi,u}$ be the sigma-algebra generated by
$\xi_v$, $0\leq v\leq u$.
Let $G_0$ be the whole sample space on which $\xi$ is defined.
For $1\leq j\leq n$, define the event
$$\begin{aligned}
G_j
& =
\bigcap_{\ell=1}^{j}
\setone{\xi_{s_{\ell}}=z_{\ell}}.
\end{aligned}$$
On $G_{j-1}$, the process satisfies
$\xi_{s_{j-1}}=z_{j-1}\in\Lambda$.
The time-homogeneous Markov property in
Definition \ref{def:fock-time-homogeneous-markov-process} and the transition formula
\eqref{eq:fock-killed-walk-transition} identify the conditional probability on this event as
$$\begin{aligned}
\fndef{G_{j-1}}
\sqfuncond{\prbexp_{\msrprb_x}}{
\fndef{\setone{\xi_{s_j}=z_j}}
}{\mblfml{F}_{\xi,s_{j-1}}}
& =
\fndef{G_{j-1}}
\fun{p_{\Lambda,s_j-s_{j-1}}}{z_{j-1},z_j}
\end{aligned}$$
for $1\leq j\leq n$.
Since $G_{j-1}\in\mblfml{F}_{\xi,s_{j-1}}$, the defining property of conditional expectation gives
$$\begin{aligned}
\fun{\msrprb_x}{G_j}
& =
\sqfun{\prbexp_{\msrprb_x}}{
\fndef{G_{j-1}}
\fndef{\setone{\xi_{s_j}=z_j}}
}
=
\sqfun{\prbexp_{\msrprb_x}}{
\fndef{G_{j-1}}
\sqfuncond{\prbexp_{\msrprb_x}}{
\fndef{\setone{\xi_{s_j}=z_j}}
}{\mblfml{F}_{\xi,s_{j-1}}}
}
\\
& =
\sqfun{\prbexp_{\msrprb_x}}{
\fndef{G_{j-1}}
\fun{p_{\Lambda,s_j-s_{j-1}}}{z_{j-1},z_j}
}
=
\fun{p_{\Lambda,s_j-s_{j-1}}}{z_{j-1},z_j}
\fun{\msrprb_x}{G_{j-1}}.
\end{aligned}$$
Since $G_0$ is the whole sample space, $\fun{\msrprb_x}{G_0}=1$.
Iterating the recursion from this initial value yields
$$\begin{aligned}
\fun{\msrprb_x}{G_n}
=
\prod_{j=1}^{n}
\fun{p_{\Lambda,s_j-s_{j-1}}}{z_{j-1},z_j}
=
\prod_{j=0}^{n-1}
\fun{p_{\Lambda,s_{j+1}-s_j}}{z_j,z_{j+1}}.
\end{aligned}$$
Because $\partial$ is absorbing according to Definition \ref{def:fock-cemetery-absorbing-state}, the identity
$z_n=y\in\Lambda$ shows that the event $G_n$ implies
$\xi_s\in\Lambda$ for every $0\leq s\leq t$.
Since $\xi$ is càdlàg and takes values in the finite state space $\Lambda\cup\setone{\partial}$, its paths have
finitely many jumps on every bounded time interval.
The inverse image of the cylinder set in
\eqref{eq:fock-bridge-cylinder-distribution} under
$\fnrestr{\xi}{\closedinterval{0}{t}}$ is $G_n$ up to a $\msrprb_x$-null set.
The definition of $\msrcal{W}_{\Lambda;x,y}^{t}$ and the formula for
$\fun{\msrprb_x}{G_n}$ prove
\eqref{eq:fock-bridge-cylinder-distribution}.
Taking no intermediate times in that formula gives \eqref{eq:fock-bridge-total-mass}.
For an indicator function of a cylinder set, the left-hand side of
\eqref{eq:fock-bridge-convolution-law} is evaluated by
\eqref{eq:fock-bridge-cylinder-distribution}; summation over the common value $z$ and the semigroup identity
in Definition \ref{def:fock-finite-state-space-submarkov} give the right-hand side.
Finite linear combinations and monotone convergence extend the identity to every nonnegative measurable
function $F$.
\end{proof}

\subsection{Open paths, loops, and the interacting normalization}\label{open-paths-loops-and-the-interacting-normalization}

The kernels from Lemma \ref{lem:fock-killed-walk-kernel} and the bridge measures from Lemma \ref{lem:fock-unnormalized-bridge-measure} provide the positive path-space reference measure for the interaction.

\begin{defn}[Finite open-path tuple]\label{def:fock-finite-open-path-tuple}
Let $\Lambda\Subset\Gamma$ be finite, let $\beta>0$, and let $m\in\monnat$.
For $1\leq i\leq m$, let $q_i\in\semigrposint$ and $x_i,y_i\in\Lambda$.
A finite open-path tuple is a tuple
$\boldsymbol{\omega}=\rbk{\omega_1,\ldots,\omega_m}$ satisfying
$$\begin{aligned}
\omega_i
\in
\Omega_{\Lambda;x_i,y_i}^{q_i\beta},
\quad
1\leq i\leq m,
\end{aligned}$$
where the path spaces are defined by \eqref{eq:fock-continuous-bridge-path-space}.
For $m=0$, the tuple is empty.
\end{defn}

\begin{defn}[Occupation field]\label{def:fock-path-occupation-field}
For $q\in\semigrposint$, $x,y\in\Lambda$, and
$\omega\in\Omega_{\Lambda;x,y}^{q\beta}$, define the occupation field of $\omega$ by
$$\begin{aligned}
\fun{n_{\omega}}{s,z}
& =
\sum_{j = 0}^{q - 1}
\fun{\fndef{\setone{z}}}{\omega_{s + j\beta}},
\quad
0 \leq s \leq \beta,
\quad
z \in \Lambda.
\end{aligned}$$
Let $\boldsymbol{\xi}=\rbk{\xi_1,\ldots,\xi_m}$ be a finite open-path tuple in the sense of
Definition \ref{def:fock-finite-open-path-tuple}, with
$\xi_i\in\Omega_{\Lambda;x_i,y_i}^{q_i\beta}$ for $1\leq i\leq m$.
For $0\leq s\leq q_i\beta$, the symbol $\xi_{i,s}$ denotes the value of the $i$-th path at time $s$.
The occupation field of $\boldsymbol{\xi}$ on $\closedinterval{0}{\beta} \times \Lambda$ is defined by
\begin{equation}
\begin{aligned}
\fun{n_{\boldsymbol{\xi}}}{s,z}
=
\sum_{i = 1}^{m}\fun{n_{\xi_i}}{s,z}
=
\sum_{i = 1}^{m}
\sum_{j = 0}^{q_i - 1}
\fun{\fndef{\setone{z}}}{\xi_{i,s + j\beta}},
\quad
\rbk{s,z}
\in
\closedinterval{0}{\beta} \times \Lambda.
\end{aligned}
\label{eq:fock-path-tuple-occupation-field}
\end{equation}
For $m = 0$, both sums in \eqref{eq:fock-path-tuple-occupation-field} are zero.
\end{defn}

For fixed \(s\in\closedinterval{0}{\beta}\) and \(z\in\Lambda\), the single-path definition is equivalently \[\begin{aligned}
\fun{n_{\omega}}{s,z}
& =
\abscard{
\set{j\in\monnat}{0\leq j\leq q-1,\ \omega_{s+j\beta}=z}}.
\end{aligned}\] It counts how many of the \(q\) values \(\omega_s,\omega_{s+\beta},\ldots,\omega_{s+\rbk{q-1}\beta}\) are equal to \(z\), and hence \[\begin{aligned}
\sum_{z\in\Lambda}\fun{n_{\omega}}{s,z}
& =
q.
\end{aligned}\] For fixed \(\rbk{s,z}\in\closedinterval{0}{\beta}\times\Lambda\), this definition is the counting identity \[\begin{aligned}
\fun{n_{\boldsymbol{\xi}}}{s,z}
& =
\abscard{
\set{\rbk{i,j}}{
1\leq i\leq m,\ 0\leq j\leq q_i-1,\ \xi_{i,s+j\beta}=z}}.
\end{aligned}\] It therefore records the total number of path values at the lattice site \(z\) for the functional-integral time parameter \(s\), and \[\begin{aligned}
\sum_{z\in\Lambda}\fun{n_{\boldsymbol{\xi}}}{s,z}
& =
\sum_{i=1}^{m}q_i.
\end{aligned}\]

The interaction energy determined by Definition \ref{def:fock-path-occupation-field} and \eqref{eq:fock-path-tuple-occupation-field} is \begin{equation}
\begin{aligned}
\fun{\Phi_{\Lambda}}{\boldsymbol{\xi}}
& =
\int_{\closedinterval{0}{\beta}}
\sum_{x,y \in \Lambda}
\fun{v}{x,y}
\rbk{
\fun{n_{\boldsymbol{\xi}}}{s,x}\fun{n_{\boldsymbol{\xi}}}{s,y}
-
\fun{\fndef{\setone{x}}}{y}\fun{n_{\boldsymbol{\xi}}}{s,x}}
\opdmsr{s}.
\end{aligned}
\label{eq:fock-path-interaction-energy}
\end{equation} The subtraction in \eqref{eq:fock-path-interaction-energy} is the normal-ordering term. For finite open-path tuples in the sense of Definition \ref{def:fock-finite-open-path-tuple}, \(\boldsymbol{\omega}=\rbk{\omega_1,\ldots,\omega_k}\) and \(\boldsymbol{\zeta}=\rbk{\zeta_1,\ldots,\zeta_{\ell}}\), define the interaction energy of their concatenated tuple by \begin{equation}
\begin{aligned}
\fun{\Phi_{\Lambda}}{\boldsymbol{\omega},\boldsymbol{\zeta}}
& =
\fun{\Phi_{\Lambda}}{\omega_1,\ldots,\omega_k,\zeta_1,\ldots,\zeta_{\ell}},
&
k,\ell
& \in
\monnat.
\end{aligned}
\label{eq:fock-concatenated-path-interaction-energy}
\end{equation}

Using the bridge path spaces defined in \eqref{eq:fock-continuous-bridge-path-space}, define the loop path space by \begin{equation}
\begin{aligned}
\Omega_{\Lambda}^{\mathrm{loop}}
& =
\bigsqcup_{\substack{q \in \semigrposint \\ x \in \Lambda}}
\Omega_{\Lambda;x,x}^{q\beta}.
\end{aligned}
\label{eq:fock-loop-path-space}
\end{equation} The positive ideal loop measure \(\msrcal{L}_{\Lambda}\) on the space in \eqref{eq:fock-loop-path-space} is defined by \begin{equation}
\begin{aligned}
\int_{\Omega_{\Lambda}^{\mathrm{loop}}}
\fun{F}{\zeta}\opdmsr{\msrcal{L}_{\Lambda}}(\zeta)
& =
\sum_{q \in \semigrposint}\frac{1}{q}
\sum_{x \in \Lambda}
\int_{\Omega_{\Lambda;x,x}^{q\beta}}
\fun{F}{\omega}\opdmsr{\msrcal{W}_{\Lambda;x,x}^{q\beta}}(\omega)
\end{aligned}
\label{eq:fock-ideal-loop-measure}
\end{equation} for every nonnegative measurable \(F\). For a finite open-path tuple \(\boldsymbol{\omega}\), define \begin{equation}
\begin{aligned}
\fun{Z_{\Lambda}}{\boldsymbol{\omega}}
=
\sum_{\ell \in \monnat}\frac{1}{\ell!}
\int_{\rbk{\Omega_{\Lambda}^{\mathrm{loop}}}^{\ell}}
\napiernum^{-\fun{\Phi_{\Lambda}}{\boldsymbol{\omega},\boldsymbol{\zeta}}}
\opdmsr{\msrcal{L}_{\Lambda}^{\otimes\ell}}(\boldsymbol{\zeta}),
\quad
Z_{\Lambda}
=
\sum_{\ell \in \monnat}\frac{1}{\ell!}
\int_{\rbk{\Omega_{\Lambda}^{\mathrm{loop}}}^{\ell}}
\napiernum^{-\fun{\Phi_{\Lambda}}{\boldsymbol{\zeta}}}
\opdmsr{\msrcal{L}_{\Lambda}^{\otimes\ell}}(\boldsymbol{\zeta}).
\end{aligned}
\label{eq:fock-open-path-loop-partition-functions}
\end{equation} Set \(F=1\) in the loop-measure definition \eqref{eq:fock-ideal-loop-measure}. The bridge-mass identity \eqref{eq:fock-bridge-total-mass}, nonnegativity of \(p_{\Lambda,q\beta}\), and the row bound \eqref{eq:fock-killed-walk-row-bound} give \[\begin{aligned}
& \int_{\Omega_{\Lambda}^{\mathrm{loop}}}\opdmsr{\msrcal{L}_{\Lambda}}(\zeta)
=
\sum_{q \in \semigrposint}\frac{1}{q}
\sum_{x \in \Lambda}
\fun{\msrcal{W}_{\Lambda;x,x}^{q\beta}}{\Omega_{\Lambda;x,x}^{q\beta}}
=
\sum_{q \in \semigrposint}\frac{1}{q}
\sum_{x \in \Lambda}
\fun{p_{\Lambda,q\beta}}{x,x}
\\
& \leq
\sum_{q \in \semigrposint}\frac{1}{q}
\sum_{x \in \Lambda}
\sum_{y \in \Lambda}
\fun{p_{\Lambda,q\beta}}{x,y}
\leq
\sum_{q \in \semigrposint}\frac{1}{q}
\sum_{x \in \Lambda}
\napiernum^{-q\beta\varepsilon}
\\ %%%%%%%%%%%%%%%%
& =
\abscard{\Lambda}
\sum_{q \in \semigrposint}\frac{\napiernum^{-q\beta\varepsilon}}{q}
=
-\abscard{\Lambda}
\log\rbk{1-\napiernum^{-\beta\varepsilon}}
<
\infty.
\end{aligned}\] For every \(q\in\semigrposint\) and \(x,y\in\Lambda\), nonnegativity and \eqref{eq:fock-killed-walk-row-bound} give the pointwise estimate \[\begin{aligned}
\fun{p_{\Lambda,q\beta}}{x,y}
\leq
\sum_{z \in \Lambda}
\fun{p_{\Lambda,q\beta}}{x,z}
\leq
\napiernum^{-q\beta\varepsilon}.
\end{aligned}\] Summing this estimate over \(q\) and evaluating the resulting geometric series give the open-path bound \[\begin{aligned}
\sum_{q \in \semigrposint}\fun{p_{\Lambda,q\beta}}{x,y}
\leq
\sum_{q \in \semigrposint}\napiernum^{-q\beta\varepsilon}
=
\napiernum^{-\beta\varepsilon}
\sum_{m \in \monnat}
\rbk{\napiernum^{-\beta\varepsilon}}^{m}
=
\frac{\napiernum^{-\beta\varepsilon}}{1-\napiernum^{-\beta\varepsilon}}
=
\frac{1}{\napiernum^{\beta\varepsilon} - 1}.
\end{aligned}\] The loop series in \eqref{eq:fock-open-path-loop-partition-functions} is consequently absolutely convergent and \(1 \leq Z_{\Lambda} < \infty\).

For \(k \in \semigrposint\) and \(\boldsymbol{x},\boldsymbol{y} \in \Lambda^k\), the reduced density kernel is \begin{equation}
\begin{aligned}
\fun{\widetilde{\rho}_{\Lambda,k}}{\boldsymbol{x};\boldsymbol{y}}
& =
\fun{\oastate[\psi_{\Lambda}]}{
\faadj{a_{y_1}}\cdots\faadj{a_{y_k}}a_{x_k}\cdots a_{x_1}}.
\end{aligned}
\label{eq:fock-reduced-density-kernel}
\end{equation} The superstable trace majorant in Proposition \ref{prop:fock-finite-volume-gibbs-existence} makes the trace defining \eqref{eq:fock-reduced-density-kernel} finite.

The proof of the open-path expansion has nine parts.

\begin{enumerate}
\def\labelenumi{\arabic{enumi}.}

\item
  Lemma \ref{lem:fock-interaction-occupation-eigenvalue} computes the interaction in the occupation-number basis.
\item
  Lemma \ref{lem:fock-finite-slice-weight-properties} records the positivity and row bound of the finite-slice weights.
\item
  Lemma \ref{lem:fock-finite-slice-trace-formula} expands the finite-slice trace in particle coordinates.
\item
  Lemma \ref{lem:fock-creation-annihilation-sector-formula} inserts the creation--annihilation strings.
\item
  Lemma \ref{lem:fock-marked-index-permutation-decomposition} reorganizes the permutation sum by the index sequences generated by repeated application of a coordinate-index permutation.
\item
  Lemma \ref{lem:fock-finite-slice-open-path-expansion} converts the coordinate products into bridge and loop weights.
\item
  Lemma \ref{lem:fock-finite-slice-domination} supplies uniform bounds, and Lemma \ref{lem:fock-continuous-time-limit} uses them to take the continuous-time limit.
\item
  Lemma \ref{lem:fock-continuous-bosonic-cycle-decomposition} identifies the operator and loop partition functions.
\item
  Lemma \ref{lem:fock-loop-marked-chain-expansion} combines these results.
\end{enumerate}

For a finite \(\Lambda\Subset\Gamma\) and \(\boldsymbol{\nu}=\seq{\nu_x}{x\in\Lambda}\in\monnat^{\Lambda}\), define \begin{equation}
\begin{aligned}
\fun{E_{\Lambda}}{\boldsymbol{\nu}}
& =
\sum_{x,y\in\Lambda}
\fun{v}{x,y}
\rbk{\nu_x\nu_y-\fun{\fndef{\setone{x}}}{y}\nu_x}.
\end{aligned}
\label{eq:fock-finite-slice-occupation-energy}
\end{equation}

\begin{lem}[Interaction in the occupation-number basis]\label{lem:fock-interaction-occupation-eigenvalue}
Let $\Lambda\Subset\Gamma$ be finite and let
$\boldsymbol{\nu}=\seq{\nu_x}{x\in\Lambda}\in\monnat^{\Lambda}$.
The occupation-number vector $\Psi_{\Lambda,\boldsymbol{\nu}}$ in
\eqref{eq:fock-occupation-number-basis} satisfies
$$\begin{aligned}
V_{\Lambda}\Psi_{\Lambda,\boldsymbol{\nu}}
& =
\fun{E_{\Lambda}}{\boldsymbol{\nu}}\Psi_{\Lambda,\boldsymbol{\nu}}.
\end{aligned}$$
Moreover, $\fun{E_{\Lambda}}{\boldsymbol{\nu}}\geq0$.
\end{lem}

\begin{proof}
If $x\neq y$, the operators $a_xa_y$ contribute
$\sqrt{\nu_x\nu_y}$ and the operators $\faadj{a_x}\faadj{a_y}$ contribute the same factor.  Hence
$$\begin{aligned}
\faadj{a_x}\faadj{a_y}a_xa_y\Psi_{\Lambda,\boldsymbol{\nu}}
& =
\nu_x\nu_y\Psi_{\Lambda,\boldsymbol{\nu}}.
\end{aligned}$$
If $x=y$, the two annihilation operators contribute $\sqrt{\nu_x}\sqrt{\nu_x-1}$ and the two creation
operators contribute the same factor, so
$$\begin{aligned}
\rbk{\faadj{a_x}}^2a_x^2\Psi_{\Lambda,\boldsymbol{\nu}}
& =
\nu_x\rbk{\nu_x-1}\Psi_{\Lambda,\boldsymbol{\nu}}.
\end{aligned}$$
These two identities give \eqref{eq:fock-finite-slice-occupation-energy} after summation over the ordered pairs
$\rbk{x,y}\in\Lambda^2$.  Every summand is nonnegative because $v$ is pointwise nonnegative and
$\nu_x\in\monnat$.
\end{proof}

To make the limiting step explicit, fix the particle sector \(N\) and an integer \(m \geq 1\), split \(\closedinterval{0}{\beta}\) into \(m\) intervals, and insert the occupation-number basis \eqref{eq:fock-occupation-number-basis} between the factors in \begin{equation}
\begin{aligned}
T_{\Lambda,m}
=
\rbk{\fnexp{-\frac{\beta}{m}
\rbk{\physham_{\txtfr,\Lambda} -
\smchemicalpotential_{\Lambda}\opfocknumber_{\Lambda}}}
\cdot
\fnexp{-\frac{\beta}{m} V_{\Lambda}}}^{m}.
\end{aligned}
\label{eq:fock-finite-slice-trotter-product}
\end{equation} The operators \(\physham_{\txtfr,\Lambda}\) and \(V_{\Lambda}\) are defined in \eqref{eq:fock-bose-hubbard-hamiltonian}. Define the one-step kernel \(b_m\) by \begin{equation}
\begin{aligned}
\fun{b_m}{x,y}
& =
\fun{\fnexp{-\frac{\sminvtemperature}{m}
\rbk{\physham[h]_{\Lambda} - \smchemicalpotential_{\Lambda}}}}{x,y}.
\end{aligned}
\label{eq:fock-finite-slice-one-step-kernel}
\end{equation} For \(\boldsymbol{x}^{(r)} \in \Lambda^N\), write \(\nu_z^{(r)} = \sum_{i = 1}^{N}\fun{\fndef{\setone{z}}}{x_i^{(r)}}\). For \(\varpi \in \grsym{N}\), let \(\Omega_{\varpi,m}^{(N)}\) be the set of tuples \(\rbk{\boldsymbol{x}^{(0)},\ldots,\boldsymbol{x}^{(m)}}\) in \(\rbk{\Lambda^N}^{m+1}\) satisfying \(x_i^{(m)} = x_{\fun{\varpi}{i}}^{(0)}\) for every \(i\), and define \begin{equation}
\begin{aligned}
\fun{B_{N,m}}{\boldsymbol{x}^{(0)},\ldots,\boldsymbol{x}^{(m)}}
& =
\prod_{r = 0}^{m - 1}
\rbk{
\prod_{i = 1}^{N}
\fun{b_m}{x_i^{(r)},x_i^{(r + 1)}}}
\napiernum^{-\frac{\sminvtemperature}{m}\fun{E_{\Lambda}}{\nu^{(r)}}}.
\end{aligned}
\label{eq:fock-finite-slice-sector-weight}
\end{equation} For each time index \(r\), the product of the kernels \(b_m\) is the one-particle contribution from \(\boldsymbol{x}^{(r)}\) to \(\boldsymbol{x}^{(r+1)}\), and the exponential is the interaction contribution of the occupation numbers \(\nu^{(r)}\). Thus \(B_{N,m}\) assigns a number to each tuple \(\rbk{\boldsymbol{x}^{(0)},\ldots,\boldsymbol{x}^{(m)}}\) in \(\Omega_{\varpi,m}^{(N)}\).

\begin{lem}[Properties of the finite-slice weights]\label{lem:fock-finite-slice-weight-properties}
Assume \eqref{eq:fock-path-chemical-potential-gap}, let $\Lambda\Subset\Gamma$ be finite, and let $\beta>0$.
For every $m\in\semigrposint$ and $x\in\Lambda$, the kernel in
\eqref{eq:fock-finite-slice-one-step-kernel} satisfies
\begin{equation}
\begin{aligned}
\fun{b_m}{x,y}
& \geq0,
&& x,y\in\Lambda,
&
\sum_{y\in\Lambda}\fun{b_m}{x,y}
& \leq
\napiernum^{-\frac{\beta\varepsilon}{m}}.
\end{aligned}
\label{eq:fock-finite-slice-one-step-row-bound}
\end{equation}
For every $N\in\monnat$, $\varpi\in\grsym{N}$, and
$\boldsymbol{x}\in\Omega_{\varpi,m}^{(N)}$, the weight in
\eqref{eq:fock-finite-slice-sector-weight} satisfies
\begin{equation}
\begin{aligned}
0
\leq
\fun{B_{N,m}}{\boldsymbol{x}}
\leq
\prod_{r=0}^{m-1}\prod_{i=1}^{N}
\fun{b_m}{x_i^{(r)},x_i^{(r+1)}}.
\end{aligned}
\label{eq:fock-finite-slice-sector-weight-bound}
\end{equation}
\end{lem}

\begin{proof}
Definitions \eqref{eq:fock-killed-kernel-definition} and
\eqref{eq:fock-finite-slice-one-step-kernel} give
$b_m=p_{\Lambda,\beta/m}$.  Nonnegativity of $p_{\Lambda,t}$ and the row estimate
\eqref{eq:fock-killed-walk-row-bound}, evaluated at $t=\beta/m$, give
\eqref{eq:fock-finite-slice-one-step-row-bound}.
Lemma \ref{lem:fock-interaction-occupation-eigenvalue} gives
$\fun{E_{\Lambda}}{\nu^{(r)}}\geq0$ for every $r$.
Consequently every factor in \eqref{eq:fock-finite-slice-sector-weight} is nonnegative, and each exponential
factor is at most $1$.  This proves \eqref{eq:fock-finite-slice-sector-weight-bound}.
\end{proof}

\begin{lem}[Finite-slice trace formula]\label{lem:fock-finite-slice-trace-formula}
Assume \eqref{eq:fock-path-chemical-potential-gap}, let $\Lambda\Subset\Gamma$ be finite, and let $\beta>0$.
For every $N\in\monnat$ and $m\in\semigrposint$, the Trotter product
\eqref{eq:fock-finite-slice-trotter-product} satisfies
\begin{equation}
\begin{aligned}
\sqfun{\trace_{P_{\Lambda,N}\sphilb{F}_{\Lambda}}}{T_{\Lambda,m}}
& =
\frac{1}{N!}
\sum_{\varpi \in \grsym{N}}
\sum_{\boldsymbol{x} \in \Omega_{\varpi,m}^{(N)}}
\fun{B_{N,m}}{\boldsymbol{x}}.
\end{aligned}
\label{eq:fock-finite-slice-trace-formula}
\end{equation}
\end{lem}

\begin{proof}
Lemma \ref{lem:fock-interaction-occupation-eigenvalue} gives the diagonal matrix element of the interaction
factor, and \eqref{eq:fock-finite-slice-one-step-kernel} gives the one-particle matrix element of the free
factor.  Identify $P_{\Lambda,N}\sphilb{F}_{\Lambda}$ with the symmetric subspace of
$\lpseq^2\rbk{\Lambda^N}$.  For $\boldsymbol{x}\in\Lambda^N$, let $\delta_{\boldsymbol{x}}$ be the function on
$\Lambda^N$ that equals $1$ at $\boldsymbol{x}$ and $0$ at every other tuple.  Insertion of the basis
$\set{\delta_{\boldsymbol{x}}}{\boldsymbol{x}\in\Lambda^N}$ between consecutive factors of
\eqref{eq:fock-finite-slice-trotter-product} assigns the product
\eqref{eq:fock-finite-slice-sector-weight} to each sequence of intermediate coordinates.
The bosonic symmetrizer $S_N$ on $\lpseq^2\rbk{\Lambda^N}$ is defined by
$$\begin{aligned}
S_N\delta_{\boldsymbol{x}}
& =
\frac{1}{N!}
\sum_{\varpi\in\grsym{N}}
\delta_{\rbk{x_{\fun{\varpi}{1}},\ldots,x_{\fun{\varpi}{N}}}}.
\end{aligned}$$
Therefore its endpoint condition is
$x_i^{(m)}=x_{\fun{\varpi}{i}}^{(0)}$.  Summing the intermediate coordinates therefore gives
\eqref{eq:fock-finite-slice-trace-formula}.  The summands are nonnegative by
Lemma \ref{lem:fock-finite-slice-weight-properties}.
\end{proof}

For \(q \in \semigrposint\), let \(\Omega_{\Lambda;x,y}^{q,m}\) consist of sequences \(\omega = \seq{\omega_r}{r = 0}^{q m}\) with \(\omega_0 = x\) and \(\omega_{q m} = y\), and define \begin{equation}
\begin{aligned}
\fun{\msrcal{W}_{\Lambda;x,y}^{q\beta,(m)}}{\setone{\omega}}
& =
\prod_{r = 0}^{q m - 1}\fun{b_m}{\omega_r,\omega_{r + 1}},
\\
\fun{n_{\omega}^{(m)}}{r,z}
& =
\sum_{j = 0}^{q - 1}
\fun{\fndef{\setone{z}}}{\omega_{r + j m}},
\quad
0 \leq r \leq m - 1,
\\
\fun{\Phi_{\Lambda}^{(m)}}{\boldsymbol{\xi}}
& =
\frac{\sminvtemperature}{m}
\sum_{r = 0}^{m - 1}
\sum_{x,y \in \Lambda}
\fun{v}{x,y}
\rbk{
\fun{n_{\boldsymbol{\xi}}^{(m)}}{r,x}
\fun{n_{\boldsymbol{\xi}}^{(m)}}{r,y}
-
\fun{\fndef{\setone{x}}}{y}
\fun{n_{\boldsymbol{\xi}}^{(m)}}{r,x}}.
\end{aligned}
\label{eq:fock-discrete-bridge-path-space}
\end{equation} The discrete loop path space and its ideal loop measure are defined by \begin{equation}
\begin{aligned}
\Omega_{\Lambda}^{\mathrm{loop},m}
& =
\bigsqcup_{\substack{q \in \semigrposint \\ z \in \Lambda}}
\Omega_{\Lambda;z,z}^{q,m},
\\
\sum_{\zeta \in \Omega_{\Lambda}^{\mathrm{loop},m}}
\fun{F}{\zeta}
\fun{\msrcal{L}_{\Lambda}^{(m)}}{\setone{\zeta}}
& =
\sum_{q \in \semigrposint}\frac{1}{q}
\sum_{z \in \Lambda}
\sum_{\zeta \in \Omega_{\Lambda;z,z}^{q,m}}
\fun{F}{\zeta}
\fun{\msrcal{W}_{\Lambda;z,z}^{q\beta,(m)}}{\setone{\zeta}}
\end{aligned}
\label{eq:fock-discrete-ideal-loop-measure}
\end{equation} for every nonnegative function \(F\) on \(\Omega_{\Lambda}^{\mathrm{loop},m}\). For finite tuples \(\boldsymbol{\omega}\) and \(\boldsymbol{\zeta}\) whose entries belong to the discrete bridge path spaces in \eqref{eq:fock-discrete-bridge-path-space}, the notation \(\fun{\Phi_{\Lambda}^{(m)}}{\boldsymbol{\omega},\boldsymbol{\zeta}}\) means that the occupation numbers of all paths in the concatenated tuple are inserted into the definition of \(\Phi_{\Lambda}^{(m)}\). The discrete open-path and loop sums are \begin{equation}
\begin{aligned}
\fun{Z_{\Lambda}^{(m)}}{\boldsymbol{\omega}}
& =
\sum_{\ell \in \monnat}\frac{1}{\ell!}
\sum_{\boldsymbol{\zeta} \in \rbk{\Omega_{\Lambda}^{\mathrm{loop},m}}^{\ell}}
\napiernum^{-\fun{\Phi_{\Lambda}^{(m)}}{\boldsymbol{\omega},\boldsymbol{\zeta}}}
\prod_{j = 1}^{\ell}
\fun{\msrcal{L}_{\Lambda}^{(m)}}{\setone{\zeta_j}},
\\
Z_{\Lambda}^{(m)}
& =
\sum_{\ell \in \monnat}\frac{1}{\ell!}
\sum_{\boldsymbol{\zeta} \in \rbk{\Omega_{\Lambda}^{\mathrm{loop},m}}^{\ell}}
\napiernum^{-\fun{\Phi_{\Lambda}^{(m)}}{\boldsymbol{\zeta}}}
\prod_{j = 1}^{\ell}
\fun{\msrcal{L}_{\Lambda}^{(m)}}{\setone{\zeta_j}}.
\end{aligned}
\label{eq:fock-discrete-open-path-loop-partition-functions}
\end{equation} Let \(\widetilde{\rho}_{\Lambda,k}^{(m)}\) denote the reduced kernel obtained from \eqref{eq:fock-reduced-density-kernel} by replacing the Gibbs density matrix with the normalized \(m\)-slice Trotter product and its trace with \(Z_{\Lambda}^{(m)}\). For \(M\in\monnat\), \(\boldsymbol{z}\in\Lambda^M\), and \(\widehat{\varpi} \in \grsym{k+M}\), its action on the terminal tuple \(\rbk{\boldsymbol{y},\boldsymbol{z}}\) is defined componentwise by \begin{equation}
\begin{aligned}
\rbk{\fun{\widehat{\varpi}}{\boldsymbol{y},\boldsymbol{z}}}_a
& =
\begin{cases}
y_{\fun{\widehat{\varpi}}{a}}, & 1 \leq \fun{\widehat{\varpi}}{a} \leq k,\\
z_{\fun{\widehat{\varpi}}{a}-k}, & k+1 \leq \fun{\widehat{\varpi}}{a} \leq k+M,
\end{cases}
\quad
1 \leq a \leq k+M.
\end{aligned}
\label{eq:fock-permuted-terminal-tuple}
\end{equation}

\begin{lem}[Creation--annihilation sector formula]\label{lem:fock-creation-annihilation-sector-formula}
Assume \eqref{eq:fock-path-chemical-potential-gap}, let $\Lambda\Subset\Gamma$ be finite, and let $\beta>0$.
For $k,m\in\semigrposint$ and $\boldsymbol{x},\boldsymbol{y}\in\Lambda^k$, the finite-slice reduced kernel is
\begin{equation}
\begin{aligned}
\fun{\widetilde{\rho}_{\Lambda,k}^{(m)}}{\boldsymbol{x};\boldsymbol{y}}
& =
\frac{1}{Z_{\Lambda}^{(m)}}
\sum_{M \in \monnat}\frac{1}{M!}
\sum_{\boldsymbol{z} \in \Lambda^M}
\sum_{\widehat{\varpi} \in \grsym{k+M}}
\sum_{\substack{
\boldsymbol{u}^{(1)},\ldots,\boldsymbol{u}^{(m-1)} \in \Lambda^{k+M}\\
\boldsymbol{u}^{(0)} = \rbk{\boldsymbol{x},\boldsymbol{z}}\\
\boldsymbol{u}^{(m)} = \fun{\widehat{\varpi}}{\boldsymbol{y},\boldsymbol{z}}}}
\fun{B_{k+M,m}}{\boldsymbol{u}^{(0)},\ldots,\boldsymbol{u}^{(m)}}.
\end{aligned}
\label{eq:fock-creation-annihilation-sector-formula}
\end{equation}
\end{lem}

\begin{proof}
If $N<k$, then $a_{x_k}\cdots a_{x_1}$ vanishes on
$P_{\Lambda,N}\sphilb{F}_{\Lambda}$.
Fix $N\geq k$ and $\boldsymbol{x}=\rbk{x_1,\ldots,x_k}$, and set
$$\begin{aligned}
\fun{\kappa_u}{\boldsymbol{x}}
& =
\abscard{\set{j\in\setone{1,\ldots,k}}{x_j=u}},
\quad u\in\Lambda.
\end{aligned}$$
Choose an occupation-number basis vector $\Psi_N=\Psi_{\Lambda,\boldsymbol{\nu}}$ from
\eqref{eq:fock-occupation-number-basis} with $\sum_{u\in\Lambda}\nu_u=N$.
If $\nu_u<\fun{\kappa_u}{\boldsymbol{x}}$ for some $u\in\Lambda$, the product
$a_{x_k}\cdots a_{x_1}$ contains more than $\nu_u$ annihilation operators at $u$.
The occupation-number formula \eqref{eq:fock-occupation-number-basis} gives
$a_u^{\nu_u+1}\Psi_N=0$.
Annihilation operators at distinct sites commute, and the factor $a_u^{\fun{\kappa_u}{\boldsymbol{x}}}$
therefore gives
$$\begin{aligned}
a_{x_k}\cdots a_{x_1}\Psi_N
& =
0.
\end{aligned}$$
It remains to consider the occupation-number basis vectors satisfying
$$\begin{aligned}
\nu_u
& \geq
\fun{\kappa_u}{\boldsymbol{x}},
\quad u\in\Lambda.
\end{aligned}$$
Identify $P_{\Lambda,N}\sphilb{F}_{\Lambda}$ with the symmetric functions in
$\lpseq^2\rbk{\Lambda^N}$.
Let $\psi_N$ be the symmetric coordinate function of $\Psi_N$, and let
$\psi_{N-k}^{\boldsymbol{x}}$ be that of $a_{x_k}\cdots a_{x_1}\Psi_N$.
For $\boldsymbol{z}=\rbk{z_1,\ldots,z_{N-k}}$, the $x_j$ are the annihilated coordinates and the $z_j$ are
the remaining coordinates; the $x_j$ need not occur among the $z_j$.
Definition \eqref{eq:fock-occupation-number-basis} shows that the two values
$\fun{\psi_{N-k}^{\boldsymbol{x}}}{\boldsymbol{z}}$ and
$\fun{\psi_N}{x_1,\ldots,x_k,z_1,\ldots,z_{N-k}}$ vanish unless
$$\begin{aligned}
\fun{\kappa_u}{\boldsymbol{x}}
+
\sum_{j=1}^{N-k}\fun{\fndef{\setone{u}}}{z_j}
& =
\nu_u,
\quad u\in\Lambda.
\end{aligned}$$
For every $\boldsymbol{z}\in\Lambda^{N-k}$, the coordinate representation of the annihilation product is
$$\begin{aligned}
\fun{\psi_{N-k}^{\boldsymbol{x}}}{\boldsymbol{z}}
& =
\sqrt{\frac{N!}{\rbk{N-k}!}}
\fun{\psi_N}{x_1,\ldots,x_k,z_1,\ldots,z_{N-k}}.
\end{aligned}$$
The adjoint string $\faadj{a_{y_1}}\cdots\faadj{a_{y_k}}$ supplies the same square-root coefficient in the
matrix element defining \eqref{eq:fock-reduced-density-kernel}.
Consequently the creation-annihilation strings contribute the factor
$N!/\rbk{N-k}!$.
Multiplication by the coefficient $1/N!$ of the bosonic symmetrizer gives
$$\begin{aligned}
\frac{1}{N!}\frac{N!}{\rbk{N-k}!}
& =
\frac{1}{\rbk{N-k}!}.
\end{aligned}$$
Set $M = N-k$, and let $\boldsymbol{z} = \rbk{z_1,\ldots,z_M} \in \Lambda^M$ denote the remaining particle
coordinates.
For $\widehat{\varpi} \in \grsym{k+M}$, the initial endpoint tuple is
$\rbk{x_1,\ldots,x_k,z_1,\ldots,z_M}$, while the terminal endpoint tuple is the
$\widehat{\varpi}$-permutation of
$\rbk{y_1,\ldots,y_k,z_1,\ldots,z_M}$.
The permutation $\widehat{\varpi}$ acts on the coordinate-index set $\setone{1,\ldots,k+M}$ and permutes the
positions in these tuples.
It does not act on the lattice $\Lambda$.
Its terminal action is \eqref{eq:fock-permuted-terminal-tuple}.
Insertion of the intermediate coordinates in the weight
\eqref{eq:fock-finite-slice-sector-weight} proves
\eqref{eq:fock-creation-annihilation-sector-formula}.
\end{proof}

Separate the coordinate-index set into \[\begin{aligned}
D_k
& =
\setone{1,\ldots,k},
&
U_M
& =
\setone{k+1,\ldots,k+M}.
\end{aligned}\] The set \(D_k\) contains the distinguished coordinate indices, and \(U_M\) contains the remaining coordinate indices.

\begin{lem}[Marked-index permutation decomposition]\label{lem:fock-marked-index-permutation-decomposition}
Let $k \in \semigrposint$, $M \in \monnat$, and $\widehat{\varpi} \in \grsym{k+M}$.
For every $i \in D_k$, define
$$\begin{aligned}
q_i
& =
\min\set{q \in \semigrposint}{\fun{\widehat{\varpi}^{q}}{i} \in D_k},
\\
a_{i,r}
& =
\fun{\widehat{\varpi}^{r}}{i},
\quad 0 \leq r \leq q_i.
\end{aligned}$$
Define the map $\varpi \colon D_k \to D_k$ by
$$\begin{aligned}
\fun{\varpi}{i}
& =
a_{i,q_i},
\quad i \in D_k.
\end{aligned}$$
The map $\varpi$ is a permutation of $D_k$, and
$$\begin{aligned}
a_{i,0}
& =
i,
&
a_{i,r}
& \in
U_M
\quad
\rbk{1 \leq r \leq q_i-1},
&
a_{i,q_i}
& =
\fun{\varpi}{i}
\in
D_k.
\end{aligned}$$
The sequence $\rbk{a_{i,0},\ldots,a_{i,q_i}}$ is the marked chain from $i$ to $\fun{\varpi}{i}$.
The restriction of $\widehat{\varpi}$ to
$U_M \setminus \bigcup_{i=1}^{k}\set{a_{i,r}}{1\leq r\leq q_i-1}$ is a permutation.
There are $\ell \in \monnat$ and points
$c_{j,0} \in U_M \setminus \bigcup_{i=1}^{k}\set{a_{i,r}}{1\leq r\leq q_i-1}$,
$1 \leq j \leq \ell$, such that the sets
$$\begin{aligned}
C_j
& =
\set{\fun{\widehat{\varpi}^{r}}{c_{j,0}}}{r \in \monnat},
\quad 1 \leq j \leq \ell,
\end{aligned}$$
are pairwise disjoint and satisfy
$$\begin{aligned}
U_M \setminus \bigcup_{i=1}^{k}\set{a_{i,r}}{1\leq r\leq q_i-1}
& =
\bigsqcup_{j=1}^{\ell}C_j.
\end{aligned}$$
The sets $C_1,\ldots,C_{\ell}$ are called the unmarked cycles.
For $1 \leq j \leq \ell$ and $0 \leq r \leq \abscard{C_j}$, define
$c_{j,r}=\fun{\widehat{\varpi}^{r}}{c_{j,0}}$.
These definitions give
$$\begin{aligned}
\setone{c_{j,0},\ldots,c_{j,\abscard{C_j}-1}}
& =
C_j,
&
c_{j,\abscard{C_j}}
& =
c_{j,0},
\\
\fun{\widehat{\varpi}}{c_{j,r}}
& =
c_{j,r+1},
&
0 \leq r
& \leq
\abscard{C_j}-1.
\end{aligned}$$
The resulting partition of $U_M$ and its cardinality are
$$\begin{aligned}
U_M
=
\rbk{\bigsqcup_{i=1}^{k}\set{a_{i,r}}{1\leq r\leq q_i-1}}
\sqcup
\rbk{\bigsqcup_{j=1}^{\ell}C_j},
\quad
M
=
\sum_{i = 1}^{k}\rbk{q_i-1}
+
\sum_{j = 1}^{\ell}\abscard{C_j}.
\end{aligned}$$
For fixed $\varpi$, $q_1,\ldots,q_k$, and $C_1,\ldots,C_{\ell}$, reindexing the sum over
$\widehat{\varpi}$ in \eqref{eq:fock-creation-annihilation-sector-formula}, including its coefficient
$1/M!$, produces
\begin{equation}
\begin{aligned}
\frac{1}{\ell!}
\prod_{j = 1}^{\ell}\frac{1}{\abscard{C_j}}.
\end{aligned}
\label{eq:fock-marked-index-permutation-coefficient}
\end{equation}
\end{lem}

\begin{proof}
For $i \in D_k$, finiteness of $\setone{1,\ldots,k+M}$ and bijectivity of $\widehat{\varpi}$ give a
$q \in \semigrposint$ such that $\fun{\widehat{\varpi}^{q}}{i}=i$.
Hence the set defining $q_i$ is nonempty.
For each $a \in \setone{1,\ldots,k+M}$, consider the set
$\set{\fun{\widehat{\varpi}^{r}}{a}}{r \in \monnat}$.
This set either intersects $D_k$ or is contained in $U_M$.
The points in its intersection with $D_k$, listed in the order in which positive powers of
$\widehat{\varpi}$ reach them, divide it into the sequences
$\rbk{a_{i,0},\ldots,a_{i,q_i}}$.
On this list, $\varpi$ sends each point to the next point and sends the last point to the first point.
Thus $\varpi$ is bijective on each intersection and is a permutation of $D_k$.
Each set generated by an index in
$U_M \setminus \bigcup_{i=1}^{k}\set{a_{i,r}}{1\leq r\leq q_i-1}$ is one of the sets $C_j$.
These alternatives prove the partition of $U_M$ in the statement.
For fixed $\varpi$, $q_1,\ldots,q_k$, and $C_1,\ldots,C_{\ell}$, the $M$ indices in $U_M$ can be assigned to
the positions $a_{i,r}$, $1\leq r\leq q_i-1$, and $c_{j,r}$,
$0\leq r\leq\abscard{C_j}-1$, in $M!$ ways.
This factor cancels the coefficient $1/M!$ in
\eqref{eq:fock-creation-annihilation-sector-formula}.
The marked chains have fixed initial indices.  The unmarked cycles can be ordered in $\ell!$ ways, and the
$j$-th cycle has $\abscard{C_j}$ choices of $c_{j,0}$.  Removing these repetitions leaves
$$\begin{aligned}
\frac{1}{\ell!}
\prod_{j = 1}^{\ell}\frac{1}{\abscard{C_j}}.
\end{aligned}$$
The factor $1/\abscard{C_j}$ agrees with the length-$\abscard{C_j}$ factor in the discrete loop measure
$\msrcal{L}_{\Lambda}^{(m)}$ defined by \eqref{eq:fock-discrete-ideal-loop-measure}.
\end{proof}

\begin{lem}[Finite-slice open-path expansion]\label{lem:fock-finite-slice-open-path-expansion}
Assume \eqref{eq:fock-path-chemical-potential-gap}, let $\Lambda\Subset\Gamma$ be finite, and let $\beta>0$.
For every $m\in\semigrposint$, the finite-slice loop series is the trace normalization:
\begin{equation}
\begin{aligned}
Z_{\Lambda}^{(m)}
& =
\sum_{N\in\monnat}
\sqfun{\trace_{P_{\Lambda,N}\sphilb{F}_{\Lambda}}}{T_{\Lambda,m}}.
\end{aligned}
\label{eq:fock-finite-slice-partition-trace-identity}
\end{equation}
For every $k,m\in\semigrposint$ and $\boldsymbol{x},\boldsymbol{y}\in\Lambda^k$, the finite-slice reduced kernel
satisfies
\begin{equation}
\begin{aligned}
\fun{\widetilde{\rho}_{\Lambda,k}^{(m)}}{\boldsymbol{x};\boldsymbol{y}}
& =
\sum_{\varpi \in \grsym{k}}
\sum_{q_1,\ldots,q_k \in \semigrposint}
\sum_{\substack{
\omega_i \in \Omega_{\Lambda;x_i,y_{\fun{\varpi}{i}}}^{q_i,m}\\
1 \leq i \leq k}}
\frac{\fun{Z_{\Lambda}^{(m)}}{\boldsymbol{\omega}}}{Z_{\Lambda}^{(m)}}
\prod_{i = 1}^{k}
\fun{\msrcal{W}_{\Lambda;x_i,y_{\fun{\varpi}{i}}}^{q_i\beta,(m)}}{\setone{\omega_i}}.
\end{aligned}
\label{eq:fock-finite-slice-open-path-expansion}
\end{equation}
\end{lem}

\begin{proof}
Lemma \ref{lem:fock-creation-annihilation-sector-formula} supplies the sector formula.
Apply Lemma \ref{lem:fock-marked-index-permutation-decomposition} to each permutation in that formula.
The reindexed coefficient is \eqref{eq:fock-marked-index-permutation-coefficient}.
For each coordinate index $a$, the initial and terminal sites are
$$\begin{aligned}
u_a^{(0)}
& =
\begin{cases}
x_a, & a \in D_k,\\
z_{a-k}, & a \in U_M,
\end{cases}
&
u_a^{(m)}
& =
\begin{cases}
y_{\fun{\widehat{\varpi}}{a}}, & \fun{\widehat{\varpi}}{a} \in D_k,\\
z_{\fun{\widehat{\varpi}}{a}-k}, & \fun{\widehat{\varpi}}{a} \in U_M.
\end{cases}
\end{aligned}$$
Concatenating the coordinate sequences along the marked chains and unmarked cycles defines
$$\begin{aligned}
\omega_{i,rm+s}
& =
u_{a_{i,r}}^{(s)},
\quad
0 \leq r \leq q_i-1,
\quad
0 \leq s \leq m,
\\
\zeta_{j,rm+s}
& =
u_{c_{j,r}}^{(s)},
\quad
0 \leq r \leq \abscard{C_j}-1,
\quad
0 \leq s \leq m.
\end{aligned}$$
The endpoint equalities used to concatenate the coordinate sequences are
$$\begin{aligned}
u_{a_{i,r}}^{(m)}
& =
z_{a_{i,r+1}-k}
=
u_{a_{i,r+1}}^{(0)},
\quad
0 \leq r \leq q_i-2,
\\
u_{c_{j,r}}^{(m)}
& =
u_{c_{j,r+1}}^{(0)},
\quad
0 \leq r \leq \abscard{C_j}-2,
\\
u_{c_{j,\abscard{C_j}-1}}^{(m)}
& =
u_{c_{j,0}}^{(0)}.
\end{aligned}$$
The endpoints of the path $\omega_i$ are
$$\begin{aligned}
\omega_{i,0}
& =
x_i,
&
\omega_{i,q_i m}
& =
y_{\fun{\varpi}{i}}.
\end{aligned}$$
In particular, the $i$-th marked path has $q_i$ blocks of $m$ time slices, while the $j$-th unmarked path is a
loop with $\abscard{C_j}m$ time slices.

The one-step part of $B_{k+M,m}$ factorizes as
$$\begin{aligned}
&
\prod_{s = 0}^{m-1}
\prod_{a = 1}^{k+M}
\fun{b_m}{u_a^{(s)},u_a^{(s+1)}}
=
\prod_{i = 1}^{k}
\prod_{t = 0}^{q_i m-1}
\fun{b_m}{\omega_{i,t},\omega_{i,t+1}}
\prod_{j = 1}^{\ell}
\prod_{t = 0}^{\abscard{C_j}m-1}
\fun{b_m}{\zeta_{j,t},\zeta_{j,t+1}}
\\ %%%%%%%%%%%%%%%%
& =
\prod_{i = 1}^{k}
\fun{\msrcal{W}_{\Lambda;x_i,y_{\fun{\varpi}{i}}}^{q_i\beta,(m)}}{\setone{\omega_i}}
\prod_{j = 1}^{\ell}
\fun{\msrcal{W}_{\Lambda;z_{c_{j,0}-k},z_{c_{j,0}-k}}^{\abscard{C_j}\beta,(m)}}{\setone{\zeta_j}}.
\end{aligned}$$
\begin{samepage}
At the $s$-th Trotter slice, the same partition of the coordinate indices gives the occupation identity
$$\begin{aligned}
\nu_x^{(s)}
& =
\sum_{a = 1}^{k+M}
\fun{\fndef{\setone{x}}}{u_a^{(s)}}
=
\sum_{i = 1}^{k}
\sum_{r = 0}^{q_i-1}
\fun{\fndef{\setone{x}}}{\omega_{i,rm+s}}
+
\sum_{j = 1}^{\ell}
\sum_{r = 0}^{\abscard{C_j}-1}
\fun{\fndef{\setone{x}}}{\zeta_{j,rm+s}}
\\
& =
\fun{n_{\boldsymbol{\omega}}^{(m)}}{s,x}
+
\fun{n_{\boldsymbol{\zeta}}^{(m)}}{s,x}.
\end{aligned}$$
\end{samepage}
Substitution of this identity into the occupation-energy factor of $B_{k+M,m}$ gives
$$\begin{aligned}
\prod_{s = 0}^{m-1}
\napiernum^{-\frac{\sminvtemperature}{m}\fun{E_{\Lambda}}{\nu^{(s)}}}
& =
\napiernum^{-\fun{\Phi_{\Lambda}^{(m)}}{\boldsymbol{\omega},\boldsymbol{\zeta}}}.
\end{aligned}$$
Combining the one-step and occupation-energy identities gives the complete factorization
$$\begin{aligned}
\fun{B_{k+M,m}}{\boldsymbol{u}^{(0)},\ldots,\boldsymbol{u}^{(m)}}
& =
\napiernum^{-\fun{\Phi_{\Lambda}^{(m)}}{\boldsymbol{\omega},\boldsymbol{\zeta}}}
\prod_{i = 1}^{k}
\fun{\msrcal{W}_{\Lambda;x_i,y_{\fun{\varpi}{i}}}^{q_i\beta,(m)}}{\setone{\omega_i}}
\prod_{j = 1}^{\ell}
\fun{\msrcal{W}_{\Lambda;z_{c_{j,0}-k},z_{c_{j,0}-k}}^{\abscard{C_j}\beta,(m)}}{\setone{\zeta_j}}.
\end{aligned}$$
For $k=0$, the sets generated by repeated application of the coordinate-index permutation are precisely the
unmarked cycles $C_1,\ldots,C_{\ell}$.
Ordering these sets and choosing one of the $\abscard{C_j}$ indices as $c_{j,0}$ produces the factor
$\ell!^{-1}\prod_{j=1}^{\ell}\abscard{C_j}^{-1}$.
This coefficient and the factorization of $B_{M,m}$ convert the sum of the sector traces in
\eqref{eq:fock-finite-slice-trace-formula} into the loop series
\eqref{eq:fock-discrete-open-path-loop-partition-functions}.
This proves \eqref{eq:fock-finite-slice-partition-trace-identity}.
The coefficient
$\frac{1}{\ell!}\prod_{j = 1}^{\ell}\frac{1}{\abscard{C_j}}$ and the path factorization rewrite the sector sum as
$$\begin{aligned}
\fun{\widetilde{\rho}_{\Lambda,k}^{(m)}}{\boldsymbol{x};\boldsymbol{y}}
& =
\frac{1}{Z_{\Lambda}^{(m)}}
\sum_{\varpi \in \grsym{k}}
\sum_{q_1,\ldots,q_k \in \semigrposint}
\sum_{\substack{
\omega_i \in \Omega_{\Lambda;x_i,y_{\fun{\varpi}{i}}}^{q_i,m}\\
1 \leq i \leq k}}
\prod_{i = 1}^{k}
\fun{\msrcal{W}_{\Lambda;x_i,y_{\fun{\varpi}{i}}}^{q_i\beta,(m)}}{\setone{\omega_i}}
\\
& \mathrel{\phantom{=}}
\times
\sum_{\ell \in \monnat}\frac{1}{\ell!}
\sum_{\boldsymbol{\zeta}\in\rbk{\Omega_{\Lambda}^{\mathrm{loop},m}}^{\ell}}
\napiernum^{-\fun{\Phi_{\Lambda}^{(m)}}{\boldsymbol{\omega},\boldsymbol{\zeta}}}
\prod_{j = 1}^{\ell}
\fun{\msrcal{L}_{\Lambda}^{(m)}}{\setone{\zeta_j}}.
\end{aligned}$$
The final two sums and the interaction factor are precisely
$\fun{Z_{\Lambda}^{(m)}}{\boldsymbol{\omega}}$ by
\eqref{eq:fock-discrete-open-path-loop-partition-functions}.
Substitution of this definition proves \eqref{eq:fock-finite-slice-open-path-expansion}.
The derivation identifies the marked endpoint condition
$\omega_{i,0} = x_i$ and $\omega_{i,q_i m} = y_{\fun{\varpi}{i}}$.
The unmarked cycles remain in the normalized loop factor.
\end{proof}

\begin{lem}[Finite-slice domination]\label{lem:fock-finite-slice-domination}
Under the assumptions of Lemma \ref{lem:fock-finite-slice-open-path-expansion}, suppose that $v$ is pointwise
nonnegative.
Every tuple $\boldsymbol{\omega}$ occurring in \eqref{eq:fock-finite-slice-open-path-expansion} satisfies
\begin{equation}
\begin{aligned}
0
\leq
\frac{\fun{Z_{\Lambda}^{(m)}}{\boldsymbol{\omega}}}{Z_{\Lambda}^{(m)}}
\leq
\napiernum^{-\fun{\Phi_{\Lambda}^{(m)}}{\boldsymbol{\omega}}}
\leq
1.
\end{aligned}
\label{eq:fock-finite-slice-normalized-loop-bound}
\end{equation}
The total loop mass, the finite-slice partition function, and the reduced kernel obey
\begin{equation}
\begin{aligned}
\sum_{\zeta\in\Omega_{\Lambda}^{\mathrm{loop},m}}
\fun{\msrcal{L}_{\Lambda}^{(m)}}{\setone{\zeta}}
& \leq
-\abscard{\Lambda}
\log\rbk{1-\napiernum^{-\beta\varepsilon}},
\\
Z_{\Lambda}^{(m)}
& \leq
\frac{1}{\rbk{1-\napiernum^{-\beta\varepsilon}}^{\abscard{\Lambda}}},
\\
0
\leq
\fun{\widetilde{\rho}_{\Lambda,k}^{(m)}}{\boldsymbol{x};\boldsymbol{y}}
& \leq
\frac{k!}{\rbk{\napiernum^{\beta\varepsilon}-1}^{k}}.
\end{aligned}
\label{eq:fock-finite-slice-uniform-bounds}
\end{equation}
\end{lem}

\begin{proof}
For the concatenated tuple $\rbk{\boldsymbol{\omega},\boldsymbol{\zeta}}$, the discrete occupation numbers
satisfy
$$\begin{aligned}
\fun{n_{\rbk{\boldsymbol{\omega},\boldsymbol{\zeta}}}^{(m)}}{r,x}
& =
\fun{n_{\boldsymbol{\omega}}^{(m)}}{r,x}
+
\fun{n_{\boldsymbol{\zeta}}^{(m)}}{r,x}.
\end{aligned}$$
Substitution into the quadratic and linear terms of $\Phi_{\Lambda}^{(m)}$ gives
$$\begin{aligned}
&
\fun{n_{\rbk{\boldsymbol{\omega},\boldsymbol{\zeta}}}^{(m)}}{r,x}
\fun{n_{\rbk{\boldsymbol{\omega},\boldsymbol{\zeta}}}^{(m)}}{r,y}
\\ %%%%%%%%%%%%%%%%
& =
\fun{n_{\boldsymbol{\omega}}^{(m)}}{r,x}
\fun{n_{\boldsymbol{\omega}}^{(m)}}{r,y}
+
\fun{n_{\boldsymbol{\zeta}}^{(m)}}{r,x}
\fun{n_{\boldsymbol{\zeta}}^{(m)}}{r,y}
+
\fun{n_{\boldsymbol{\omega}}^{(m)}}{r,x}
\fun{n_{\boldsymbol{\zeta}}^{(m)}}{r,y}
+
\fun{n_{\boldsymbol{\zeta}}^{(m)}}{r,x}
\fun{n_{\boldsymbol{\omega}}^{(m)}}{r,y}.
\end{aligned}$$
The symmetry and pointwise nonnegativity of $v$ identify the two mixed quadratic terms and give
$$\begin{aligned}
\fun{I_{\Lambda}^{(m)}}{\boldsymbol{\omega},\boldsymbol{\zeta}}
& =
\frac{2\sminvtemperature}{m}
\sum_{r = 0}^{m - 1}
\sum_{x,y \in \Lambda}
\fun{v}{x,y}
\fun{n_{\boldsymbol{\omega}}^{(m)}}{r,x}
\fun{n_{\boldsymbol{\zeta}}^{(m)}}{r,y}
\geq
0,
\\
\fun{\Phi_{\Lambda}^{(m)}}{\boldsymbol{\omega},\boldsymbol{\zeta}}
& =
\fun{\Phi_{\Lambda}^{(m)}}{\boldsymbol{\omega}}
+
\fun{\Phi_{\Lambda}^{(m)}}{\boldsymbol{\zeta}}
+
\fun{I_{\Lambda}^{(m)}}{\boldsymbol{\omega},\boldsymbol{\zeta}}
\geq
\fun{\Phi_{\Lambda}^{(m)}}{\boldsymbol{\omega}}
+
\fun{\Phi_{\Lambda}^{(m)}}{\boldsymbol{\zeta}}.
\end{aligned}$$
For every tuple $\boldsymbol{\xi}$, the discrete self-energy is nonnegative because
$$\begin{aligned}
&
\sum_{x,y \in \Lambda}
\fun{v}{x,y}
\rbk{
\fun{n_{\boldsymbol{\xi}}^{(m)}}{r,x}
\fun{n_{\boldsymbol{\xi}}^{(m)}}{r,y}
-
\fun{\fndef{\setone{x}}}{y}
\fun{n_{\boldsymbol{\xi}}^{(m)}}{r,x}}
\\
& =
\sum_{x \in \Lambda}
\fun{v}{x,x}
\fun{n_{\boldsymbol{\xi}}^{(m)}}{r,x}
\rbk{\fun{n_{\boldsymbol{\xi}}^{(m)}}{r,x}-1}
+
\sum_{\substack{x,y \in \Lambda\\x\neq y}}
\fun{v}{x,y}
\fun{n_{\boldsymbol{\xi}}^{(m)}}{r,x}
\fun{n_{\boldsymbol{\xi}}^{(m)}}{r,y}
\geq
0.
\end{aligned}$$
The discrete loop series in
\eqref{eq:fock-discrete-open-path-loop-partition-functions} satisfies
$$\begin{aligned}
&
\fun{Z_{\Lambda}^{(m)}}{\boldsymbol{\omega}}
=
\sum_{\ell \in \monnat}\frac{1}{\ell!}
\sum_{\boldsymbol{\zeta}\in\rbk{\Omega_{\Lambda}^{\mathrm{loop},m}}^{\ell}}
\napiernum^{-\fun{\Phi_{\Lambda}^{(m)}}{\boldsymbol{\omega},\boldsymbol{\zeta}}}
\prod_{j = 1}^{\ell}
\fun{\msrcal{L}_{\Lambda}^{(m)}}{\setone{\zeta_j}}
\\
& \leq
\napiernum^{-\fun{\Phi_{\Lambda}^{(m)}}{\boldsymbol{\omega}}}
\sum_{\ell \in \monnat}\frac{1}{\ell!}
\sum_{\boldsymbol{\zeta}\in\rbk{\Omega_{\Lambda}^{\mathrm{loop},m}}^{\ell}}
\napiernum^{-\fun{\Phi_{\Lambda}^{(m)}}{\boldsymbol{\zeta}}}
\prod_{j = 1}^{\ell}
\fun{\msrcal{L}_{\Lambda}^{(m)}}{\setone{\zeta_j}}
=
\napiernum^{-\fun{\Phi_{\Lambda}^{(m)}}{\boldsymbol{\omega}}}
Z_{\Lambda}^{(m)}
\leq
Z_{\Lambda}^{(m)}.
\end{aligned}$$
This proves \eqref{eq:fock-finite-slice-normalized-loop-bound}.

The total mass of the discrete ideal loop measure can be computed from
\eqref{eq:fock-discrete-ideal-loop-measure}.
The semigroup property and the row estimate \eqref{eq:fock-killed-walk-row-bound} give
$$\begin{aligned}
&
\sum_{\zeta\in\Omega_{\Lambda}^{\mathrm{loop},m}}
\fun{\msrcal{L}_{\Lambda}^{(m)}}{\setone{\zeta}}
=
\sum_{q \in \semigrposint}\frac{1}{q}
\sum_{z \in \Lambda}
\sum_{\zeta\in\Omega_{\Lambda;z,z}^{q,m}}
\prod_{r = 0}^{qm-1}
\fun{b_m}{\zeta_r,\zeta_{r+1}}
=
\sum_{q \in \semigrposint}\frac{1}{q}
\sum_{z \in \Lambda}
\fun{b_m^{qm}}{z,z}
\\
& =
\sum_{q \in \semigrposint}\frac{1}{q}
\sum_{z \in \Lambda}
\fun{p_{\Lambda,q\beta}}{z,z}
\leq
\sum_{q \in \semigrposint}\frac{1}{q}
\sum_{z,w \in \Lambda}
\fun{p_{\Lambda,q\beta}}{z,w}
\leq
\abscard{\Lambda}
\sum_{q \in \semigrposint}
\frac{\napiernum^{-q\beta\varepsilon}}{q}
=
-\abscard{\Lambda}
\log\rbk{1-\napiernum^{-\beta\varepsilon}}.
\end{aligned}$$
The equality $b_m^{qm}=p_{\Lambda,q\beta}$ follows directly from the definitions of $b_m$ and
$p_{\Lambda,t}$.
Dropping the nonnegative interaction energy and summing the exponential series gives
$$\begin{aligned}
&Z_{\Lambda}^{(m)}
\leq
\sum_{\ell \in \monnat}\frac{1}{\ell!}
\rbk{
\sum_{\zeta\in\Omega_{\Lambda}^{\mathrm{loop},m}}
\fun{\msrcal{L}_{\Lambda}^{(m)}}{\setone{\zeta}}}^{\ell}
=
\fnexp{
\sum_{\zeta\in\Omega_{\Lambda}^{\mathrm{loop},m}}
\fun{\msrcal{L}_{\Lambda}^{(m)}}{\setone{\zeta}}}
\\ %%%%%%%%%%%%%%%%
& \leq
\fnexp{
-\abscard{\Lambda}
\fun{\log}{1 - \napiernum^{-\sminvtemperature\omega \varepsilon}}}
=
\rbk{1-\napiernum^{-\beta\varepsilon}}^{-\abscard{\Lambda}}.
\end{aligned}$$

We apply the normalized-loop estimate to the finite-slice reduced kernel.
Positivity of the bridge weights gives
$$\begin{aligned}
&
0
\leq
\fun{\widetilde{\rho}_{\Lambda,k}^{(m)}}{\boldsymbol{x};\boldsymbol{y}}
\leq
\sum_{\varpi \in \grsym{k}}
\sum_{q_1,\ldots,q_k \in \semigrposint}
\sum_{\substack{
\omega_i \in \Omega_{\Lambda;x_i,y_{\fun{\varpi}{i}}}^{q_i,m}\\
1 \leq i \leq k}}
\prod_{i = 1}^{k}
\fun{\msrcal{W}_{\Lambda;x_i,y_{\fun{\varpi}{i}}}^{q_i\beta,(m)}}{\setone{\omega_i}}
\\ %%%%%%%%%%%%%%%%
& =
\sum_{\varpi \in \grsym{k}}
\prod_{i = 1}^{k}
\rbk{
\sum_{q_i \in \semigrposint}
\sum_{\omega_i \in \Omega_{\Lambda;x_i,y_{\fun{\varpi}{i}}}^{q_i,m}}
\fun{\msrcal{W}_{\Lambda;x_i,y_{\fun{\varpi}{i}}}^{q_i\beta,(m)}}{\setone{\omega_i}}}
\\ %%%%%%%%%%%%%%%%
& =
\sum_{\varpi \in \grsym{k}}
\prod_{i = 1}^{k}
\rbk{
\sum_{q_i \in \semigrposint}
\fun{b_m^{q_i m}}{x_i,y_{\fun{\varpi}{i}}}}
=
\sum_{\varpi \in \grsym{k}}
\prod_{i = 1}^{k}
\rbk{
\sum_{q_i \in \semigrposint}
\fun{p_{\Lambda,q_i\beta}}{x_i,y_{\fun{\varpi}{i}}}}
\\
& \leq
\sum_{\varpi \in \grsym{k}}
\prod_{i = 1}^{k}
\rbk{
\sum_{q_i \in \semigrposint}
\napiernum^{-q_i\beta\varepsilon}}
=
k!
\rbk{
\sum_{q \in \semigrposint}
\napiernum^{-q\beta\varepsilon}}^{k}
=
\frac{k!}{\rbk{\napiernum^{\beta\varepsilon}-1}^{k}}.
\end{aligned}$$
The marked-chain estimate is independent of $m$, the volume, and the endpoint sites.
The loop-mass calculation, the exponential-series estimate, and the marked-chain estimate prove
\eqref{eq:fock-finite-slice-uniform-bounds}.
\end{proof}

The passage from finite-slice paths to continuous time is isolated in the next lemma.

\begin{lem}[Continuous-time limit]\label{lem:fock-continuous-time-limit}
Under the assumptions of Lemma \ref{lem:fock-finite-slice-domination}, the discrete loop partition function
converges to the loop partition function in \eqref{eq:fock-open-path-loop-partition-functions}, and
\begin{equation}
\begin{aligned}
\lim_{m\to\infty}
\fun{\widetilde{\rho}_{\Lambda,k}^{(m)}}{\boldsymbol{x};\boldsymbol{y}}
& =
\sum_{\varpi \in \grsym{k}}
\sum_{q_1,\ldots,q_k \in \semigrposint}
\int_{\prod_{i = 1}^{k}\Omega_{\Lambda;x_i,y_{\fun{\varpi}{i}}}^{q_i\beta}}
\frac{\fun{Z_{\Lambda}}{\boldsymbol{\omega}}}{Z_{\Lambda}}
\prod_{i = 1}^{k}
\opdmsr{\msrcal{W}_{\Lambda;x_i,y_{\fun{\varpi}{i}}}^{q_i\beta}}(\omega_i).
\end{aligned}
\label{eq:fock-continuous-time-open-path-limit}
\end{equation}
\end{lem}

\begin{proof}
For a path $\omega\in\Omega_{\Lambda;x,y}^{q\beta}$, let its $m$-skeleton be
$\rbk{\omega_0,\ldots,\omega_{qm}}$ with $\omega_r=\omega_{r\beta/m}$.
The joint distribution of the path values at the times $r\beta/m$, $0\leq r\leq qm$, equals the finite-slice
measure because
$$\begin{aligned}
&
\fun{\msrcal{W}_{\Lambda;x,y}^{q\beta}}{
\set{\omega}{\omega_{r\beta/m}=\omega_r,\quad 0\leq r\leq qm}}
\\
& =
\prod_{r = 0}^{qm-1}
\fun{p_{\Lambda,\beta/m}}{\omega_r,\omega_{r+1}}
=
\prod_{r = 0}^{qm-1}
\fun{b_m}{\omega_r,\omega_{r+1}}
=
\fun{\msrcal{W}_{\Lambda;x,y}^{q\beta,(m)}}{
\setone{\rbk{\omega_0,\ldots,\omega_{qm}}}}.
\end{aligned}$$
For $\msrcal{W}_{\Lambda;x,y}^{q\beta}$-almost every path, the number of jumps in
$\closedinterval{0}{q\beta}$ is finite.
The left Riemann sums of the piecewise-constant occupation functions defined in
\eqref{eq:fock-path-tuple-occupation-field} satisfy
$$\begin{aligned}
&
\lim_{m \to \infty}
\frac{\sminvtemperature}{m}
\sum_{r = 0}^{m-1}
\sum_{x,y\in\Lambda}
\fun{v}{x,y}
\rbk{
\fun{n_{\boldsymbol{\xi}}}{\frac{r\beta}{m},x}
\fun{n_{\boldsymbol{\xi}}}{\frac{r\beta}{m},y}
-
\fun{\fndef{\setone{x}}}{y}
\fun{n_{\boldsymbol{\xi}}}{\frac{r\beta}{m},x}}
\\ %%%%%%%%%%%%%%%%
& =
\int_{\closedinterval{0}{\beta}}
\sum_{x,y\in\Lambda}
\fun{v}{x,y}
\rbk{
\fun{n_{\boldsymbol{\xi}}}{s,x}
\fun{n_{\boldsymbol{\xi}}}{s,y}
-
\fun{\fndef{\setone{x}}}{y}
\fun{n_{\boldsymbol{\xi}}}{s,x}}
\opdmsr{s}
=
\fun{\Phi_{\Lambda}}{\boldsymbol{\xi}}.
\end{aligned}$$
Fix $s\in\semigrposint$, endpoints $x_a,y_a\in\Lambda$, and lengths
$q_a\in\semigrposint$ for $1\leq a\leq s$.
Since $0 \leq \napiernum^{-\Phi_{\Lambda}^{(m)}} \leq 1$,
dominated convergence for this fixed path tuple gives
\begin{equation}
\begin{aligned}
&
\lim_{m \to \infty}
\sum_{\substack{
\xi_a^{(m)}\in\Omega_{\Lambda;x_a,y_a}^{q_a,m}\\
1\leq a\leq s}}
\napiernum^{-\fun{\Phi_{\Lambda}^{(m)}}{\boldsymbol{\xi}^{(m)}}}
\prod_{a=1}^{s}
\fun{\msrcal{W}_{\Lambda;x_a,y_a}^{q_a\beta,(m)}}{
\setone{\xi_a^{(m)}}}
\\ %%%%%%%%%%%%%%%%
& =
\int_{\prod_{a=1}^{s}\Omega_{\Lambda;x_a,y_a}^{q_a\beta}}
\napiernum^{-\fun{\Phi_{\Lambda}}{\boldsymbol{\xi}}}
\prod_{a=1}^{s}
\opdmsr{\msrcal{W}_{\Lambda;x_a,y_a}^{q_a\beta}}(\xi_a).
\end{aligned}
\label{eq:fock-fixed-tuple-skeleton-limit}
\end{equation}
The discrete sum ranges over skeletons with the prescribed endpoints.

The loop-count, loop-length, and marked-length tails are uniform in $m$.
The loop-count tail satisfies
$$\begin{aligned}
0
& \leq
\sum_{\ell=L+1}^{\infty}\frac{1}{\ell!}
\sum_{\boldsymbol{\zeta}\in\rbk{\Omega_{\Lambda}^{\mathrm{loop},m}}^{\ell}}
\napiernum^{-\fun{\Phi_{\Lambda}^{(m)}}{\boldsymbol{\zeta}}}
\prod_{j=1}^{\ell}
\fun{\msrcal{L}_{\Lambda}^{(m)}}{\setone{\zeta_j}}
\leq
\sum_{\ell=L+1}^{\infty}\frac{1}{\ell!}
\rbk{
-\abscard{\Lambda}
\log\rbk{1-\napiernum^{-\beta\varepsilon}}}^{\ell}
\xrightarrow{L \to \infty}
0.
\end{aligned}$$
The loop-length tail satisfies
$$\begin{aligned}
0
& \leq
\sum_{q=Q+1}^{\infty}\frac{1}{q}
\sum_{z\in\Lambda}
\sum_{\zeta\in\Omega_{\Lambda;z,z}^{q,m}}
\fun{\msrcal{W}_{\Lambda;z,z}^{q\beta,(m)}}{\setone{\zeta}}
=
\sum_{q=Q+1}^{\infty}\frac{1}{q}
\sum_{z\in\Lambda}
\fun{p_{\Lambda,q\beta}}{z,z}
\leq
\abscard{\Lambda}
\sum_{q=Q+1}^{\infty}
\frac{\napiernum^{-q\beta\varepsilon}}{q}
\xrightarrow{Q \to \infty}
0.
\end{aligned}$$
For loop tuples, summing over the position of a loop whose length exceeds $Q\beta$ gives
$$\begin{aligned}
&
\sum_{\ell=1}^{\infty}\frac{1}{\ell!}
\sum_{\substack{
\boldsymbol{\zeta}\in\rbk{\Omega_{\Lambda}^{\mathrm{loop},m}}^{\ell}\\
\zeta_j\in\bigsqcup_{\substack{q\in\semigrposint,\ q>Q\\z\in\Lambda}}\Omega_{\Lambda;z,z}^{q,m}
\text{ for some }1\leq j\leq\ell}}
\prod_{j=1}^{\ell}
\fun{\msrcal{L}_{\Lambda}^{(m)}}{\setone{\zeta_j}}
\\
& \leq
\rbk{
\abscard{\Lambda}
\sum_{q=Q+1}^{\infty}
\frac{\napiernum^{-q\beta\varepsilon}}{q}}
\fnexp{
\abscard{\Lambda}
\sum_{q\in\semigrposint}
\frac{\napiernum^{-q\beta\varepsilon}}{q}}
\xrightarrow{Q \to \infty}
0.
\end{aligned}$$
For each marked path, the length tail satisfies
$$\begin{aligned}
0
& \leq
\sum_{q=Q+1}^{\infty}
\sum_{\omega\in\Omega_{\Lambda;x,y}^{q,m}}
\fun{\msrcal{W}_{\Lambda;x,y}^{q\beta,(m)}}{\setone{\omega}}
=
\sum_{q=Q+1}^{\infty}
\fun{p_{\Lambda,q\beta}}{x,y}
\leq
\sum_{q=Q+1}^{\infty}\napiernum^{-q\beta\varepsilon}
\xrightarrow{Q \to \infty}
0.
\end{aligned}$$
The sum over marked lengths with at least one $q_i>Q$ is bounded uniformly in the endpoints by
$$\begin{aligned}
&
\sum_{\varpi\in\grsym{k}}
\sum_{\substack{q_1,\ldots,q_k\in\semigrposint\\q_i>Q\text{ for some }1\leq i\leq k}}
\prod_{i=1}^{k}
\fun{p_{\Lambda,q_i\beta}}{x_i,y_{\fun{\varpi}{i}}}
\leq
k!k
\rbk{
\sum_{q=Q+1}^{\infty}\napiernum^{-q\beta\varepsilon}}
\rbk{
\sum_{q\in\semigrposint}\napiernum^{-q\beta\varepsilon}}^{k-1}
\xrightarrow{Q \to \infty}
0.
\end{aligned}$$
The fixed-tuple convergence, the three tail estimates, and the finiteness of $\grsym{k}$ permit the limit
$m\to\infty$ to pass through the loop-count, loop-length, and marked-length sums in
\eqref{eq:fock-finite-slice-open-path-expansion}.
Applying the same argument with no marked paths gives $Z_{\Lambda}^{(m)}\to Z_{\Lambda}$.
Applying it with the $k$ marked paths gives \eqref{eq:fock-continuous-time-open-path-limit}.
\end{proof}

\begin{lem}[Continuous bosonic cycle decomposition]\label{lem:fock-continuous-bosonic-cycle-decomposition}
Assume \eqref{eq:fock-path-chemical-potential-gap}, suppose that $v$ is pointwise nonnegative, let
$\Lambda\Subset\Gamma$ be finite, and let $\beta>0$.
The operator partition function in \eqref{eq:fock-finite-volume-gibbs-state} equals the loop partition function:
\begin{equation}
\begin{aligned}
Z_{\Lambda,\beta}
& =
Z_{\Lambda}.
\end{aligned}
\label{eq:fock-continuous-cycle-partition-identity}
\end{equation}
For every $k\in\semigrposint$ and $\boldsymbol{x},\boldsymbol{y}\in\Lambda^k$, the reduced kernel satisfies
\begin{equation}
\begin{aligned}
\fun{\widetilde{\rho}_{\Lambda,k}}{\boldsymbol{x};\boldsymbol{y}}
& =
\sum_{\varpi \in \grsym{k}}
\sum_{q_1,\ldots,q_k \in \semigrposint}
\int_{\prod_{i = 1}^{k}\Omega_{\Lambda;x_i,y_{\fun{\varpi}{i}}}^{q_i\beta}}
\frac{\fun{Z_{\Lambda}}{\boldsymbol{\omega}}}{Z_{\Lambda}}
\prod_{i = 1}^{k}
\opdmsr{\msrcal{W}_{\Lambda;x_i,y_{\fun{\varpi}{i}}}^{q_i\beta}}(\omega_i).
\end{aligned}
\label{eq:fock-continuous-cycle-reduced-kernel}
\end{equation}
\end{lem}

\begin{proof}
Apply the fixed-tuple limit \eqref{eq:fock-fixed-tuple-skeleton-limit} to the finite-slice trace formula
\eqref{eq:fock-finite-slice-trace-formula}.
The trace majorant in Proposition \ref{prop:fock-finite-volume-gibbs-existence} permits summation over the
particle sectors and gives
\begin{equation}
\begin{aligned}
Z_{\Lambda,\beta}
& =
\sum_{N \in \monnat}\frac{1}{N!}
\sum_{\varpi \in \grsym{N}}
\sum_{\boldsymbol{x} \in \Lambda^N}
\int_{\prod_{i = 1}^{N}\Omega_{\Lambda;x_i,x_{\fun{\varpi}{i}}}^{\beta}}
\napiernum^{-\fun{\Phi_{\Lambda}}{\omega_1,\ldots,\omega_N}}
\prod_{i = 1}^{N}
\opdmsr{\msrcal{W}_{\Lambda;x_i,x_{\fun{\varpi}{i}}}^{\beta}}(\omega_i).
\end{aligned}
\label{eq:fock-precycle-partition-function}
\end{equation}
Let $n_q$ be the number of cycles of length $q$ in $\varpi$.
The cycle lengths exhaust all particle-coordinate indices and satisfy
\begin{equation}
\begin{aligned}
\sum_{q \in \semigrposint}q n_q
& =
N,
\end{aligned}
\label{eq:fock-bosonic-cycle-length-constraint}
\end{equation}
The number of permutations with these cycle counts is obtained by arranging the $N$ particle-coordinate indices
in a row.
There are $N!$ arrangements.
Each of the $n_q$ cycles of length $q$ has $q$ cyclic rotations representing the same cycle, and the cycles
of the same length have $n_q!$ orderings.
Division by these repetitions gives
\begin{equation}
\begin{aligned}
\abscard{\setone{\varpi \in \grsym{N}:\varpi\text{ has $n_q$ cycles of length $q$}}}
& =
\frac{N!}{\prod_{q \in \semigrposint}q^{n_q} n_q!},
\\
\frac{1}{N!}
\frac{N!}{\prod_{q \in \semigrposint}q^{n_q} n_q!}
& =
\prod_{q \in \semigrposint}
\frac{1}{n_q!}\rbk{\frac{1}{q}}^{n_q}.
\end{aligned}
\label{eq:fock-bosonic-cycle-count}
\end{equation}

We calculate how a single cycle becomes one loop.
Fix a cycle $\rbk{i_1\,i_2\,\cdots\,i_q}$ and use the convention $i_{q+1}=i_1$.
For paths $\omega_{i_r}\in\Omega_{\Lambda;x_{i_r},x_{i_{r+1}}}^{\beta}$, define their concatenated loop
$\zeta\in\Omega_{\Lambda;x_{i_1},x_{i_1}}^{q\beta}$ by
$$\begin{aligned}
\zeta_{s+\rbk{r-1}\beta}
& =
\omega_{i_r,s},
\quad
0\leq s<\beta,
\quad
1\leq r\leq q,
\\
\zeta_{q\beta}
& =
x_{i_1}.
\end{aligned}$$
Repeated application of the bridge convolution law
\eqref{eq:fock-bridge-convolution-law} gives, for every nonnegative function $F$ of the
concatenated loop,
$$\begin{aligned}
&
\sum_{x_{i_1},\ldots,x_{i_q}\in\Lambda}
\int_{\prod_{r=1}^{q}
\Omega_{\Lambda;x_{i_r},x_{i_{r+1}}}^{\beta}}
\fun{F}{\zeta}
\prod_{r=1}^{q}
\opdmsr{\msrcal{W}_{\Lambda;x_{i_r},x_{i_{r+1}}}^{\beta}}(\omega_{i_r})
=
\sum_{x_{i_1}\in\Lambda}
\int_{\Omega_{\Lambda;x_{i_1},x_{i_1}}^{q\beta}}
\fun{F}{\zeta}
\opdmsr{\msrcal{W}_{\Lambda;x_{i_1},x_{i_1}}^{q\beta}}(\zeta).
\end{aligned}$$
The loop occupation at a time $s\in\closedinterval{0}{\beta}$ is the sum of the occupations of its
constituent paths at that time:
$$\begin{aligned}
\fun{n_{\zeta}}{s,x}
& =
\sum_{r=1}^{q}
\fun{\fndef{\setone{x}}}{\omega_{i_r,s}}.
\end{aligned}$$
If $\boldsymbol{\zeta}$ is the tuple obtained by concatenating every cycle of $\varpi$, summation of this
identity over all cycles gives
$$\begin{aligned}
\fun{n_{\boldsymbol{\zeta}}}{s,x}
=
\sum_{i=1}^{N}
\fun{\fndef{\setone{x}}}{\omega_{i,s}},
\quad
\fun{\Phi_{\Lambda}}{\boldsymbol{\zeta}}
=
\fun{\Phi_{\Lambda}}{\omega_1,\ldots,\omega_N}.
\end{aligned}$$

For each $q\in\semigrposint$, define the length-$q\beta$ part of the ideal loop measure by
$$\begin{aligned}
\int_{\Omega_{\Lambda}^{\mathrm{loop}}}
\fun{F}{\zeta}
\opdmsr{\msrcal{L}_{\Lambda,q}}(\zeta)
& =
\frac{1}{q}
\sum_{z\in\Lambda}
\int_{\Omega_{\Lambda;z,z}^{q\beta}}
\fun{F}{\zeta}
\opdmsr{\msrcal{W}_{\Lambda;z,z}^{q\beta}}(\zeta).
\end{aligned}$$
The loop-measure definition \eqref{eq:fock-ideal-loop-measure} gives
$\msrcal{L}_{\Lambda}
=
\sum_{q\in\semigrposint}
\msrcal{L}_{\Lambda,q}$.
For fixed cycle counts satisfying \eqref{eq:fock-bosonic-cycle-length-constraint}, the coefficient
\eqref{eq:fock-bosonic-cycle-count}, the bridge convolution law
\eqref{eq:fock-bridge-convolution-law}, and the energy identity compute the
corresponding contribution to \eqref{eq:fock-precycle-partition-function} as
$$\begin{aligned}
&
\prod_{q\in\semigrposint}\frac{1}{n_q!}
\int_{\prod_{q\in\semigrposint}
\rbk{\Omega_{\Lambda}^{\mathrm{loop}}}^{n_q}}
\napiernum^{-\fun{\Phi_{\Lambda}}{\boldsymbol{\zeta}}}
\prod_{q\in\semigrposint}
\opdmsr{\msrcal{L}_{\Lambda,q}^{\otimes n_q}}(\boldsymbol{\zeta}_q).
\end{aligned}$$
Only finitely many $n_q$ are nonzero because their weighted sum is $N$.
Let $\ell=\sum_{q\in\semigrposint}n_q$.
The multinomial expansion of
$\msrcal{L}_{\Lambda}^{\otimes\ell}$ contains
$\ell!/\prod_{q\in\semigrposint}n_q!$ orderings with these fixed counts.
Multiplication by $1/\ell!$ gives
$$\begin{aligned}
\frac{1}{\ell!}
\frac{\ell!}{\prod_{q\in\semigrposint}n_q!}
& =
\prod_{q\in\semigrposint}\frac{1}{n_q!}.
\end{aligned}$$
Summing over all cycle counts and then over $\ell$ yields
\begin{equation}
\begin{aligned}
Z_{\Lambda,\beta}
& =
\sum_{\ell\in\monnat}\frac{1}{\ell!}
\int_{\rbk{\Omega_{\Lambda}^{\mathrm{loop}}}^{\ell}}
\napiernum^{-\fun{\Phi_{\Lambda}}{\boldsymbol{\zeta}}}
\opdmsr{\msrcal{L}_{\Lambda}^{\otimes\ell}}(\boldsymbol{\zeta})
=
Z_{\Lambda}.
\end{aligned}
\label{eq:fock-loop-grand-partition-identity}
\end{equation}
The quantity $Z_{\Lambda,\beta}$ in \eqref{eq:fock-loop-grand-partition-identity} is the operator trace defined
in \eqref{eq:fock-finite-volume-gibbs-state}.
The quantity $Z_{\Lambda}$ is the loop series defined in
\eqref{eq:fock-open-path-loop-partition-functions}.
Equation \eqref{eq:fock-loop-grand-partition-identity} proves
\eqref{eq:fock-continuous-cycle-partition-identity}.

For the reduced kernel, the $k$ creation-annihilation pairs cut the permutation cycles that contain the marked
particle-coordinate indices.
Fix $\varpi \in \grsym{k}$, marked lengths $q_1,\ldots,q_k$, and unmarked cycle counts
$\seq{n_q}{q \in \semigrposint}$ with
$$\begin{aligned}
N
& =
\sum_{i = 1}^{k}q_i
+
\sum_{q \in \semigrposint}q n_q.
\end{aligned}$$
The factor $N!/\rbk{N-k}!$ selects and orders the $k$ distinguished particle-coordinate indices.
After the marked endpoints and $\varpi$ are fixed, the remaining $N-k$ particle-coordinate indices fill the
$q_i-1$ intermediate positions of the $i$-th marked chain, $1\leq i\leq k$, and the positions in the unmarked
cycles.
There is no cyclic-rotation quotient for a marked chain because its initial endpoint is prescribed.
The number of arrangements of the remaining particle-coordinate indices is
$$\begin{aligned}
\frac{\rbk{N-k}!}
{\prod_{q\in\semigrposint}q^{n_q}n_q!}.
\end{aligned}$$
Multiplication by the selection factor and the bosonic symmetrizer gives
$$\begin{aligned}
\frac{N!}{\rbk{N - k}!}
\frac{1}{N!}
\frac{\rbk{N - k}!}{\prod_{q \in \semigrposint}q^{n_q} n_q!}
& =
\prod_{q \in \semigrposint}
\frac{1}{n_q!}\rbk{\frac{1}{q}}^{n_q}.
\end{aligned}$$
For fixed marked paths
$\boldsymbol{\omega}=\rbk{\omega_1,\ldots,\omega_k}$, summation over all unmarked cycle counts is
$$\begin{aligned}
\sum_{\ell\in\monnat}\frac{1}{\ell!}
\int_{\rbk{\Omega_{\Lambda}^{\mathrm{loop}}}^{\ell}}
\napiernum^{-\fun{\Phi_{\Lambda}}{\boldsymbol{\omega},\boldsymbol{\zeta}}}
\opdmsr{\msrcal{L}_{\Lambda}^{\otimes\ell}}(\boldsymbol{\zeta})
=
\fun{Z_{\Lambda}}{\boldsymbol{\omega}}.
\end{aligned}$$
Substitution of this identity into the marked-chain sum gives
$$\begin{aligned}
\fun{\widetilde{\rho}_{\Lambda,k}}{\boldsymbol{x};\boldsymbol{y}}
& =
\sum_{\varpi \in \grsym{k}}
\sum_{q_1,\ldots,q_k\in\semigrposint}
\int_{\prod_{i=1}^{k}
\Omega_{\Lambda;x_i,y_{\fun{\varpi}{i}}}^{q_i\beta}}
\frac{\fun{Z_{\Lambda}}{\boldsymbol{\omega}}}{Z_{\Lambda}}
\prod_{i=1}^{k}
\opdmsr{\msrcal{W}_{\Lambda;x_i,y_{\fun{\varpi}{i}}}^{q_i\beta}}(\omega_i).
\end{aligned}$$
This identity is \eqref{eq:fock-continuous-cycle-reduced-kernel}.
The loop and marked-chain tail bounds prove absolute convergence of every rearrangement in the calculation.
\end{proof}

\begin{lem}[Open-path and loop expansion]\label{lem:fock-loop-marked-chain-expansion}
Assume \eqref{eq:fock-path-chemical-potential-gap}, suppose that $v$ is pointwise nonnegative, let
$\Lambda\Subset\Gamma$ be finite, and let $\beta>0$.
For every $k\in\semigrposint$ and $\boldsymbol{x},\boldsymbol{y}\in\Lambda^k$, the reduced kernel satisfies
\begin{equation}
\begin{aligned}
\fun{\widetilde{\rho}_{\Lambda,k}}{\boldsymbol{x};\boldsymbol{y}}
& =
\sum_{\varpi \in \grsym{k}}
\sum_{q_1,\ldots,q_k \in \semigrposint}
\int_{\prod_{i = 1}^{k}\Omega_{\Lambda;x_i,y_{\fun{\varpi}{i}}}^{q_i\beta}}
\frac{\fun{Z_{\Lambda}}{\boldsymbol{\omega}}}{Z_{\Lambda}}
\prod_{i = 1}^{k}
\opdmsr{\msrcal{W}_{\Lambda;x_i,y_{\fun{\varpi}{i}}}^{q_i\beta}}(\omega_i).
\end{aligned}
\label{eq:fock-reduced-density-open-path-expansion}
\end{equation}
\end{lem}

\begin{proof}
Lemma \ref{lem:fock-finite-slice-trace-formula} and
Lemma \ref{lem:fock-creation-annihilation-sector-formula} provide the coordinate formulas used in the
finite-slice expansion.
Lemma \ref{lem:fock-finite-slice-open-path-expansion} gives the exact finite-slice formula.
Its trace normalization is \eqref{eq:fock-finite-slice-partition-trace-identity}.
Lemma \ref{lem:fock-finite-slice-domination} supplies bounds that are uniform in the slice number.
Lemma \ref{lem:fock-continuous-time-limit} identifies the limit of the finite-slice reduced kernels with the
right-hand side of \eqref{eq:fock-reduced-density-open-path-expansion}.
The Trotter product \eqref{eq:fock-finite-slice-trotter-product} and the operator
$K_{\Lambda,\smchemicalpotential_{\Lambda}}$ in \eqref{eq:fock-bose-hubbard-hamiltonian} satisfy
$$\begin{aligned}
\lim_{m\to\infty}
\fnrestr{T_{\Lambda,m}}{P_{\Lambda,N}\sphilb{F}_{\Lambda}}
& =
\fnrestr{
\napiernum^{-\beta K_{\Lambda,\smchemicalpotential_{\Lambda}}}
}{P_{\Lambda,N}\sphilb{F}_{\Lambda}}
\end{aligned}$$
in the operator norm for every $N\in\monnat$.
The trace majorant in Proposition \ref{prop:fock-finite-volume-gibbs-existence} permits the particle-sector sum
and the limit to be interchanged.
The resulting limit is
$$\begin{aligned}
\lim_{m\to\infty}
\fun{\widetilde{\rho}_{\Lambda,k}^{(m)}}{\boldsymbol{x};\boldsymbol{y}}
& =
\fun{\widetilde{\rho}_{\Lambda,k}}{\boldsymbol{x};\boldsymbol{y}}.
\end{aligned}$$
Comparison with \eqref{eq:fock-continuous-time-open-path-limit} proves
\eqref{eq:fock-reduced-density-open-path-expansion}.
Lemma \ref{lem:fock-continuous-bosonic-cycle-decomposition} gives the equality of the operator and loop
normalizations in \eqref{eq:fock-continuous-cycle-partition-identity} and independently verifies the reduced
kernel formula in \eqref{eq:fock-continuous-cycle-reduced-kernel}.
\end{proof}

Let \(\oastate[\psi_{\txtfr,\Lambda}]\) be the Gibbs state at inverse temperature \(\beta\) for the free grand-canonical Hamiltonian \(\physham_{\txtfr,\Lambda}-\smchemicalpotential_{\Lambda}\opfocknumber_{\Lambda}\). It has the same hopping and chemical potential as the interacting Gibbs state \(\oastate[\psi_{\Lambda}]\) in \eqref{eq:fock-finite-volume-gibbs-state}. Let \(\widetilde{\rho}_{\txtfr,\Lambda,k}\) be its reduced density kernel defined by \eqref{eq:fock-reduced-density-kernel}. Equation \eqref{eq:fock-path-interaction-energy} is identically zero when \(v=0\). Let \(Z_{\txtfr,\Lambda}\) and \(\fun{Z_{\txtfr,\Lambda}}{\boldsymbol{\omega}}\) be the two loop partition functions of this free gas, obtained from \eqref{eq:fock-open-path-loop-partition-functions} by setting \(v=0\). In particular, \begin{equation}
\begin{aligned}
\fun{\widetilde{\rho}_{\txtfr,\Lambda,k}}{\boldsymbol{x};\boldsymbol{y}}
& =
\fun{\oastate[\psi_{\txtfr,\Lambda}]}{
\faadj{a_{y_1}}\cdots\faadj{a_{y_k}}a_{x_k}\cdots a_{x_1}}.
\end{aligned}
\label{eq:fock-free-reduced-density-kernel}
\end{equation}

\begin{lem}[Repulsive loop domination]\label{lem:fock-repulsive-loop-domination}
Assume \eqref{eq:fock-path-chemical-potential-gap}, suppose that $v$ is pointwise nonnegative, let
$\Lambda\Subset\Gamma$ be finite, and let $\beta>0$.
For every finite open-path tuple $\boldsymbol{\omega}$ in the sense of
Definition \ref{def:fock-finite-open-path-tuple}, the normalized loop factor satisfies
\begin{equation}
\begin{aligned}
\frac{\fun{Z_{\Lambda}}{\boldsymbol{\omega}}}{Z_{\Lambda}}
& \leq
\napiernum^{-\fun{\Phi_{\Lambda}}{\boldsymbol{\omega}}}
\leq
1.
\end{aligned}
\label{eq:fock-interacting-loop-cancellation}
\end{equation}
For every $k\in\semigrposint$ and $\boldsymbol{x},\boldsymbol{y}\in\Lambda^k$, the interacting reduced density
kernel in \eqref{eq:fock-reduced-density-kernel} is bounded by the free kernel in
\eqref{eq:fock-free-reduced-density-kernel} with the same hopping and chemical potential:
\begin{equation}
\begin{aligned}
0
\leq
\fun{\widetilde{\rho}_{\Lambda,k}}{\boldsymbol{x};\boldsymbol{y}}
& \leq
\fun{\widetilde{\rho}_{\txtfr,\Lambda,k}}{\boldsymbol{x};\boldsymbol{y}}.
\end{aligned}
\label{eq:fock-reduced-density-domination}
\end{equation}
\end{lem}

\begin{proof}
Use the open-path representation in Lemma \ref{lem:fock-loop-marked-chain-expansion}.
For two finite open-path tuples in the sense of Definition \ref{def:fock-finite-open-path-tuple}, define their
nonnegative cross energy by
$$\begin{aligned}
\fun{I_{\Lambda}}{\boldsymbol{\omega},\boldsymbol{\zeta}}
& =
2\int_{\closedinterval{0}{\beta}}
\sum_{x,y \in \Lambda}
\fun{v}{x,y}
\fun{n_{\boldsymbol{\omega}}}{s,x}
\fun{n_{\boldsymbol{\zeta}}}{s,y}
\opdmsr{s}.
\end{aligned}$$
Expanding the concatenated-tuple definition
\eqref{eq:fock-concatenated-path-interaction-energy} with the occupation-field definition
\eqref{eq:fock-path-tuple-occupation-field} and the interaction energy
\eqref{eq:fock-path-interaction-energy} starts from
$$\begin{aligned}
\fun{n_{\rbk{\boldsymbol{\omega},\boldsymbol{\zeta}}}}{s,x}
& =
\fun{n_{\boldsymbol{\omega}}}{s,x}
+
\fun{n_{\boldsymbol{\zeta}}}{s,x},
\\
\fun{n_{\rbk{\boldsymbol{\omega},\boldsymbol{\zeta}}}}{s,x}
\fun{n_{\rbk{\boldsymbol{\omega},\boldsymbol{\zeta}}}}{s,y}
& =
\fun{n_{\boldsymbol{\omega}}}{s,x}
\fun{n_{\boldsymbol{\omega}}}{s,y}
+
\fun{n_{\boldsymbol{\zeta}}}{s,x}
\fun{n_{\boldsymbol{\zeta}}}{s,y}
\\
& +
\fun{n_{\boldsymbol{\omega}}}{s,x}
\fun{n_{\boldsymbol{\zeta}}}{s,y}
+
\fun{n_{\boldsymbol{\zeta}}}{s,x}
\fun{n_{\boldsymbol{\omega}}}{s,y}.
\end{aligned}$$
The linear normal-ordering term splits without a mixed contribution.
Symmetry of $v$ makes the two mixed quadratic terms equal.
Integration over $s$ gives
\begin{equation}
\begin{aligned}
\fun{\Phi_{\Lambda}}{\boldsymbol{\omega},\boldsymbol{\zeta}}
=
\fun{\Phi_{\Lambda}}{\boldsymbol{\omega}}
+
\fun{\Phi_{\Lambda}}{\boldsymbol{\zeta}}
+
\fun{I_{\Lambda}}{\boldsymbol{\omega},\boldsymbol{\zeta}}
\geq
\fun{\Phi_{\Lambda}}{\boldsymbol{\omega}}
+
\fun{\Phi_{\Lambda}}{\boldsymbol{\zeta}}.
\end{aligned}
\label{eq:fock-repulsive-path-splitting}
\end{equation}
The inequality uses pointwise nonnegativity of $v$.
Insert \eqref{eq:fock-repulsive-path-splitting} into the first member of
\eqref{eq:fock-open-path-loop-partition-functions}.
The normalized background contribution is bounded by
$$\begin{aligned}
\frac{\fun{Z_{\Lambda}}{\boldsymbol{\omega}}}{Z_{\Lambda}}
\leq
\napiernum^{-\fun{\Phi_{\Lambda}}{\boldsymbol{\omega}}}
\frac{1}{Z_{\Lambda}}
\sum_{\ell \in \monnat}\frac{1}{\ell!}
\int_{\rbk{\Omega_{\Lambda}^{\mathrm{loop}}}^{\ell}}
\napiernum^{-\fun{\Phi_{\Lambda}}{\boldsymbol{\zeta}}}
\opdmsr{\msrcal{L}_{\Lambda}^{\otimes\ell}}(\boldsymbol{\zeta})
=
\napiernum^{-\fun{\Phi_{\Lambda}}{\boldsymbol{\omega}}}.
\end{aligned}$$
The final loop series is exactly $Z_{\Lambda}$.
All bridge measures in \eqref{eq:fock-reduced-density-open-path-expansion} are positive.
The bound \eqref{eq:fock-interacting-loop-cancellation} gives
$$\begin{aligned}
0
\leq
\fun{\widetilde{\rho}_{\Lambda,k}}{\boldsymbol{x};\boldsymbol{y}}
& \leq
\sum_{\varpi \in \grsym{k}}
\sum_{q_1,\ldots,q_k\in\semigrposint}
\int_{\prod_{i=1}^{k}
\Omega_{\Lambda;x_i,y_{\fun{\varpi}{i}}}^{q_i\beta}}
\prod_{i=1}^{k}
\opdmsr{\msrcal{W}_{\Lambda;x_i,y_{\fun{\varpi}{i}}}^{q_i\beta}}(\omega_i).
\end{aligned}$$
For this free gas, $v=0$ implies
$$\begin{aligned}
\fun{Z_{\txtfr,\Lambda}}{\boldsymbol{\omega}}
=
Z_{\txtfr,\Lambda},
\quad
\frac{\fun{Z_{\txtfr,\Lambda}}{\boldsymbol{\omega}}}
{Z_{\txtfr,\Lambda}}
=
1.
\end{aligned}$$
The right-hand side of the reduced-kernel bound is consequently
$\fun{\widetilde{\rho}_{\txtfr,\Lambda,k}}{\boldsymbol{x};\boldsymbol{y}}$, which proves
\eqref{eq:fock-reduced-density-domination}.
\end{proof}

\subsection{Factorial and exponential local moments}\label{factorial-and-exponential-local-moments}

The free Bose gas in the comparison has one-particle covariance \[\begin{aligned}
C_{\txtfr,\Lambda}
& =
\rbk{\napiernum^{\beta\rbk{\physham[h]_{\Lambda} - \smchemicalpotential_{\Lambda}}} - 1}^{-1}.
\end{aligned}\] The gap \eqref{eq:fock-path-chemical-potential-gap} implies \(\physham[h]_{\Lambda} - \smchemicalpotential_{\Lambda} \geq \varepsilon\) and \[\begin{aligned}
\bkt{\delta_x}{C_{\txtfr,\Lambda}\delta_x}
& \leq
\rbk{\napiernum^{\beta\varepsilon} - 1}^{-1}.
\end{aligned}\]

\begin{thm}[Exponential replacement of the DLL moment hypothesis]\label{thm:fock-all-temperature-local-number}
Assume \eqref{eq:fock-path-chemical-potential-gap} and pointwise nonnegativity of $v$.
For every $\beta > 0$, $k \in \monnat$, and $0 < a < \beta\varepsilon$, the finite-volume Gibbs states defined
by \eqref{eq:fock-finite-volume-gibbs-state} satisfy
\begin{equation}
\begin{aligned}
\sup_{\Lambda \Subset \Gamma}
\sup_{x \in \Lambda}
\fun{\oastate[\psi_{\Lambda}]}{
\prod_{j = 0}^{k - 1}\rbk{N_x - j}}
& \leq
\frac{k!}{\rbk{\napiernum^{\beta\varepsilon} - 1}^{k}},
\end{aligned}
\label{eq:fock-factorial-local-number-bound}
\end{equation}
where the product is $1$ for $k = 0$, and
\begin{equation}
\begin{aligned}
\sup_{\Lambda \Subset \Gamma}
\sup_{x \in \Lambda}
\fun{\oastate[\psi_{\Lambda}]}{\napiernum^{aN_x}}
& \leq
\frac{\napiernum^{\beta\varepsilon} - 1}
{\napiernum^{\beta\varepsilon} - \napiernum^{a}}.
\end{aligned}
\label{eq:fock-exponential-local-number-bound}
\end{equation}
\end{thm}

\begin{proof}
For $k \in \semigrposint$, evaluate \eqref{eq:fock-reduced-density-domination} at
$\boldsymbol{x} = \boldsymbol{y} = \rbk{x,\ldots,x}$.
The domination is Lemma \ref{lem:fock-repulsive-loop-domination}.
The occupation identity and the free Bose Wick rule are
$$\begin{aligned}
\rbk{\faadj{a_x}}^{k}a_x^{k}
=
\prod_{j = 0}^{k - 1}\rbk{N_x - j},
\quad
\fun{\oastate[\psi_{\txtfr,\Lambda}]}{
\rbk{\faadj{a_x}}^{k}a_x^{k}}
=
k! \rbk{\bkt{\delta_x}{C_{\txtfr,\Lambda}\delta_x}}^{k}.
\end{aligned}$$
Indeed, fix the basis vector $\Psi_{\Lambda,\boldsymbol{\nu}}$ in
\eqref{eq:fock-occupation-number-basis} with $\nu_x=n$.
For $n\geq k$, let $\boldsymbol{\nu}^{(k)}$ be the multi-index obtained by replacing $\nu_x$ with $n-k$.
Repeated annihilation and creation give
$$\begin{aligned}
a_x^k\Psi_{\Lambda,\boldsymbol{\nu}}
& =
\sqrt{n\rbk{n-1}\cdots\rbk{n-k+1}}
\Psi_{\Lambda,\boldsymbol{\nu}^{(k)}},
\\
\rbk{\faadj{a_x}}^ka_x^k\Psi_{\Lambda,\boldsymbol{\nu}}
& =
n\rbk{n-1}\cdots\rbk{n-k+1}
\Psi_{\Lambda,\boldsymbol{\nu}},
\end{aligned}$$
while $a_x^k\Psi_{\Lambda,\boldsymbol{\nu}}=0$ when $n<k$.
The two diagonal operators consequently agree.
The free Bose Wick rule sums over all bijections between the $k$ creation and $k$ annihilation entries.
Every contraction equals $\bkt{\delta_x}{C_{\txtfr,\Lambda}\delta_x}$, and the number of bijections is
$\abscard{\grsym{k}} = k!$.
The domination and covariance estimates now give
$$\begin{aligned}
&
\fun{\oastate[\psi_{\Lambda}]}{
\prod_{j=0}^{k-1}\rbk{N_x-j}}
=
\fun{\widetilde{\rho}_{\Lambda,k}}{
\rbk{x,\ldots,x};\rbk{x,\ldots,x}}
\\
& \leq
\fun{\widetilde{\rho}_{\txtfr,\Lambda,k}}{
\rbk{x,\ldots,x};\rbk{x,\ldots,x}}
=
k!\rbk{\bkt{\delta_x}{C_{\txtfr,\Lambda}\delta_x}}^k
\leq
\frac{k!}{\rbk{\napiernum^{\beta\varepsilon}-1}^k}.
\end{aligned}$$
Taking the two suprema proves \eqref{eq:fock-factorial-local-number-bound}.
The case $k = 0$ is the normalization of the state.

For every $n \in \monnat$, the binomial theorem gives the factorial generating identity
\begin{equation}
\begin{aligned}
&
\napiernum^{a n}
=
\rbk{1+\napiernum^{a}-1}^{n}
=
\sum_{k=0}^{n}
\binom{n}{k}\rbk{\napiernum^{a}-1}^{k}
\\
& =
\sum_{k=0}^{n}
\frac{\rbk{\napiernum^{a}-1}^{k}}{k!}
\frac{n!}{\rbk{n-k}!}
=
\sum_{k \in \monnat}
\frac{\rbk{\napiernum^{a} - 1}^{k}}{k!}
\prod_{j = 0}^{k - 1}\rbk{n - j},
\end{aligned}
\label{eq:fock-factorial-generating-identity}
\end{equation}
where the final infinite sum equals the finite sum because the falling factorial vanishes for $k>n$.
Apply \eqref{eq:fock-factorial-generating-identity} to the spectral resolution of $N_x$.
Every summand is positive.
Monotone convergence and \eqref{eq:fock-factorial-local-number-bound} give
$$\begin{aligned}
\fun{\oastate[\psi_{\Lambda}]}{\napiernum^{aN_x}}
=
\sum_{k \in \monnat}
\frac{\rbk{\napiernum^{a}-1}^{k}}{k!}
\fun{\oastate[\psi_{\Lambda}]}{
\prod_{j=0}^{k-1}\rbk{N_x-j}}
\leq
\sum_{k \in \monnat}
\rbk{
\frac{\napiernum^{a}-1}
{\napiernum^{\beta\varepsilon}-1}}^{k}
=
\frac{1}
{1-\frac{\napiernum^{a}-1}{\napiernum^{\beta\varepsilon}-1}}
=
\frac{\napiernum^{\beta\varepsilon} - 1}
{\napiernum^{\beta\varepsilon} - \napiernum^{a}}.
\end{aligned}$$
The geometric series converges because $a < \beta\varepsilon$.
This proves \eqref{eq:fock-exponential-local-number-bound}.
\end{proof}

Theorem \ref{thm:fock-all-temperature-local-number} is uniform in the volume and the lattice site. It is an all-temperature statement but not an all-density statement. For \(\Gamma = \ringratint^{d}\) and nearest-neighbor hopping, one may take \(C_h = 2d\), and the hypothesis is a uniform version of \(\smchemicalpotential < -2d\) in \cite[Remark 5.3]{DeuchertLampartLemm001}.

\begin{defn}[Exponential spatial weights]\label{def:fock-exponential-spatial-weights}
Fix $b>0$.
For $x,y\in\Gamma$ and $\Delta\Subset\Gamma$ or $\Delta=\Gamma$, define
\begin{equation}
\begin{aligned}
\fun{q_x}{y}
& =
\napiernum^{-b\fun{\topdist}{x,y}},
&
Q_{x,\Delta}
& =
\sum_{y\in\Delta}\fun{q_x}{y}N_y,
&
S_b
& =
\sup_{x\in\Gamma}\sum_{y\in\Gamma}\fun{q_x}{y}.
\end{aligned}
\label{eq:fock-exponential-spatial-weights}
\end{equation}
\end{defn}

Let \(P_{\Delta,n}\) be the particle-sector projection defined by \eqref{eq:fock-particle-sector-projection} with \(X=\Delta\). The inequalities \(0<\fun{q_x}{y}\leq1\) give the sector bound \begin{equation}
\begin{aligned}
0
\leq
Q_{x,\Delta}P_{\Delta,n}
\leq
nP_{\Delta,n},
\quad
n\in\monnat.
\end{aligned}
\label{eq:fock-weighted-number-sector-bound}
\end{equation}

\begin{lem}[Initial weighted exponential moment]
For every $b>0$, the constant $S_b$ in \eqref{eq:fock-exponential-spatial-weights} is finite.
Assume \eqref{eq:fock-path-chemical-potential-gap} and pointwise nonnegativity of $v$.
Fix $\beta>0$ and $a\in\openinterval{0}{\beta\varepsilon}$, and define
\begin{equation}
\begin{aligned}
M_a
& =
\frac{\napiernum^{\beta\varepsilon}-1}
{\napiernum^{\beta\varepsilon}-\napiernum^a}.
\end{aligned}
\label{eq:fock-weighted-moment-constant}
\end{equation}
For every finite $\Lambda\Subset\Gamma$, $\Delta\in\setone{\Lambda,\Gamma}$, and $x\in\Gamma$,
the state $\oastate[\overline{\psi}_{\Lambda}]$ defined by
\eqref{eq:fock-vacuum-extended-gibbs-state} satisfies
\begin{equation}
\begin{aligned}
\fun{\oastate[\overline{\psi}_{\Lambda}]}{
\napiernum^{\frac{a}{S_b}Q_{x,\Delta}}}
& \leq
M_a.
\end{aligned}
\label{eq:fock-initial-weighted-exponential-moment}
\end{equation}
\end{lem}

\begin{proof}
For every $x\in\Gamma$, decomposition into the sets of points at fixed graph distance from $x$ and
\eqref{eq:fock-polynomial-ball-volume-growth} give
$$\begin{aligned}
\sum_{y\in\Gamma}\napiernum^{-b\fun{\topdist}{x,y}}
& =
1+
\sum_{n=1}^{\infty}
\napiernum^{-bn}
\abscard{\set{y\in\Gamma}{\fun{\topdist}{x,y}=n}}
\\
& \leq
1+
\sum_{n=1}^{\infty}
\napiernum^{-bn}\abscard{x[n]}
\leq
1+\sigma
\sum_{n=1}^{\infty}n^d\napiernum^{-bn}
<
\infty.
\end{aligned}$$
Taking the supremum over $x$ proves $S_b<\infty$.
The condition $0<a<\beta\varepsilon$ gives
$$\begin{aligned}
M_a-1
& =
\frac{\napiernum^a-1}
{\napiernum^{\beta\varepsilon}-\napiernum^a}
>
0.
\end{aligned}$$
The estimate \eqref{eq:fock-exponential-local-number-bound} states that
$\fun{\oastate[\overline{\psi}_{\Lambda}]}{\napiernum^{aN_y}}\leq M_a$ for $y\in\Lambda$.
For $y\in \Delta\setminus\Lambda$, definition
\eqref{eq:fock-vacuum-extended-gibbs-state} gives
$\fun{\oastate[\overline{\psi}_{\Lambda}]}{\napiernum^{aN_y}}=1\leq M_a$.
Consequently,
$$\begin{aligned}
\fun{\oastate[\overline{\psi}_{\Lambda}]}{\napiernum^{aN_y}}
& \leq
M_a,
\quad
y\in \Delta.
\end{aligned}$$
For $\boldsymbol{n}=\rbk{n_y}_{y\in \Delta}\in\monnat^{\Delta}$, define the joint probability mass by
$$\begin{aligned}
\fun{\nu_{\Lambda,\Delta}}{\boldsymbol{n}}
& =
\fun{\oastate[\overline{\psi}_{\Lambda}]}{
\prod_{y\in \Delta}\fun{\fndef{\setone{n_y}}}{N_y}}.
\end{aligned}$$
Here $\fun{\fndef{\setone{n_y}}}{N_y}$ is the spectral projection of $N_y$ for the eigenvalue $n_y$.
The commutativity of the number operators shows that the product is their joint spectral projection,
and normalization of the state gives
$\sum_{\boldsymbol{n}\in\monnat^{\Delta}}\fun{\nu_{\Lambda,\Delta}}{\boldsymbol{n}}=1$.
Since the definition of $S_b$ gives
$$\begin{aligned}
0
\leq
\sum_{y\in \Delta}\frac{\fun{q_x}{y}}{S_b}
\leq
1,
\end{aligned}$$
generalized Hölder on the probability space
$\rbk{\monnat^{\Delta},\nu_{\Lambda,\Delta}}$ uses the exponents
$S_b/\fun{q_x}{y}$, $y\in \Delta$.
If $\sum_{y\in \Delta}\fun{q_x}{y}/S_b<1$, include the constant function $1$ with exponent
$\rbk{1-\sum_{y\in \Delta}\fun{q_x}{y}/S_b}^{-1}$; if the sum is $1$, omit this factor.
Thus
$$\begin{aligned}
&
\fun{\oastate[\overline{\psi}_{\Lambda}]}{
\napiernum^{\frac{a}{S_b}Q_{x,\Delta}}}
=
\sum_{\boldsymbol{n}\in\monnat^{\Delta}}
\fun{\nu_{\Lambda,\Delta}}{\boldsymbol{n}}
\prod_{y\in \Delta}
\rbk{\napiernum^{a n_y}}^{\frac{\fun{q_x}{y}}{S_b}}
\\ %%%%%%%%%%%%%%%%
& \leq
\rbk{
\sum_{\boldsymbol{n}\in\monnat^{\Delta}}
\fun{\nu_{\Lambda,\Delta}}{\boldsymbol{n}}}^{
1-\sum_{y\in \Delta}\fun{q_x}{y}/S_b}
\prod_{y \in \Delta}
\rbk{
\sum_{\boldsymbol{n}\in\monnat^{\Delta}}
\fun{\nu_{\Lambda,\Delta}}{\boldsymbol{n}}
\napiernum^{a n_y}}^{\frac{\fun{q_x}{y}}{S_b}}
\\ %%%%%%%%%%%%%%%%
& =
\prod_{y \in \Delta}
\rbk{
\fun{\oastate[\overline{\psi}_{\Lambda}]}{
\napiernum^{aN_y}}}^{\frac{\fun{q_x}{y}}{S_b}}
\leq
M_a^{\sum_{y\in \Delta} \frac{\fun{q_x}{y}}{S_b}}
\leq
M_a.
\end{aligned}$$
This proves \eqref{eq:fock-initial-weighted-exponential-moment}.
\end{proof}

\begin{lem}[Exponential conjugation by weighted occupation]
Let $q_x$ and $Q_{x,\Delta}$ be defined by
\eqref{eq:fock-exponential-spatial-weights}.
For $\Delta\Subset\Gamma$ or $\Delta=\Gamma$, let
$\sphilb{D}_{\Delta,\txtfin}$ be defined by
\eqref{eq:fock-finite-particle-subspace} with $X=\Delta$.
For $x\in\Gamma$, $u,v\in\Delta$, and $s\in\fldreal$, the following identities hold on
$\sphilb{D}_{\Delta,\txtfin}$:
\begin{equation}
\begin{aligned}
\napiernum^{sQ_{x,\Delta}}\faadj{a_u}\napiernum^{-sQ_{x,\Delta}}
& =
\napiernum^{s\fun{q_x}{u}}\faadj{a_u},
\\
\napiernum^{sQ_{x,\Delta}}a_v\napiernum^{-sQ_{x,\Delta}}
& =
\napiernum^{-s\fun{q_x}{v}}a_v.
\end{aligned}
\label{eq:fock-exponential-creation-annihilation-conjugation}
\end{equation}
\end{lem}

\begin{proof}
The local-number commutators in
\eqref{eq:fock-creation-annihilation-number-commutators} give
$$\begin{aligned}
\commutator{Q_{x,\Delta}}{\faadj{a_u}}
& =
\sum_{y\in\Delta}
\fun{q_x}{y}\commutator{N_y}{\faadj{a_u}}
=
\fun{q_x}{u}\faadj{a_u},
\\
\commutator{Q_{x,\Delta}}{a_v}
& =
\sum_{y\in\Delta}
\fun{q_x}{y}\commutator{N_y}{a_v}
=
-\fun{q_x}{v}a_v.
\end{aligned}$$
Choose the pair $\rbk{A,\lambda}$ from the following two possibilities:
$$\begin{aligned}
\rbk{A,\lambda}
& \in
\setone{
\rbk{\faadj{a_u},\fun{q_x}{u}},
\rbk{a_v,-\fun{q_x}{v}}}.
\end{aligned}$$
The two commutator identities give $\commutator{Q_{x,\Delta}}{A}=\lambda A$.
The sector bound \eqref{eq:fock-weighted-number-sector-bound} shows that
$Q_{x,\Delta}$ is bounded on each particle-number summand and preserves
$\sphilb{D}_{\Delta,\txtfin}$.
For $n\in\monnat$, the first choice of $A$ maps the $n$-particle summand into the
$\rbk{n+1}$-particle summand.
For $n\in\semigrposint$, the second choice maps the $n$-particle summand into the
$\rbk{n-1}$-particle summand, and it vanishes on the vacuum summand.
Thus all operators in the following derivative act on $\sphilb{D}_{\Delta,\txtfin}$.
For $\Xi\in\sphilb{D}_{\Delta,\txtfin}$, the common calculation for both choices of
$\rbk{A,\lambda}$ is
$$\begin{aligned}
\opod{s}
\rbk{
\napiernum^{-s\lambda}
\napiernum^{sQ_{x,\Delta}}
A
\napiernum^{-sQ_{x,\Delta}}\Xi}
& =
\napiernum^{-s\lambda}
\napiernum^{sQ_{x,\Delta}}
\rbk{
-\lambda A+\commutator{Q_{x,\Delta}}{A}}
\napiernum^{-sQ_{x,\Delta}}\Xi
=
0.
\end{aligned}$$
At $s=0$, the vector inside the derivative equals $A\Xi$.
Hence the vector is independent of $s$.
Multiplication by $\napiernum^{s\lambda}$ gives
$$\begin{aligned}
\napiernum^{sQ_{x,\Delta}}
A
\napiernum^{-sQ_{x,\Delta}}\Xi
& =
\napiernum^{s\lambda}A\Xi.
\end{aligned}$$
The two choices of $\rbk{A,\lambda}$ prove
\eqref{eq:fock-exponential-creation-annihilation-conjugation}.
\end{proof}

\begin{lem}[Weighted hopping commutator]
Let $C_h$ be defined by \eqref{eq:fock-uniform-degree-bound}, and let
$q_x$, $Q_{x,\Delta}$, and $S_b$ be defined by
\eqref{eq:fock-exponential-spatial-weights}.
Fix $a,b>0$, and define
\begin{equation}
\begin{aligned}
c_{a,b}
& =
C_h\rbk{\napiernum^b-1}
\fnexp{\frac{a\rbk{\napiernum^b-1}}{2S_b}}.
\end{aligned}
\label{eq:fock-weighted-commutator-constant}
\end{equation}
For $\Delta\Subset\Gamma$ or $\Delta=\Gamma$, let
$\sphilb{D}_{\Delta,\txtfin}$ be the finite-particle subspace defined by
\eqref{eq:fock-finite-particle-subspace} with $X=\Delta$.
For $x\in\Gamma$, $0<\theta\leq a/S_b$, and
$\Xi\in\sphilb{D}_{\Delta,\txtfin}$, one has
\begin{equation}
\begin{aligned}
\mathord{\pm}\imunit
\bkt{\Xi}{
\commutator{H_{\Delta}}{\napiernum^{\theta Q_{x,\Delta}}}\Xi}
& \leq
c_{a,b}\theta
\bkt{\Xi}{Q_{x,\Delta}\napiernum^{\theta Q_{x,\Delta}}\Xi}.
\end{aligned}
\label{eq:fock-weighted-exponential-commutator}
\end{equation}
\end{lem}

\begin{proof}
The density interaction commutes with $Q_{x,\Delta}$.
Fix $u,v\in\Delta$ with $u\sim v$.
Equation \eqref{eq:fock-exponential-creation-annihilation-conjugation} and insertion of the identity
$\napiernum^{-\theta Q_{x,\Delta}}\napiernum^{\theta Q_{x,\Delta}}$
between $\faadj{a_u}$ and $a_v$ give the following conjugation formula on
$\sphilb{D}_{\Delta,\txtfin}$:
\begin{equation}
\begin{aligned}
\napiernum^{\theta Q_{x,\Delta}}
\faadj{a_u}a_v
\napiernum^{-\theta Q_{x,\Delta}}
=
\rbk{\napiernum^{\theta Q_{x,\Delta}}
\faadj{a_u}\napiernum^{-\theta Q_{x,\Delta}}}
\rbk{\napiernum^{\theta Q_{x,\Delta}}
a_v\napiernum^{-\theta Q_{x,\Delta}}}
=
\napiernum^{\theta\rbk{\fun{q_x}{u} - \fun{q_x}{v}}}\faadj{a_u}a_v.
\end{aligned}
\label{eq:fock-hopping-weight-conjugation}
\end{equation}
The triangle inequality for the graph distance gives
$\abs{\fun{\topdist}{x,u}-\fun{\topdist}{x,v}}\leq1$.
If $\fun{\topdist}{x,u}\geq\fun{\topdist}{x,v}$, then
$$\begin{aligned}
\abs{\fun{q_x}{u} - \fun{q_x}{v}}
& =
\fun{q_x}{u}
\rbk{
\napiernum^{b\rbk{
\fun{\topdist}{x,u}-\fun{\topdist}{x,v}}}
-1}
\leq
\rbk{\napiernum^b-1}\fun{q_x}{u}.
\end{aligned}$$
In this case $\fun{q_x}{u}\leq\fun{q_x}{v}$.
Interchanging $u$ and $v$ in the other case proves
\begin{equation}
\begin{aligned}
\abs{\fun{q_x}{u} - \fun{q_x}{v}}
& \leq
\rbk{\napiernum^{b} - 1}
\min\setone{\fun{q_x}{u},\fun{q_x}{v}},
\quad
u\sim v.
\end{aligned}
\label{eq:fock-q-edge-difference-bound}
\end{equation}
Multiplying \eqref{eq:fock-hopping-weight-conjugation} on the right by $\napiernum^{\theta Q_{x,\Delta}}$ and subtracting
the resulting equality from $\faadj{a_u}a_v\napiernum^{\theta Q_{x,\Delta}}$ gives
\begin{equation}
\begin{aligned}
\commutator{\faadj{a_u}a_v}{\napiernum^{\theta Q_{x,\Delta}}}
& =
\rbk{
1-
\napiernum^{\theta\rbk{\fun{q_x}{u}-\fun{q_x}{v}}}}
\faadj{a_u}a_v
\napiernum^{\theta Q_{x,\Delta}}.
\end{aligned}
\label{eq:fock-weighted-hopping-commutator}
\end{equation}
Fix $\Xi\in\sphilb{D}_{\Delta,\txtfin}$ and set
$\Phi=\fnexp{\theta Q_{x,\Delta}/2}\Xi$.
The sector bound \eqref{eq:fock-weighted-number-sector-bound} shows that $Q_{x,\Delta}$ preserves every
particle-number summand and is bounded on each such summand.
Hence $\Phi\in\sphilb{D}_{\Delta,\txtfin}$, so every inner product in the following calculation is defined.
For $s\in\fldreal$,
$$\begin{aligned}
\abs{1-\napiernum^s}\napiernum^{-s/2}
& =
2\abs{\sinh\rbk{s/2}}
\leq
\abs{s}\napiernum^{\abs{s}/2}.
\end{aligned}$$
The definition of $\Phi$ gives $\Xi=\fnexp{-\theta Q_{x,\Delta}/2}\Phi$.
Substitute this equality into \eqref{eq:fock-weighted-hopping-commutator} and apply the scalar inequality with
$s
=
\theta\rbk{\fun{q_x}{u}-\fun{q_x}{v}}$.
This gives
$$\begin{aligned}
&
\abs{
\bkt{\Xi}{
\commutator{\faadj{a_u}a_v}{\napiernum^{\theta Q_{x,\Delta}}}
\Xi}}
=
\abs{
1-\napiernum^{\theta\rbk{\fun{q_x}{u}-\fun{q_x}{v}}}}
\napiernum^{-\frac{\theta}{2}
\rbk{\fun{q_x}{u}-\fun{q_x}{v}}}
\abs{\bkt{\Phi}{\faadj{a_u}a_v\Phi}}
\\
& \leq
\frac{\theta}{2}
\abs{\fun{q_x}{u}-\fun{q_x}{v}}
\napiernum^{
\frac{\theta}{2}
\abs{\fun{q_x}{u}-\fun{q_x}{v}}}
\rbk{
\bkt{\Phi}{N_u\Phi}
+
\bkt{\Phi}{N_v\Phi}},
\end{aligned}$$
where the last line uses
$$\begin{aligned}
2\abs{\bkt{\Phi}{\faadj{a_u}a_v\Phi}}
=
2\abs{\bkt{a_u\Phi}{a_v\Phi}}
\leq
\norm{a_u\Phi}^{2}+\norm{a_v\Phi}^{2}
=
\bkt{\Phi}{N_u\Phi}+\bkt{\Phi}{N_v\Phi}.
\end{aligned}$$
Equation \eqref{eq:fock-q-edge-difference-bound} and
$\min\setone{\fun{q_x}{u},\fun{q_x}{v}}\leq\fun{q_x}{u}$ and
$\min\setone{\fun{q_x}{u},\fun{q_x}{v}}\leq\fun{q_x}{v}$ give
$$\begin{aligned}
&
\abs{\fun{q_x}{u}-\fun{q_x}{v}}
\rbk{
\bkt{\Phi}{N_u\Phi}
+
\bkt{\Phi}{N_v\Phi}}
\\ %%%%%%%%%%%%%%%%
&\leq
\rbk{\napiernum^b-1}
\min\setone{\fun{q_x}{u},\fun{q_x}{v}}
\rbk{
\bkt{\Phi}{N_u\Phi}
+
\bkt{\Phi}{N_v\Phi}}
\\
& \leq
\rbk{\napiernum^{b}-1}
\rbk{
\fun{q_x}{u}\bkt{\Phi}{N_u\Phi}
+
\fun{q_x}{v}\bkt{\Phi}{N_v\Phi}}.
\end{aligned}$$
Equation \eqref{eq:fock-q-edge-difference-bound},
$\min\setone{\fun{q_x}{u},\fun{q_x}{v}}\leq1$, and $0<\theta\leq a/S_b$ also give
$$\begin{aligned}
\theta\abs{\fun{q_x}{u}-\fun{q_x}{v}}
& \leq
\frac{a\rbk{\napiernum^b-1}}{S_b}.
\end{aligned}$$
Summing the hopping terms and applying the uniform degree bound
\eqref{eq:fock-uniform-degree-bound} gives
$$\begin{aligned}
\abs{
\bkt{\Xi}{
\imunit\commutator{H_{\Delta}}{\napiernum^{\theta Q_{x,\Delta}}}
\Xi}}
\leq
c_{a,b}\theta
\sum_{y\in \Delta}
\fun{q_x}{y}
\bkt{\Phi}{N_y\Phi}
=
c_{a,b}\theta
\bkt{\Xi}{
Q_{x,\Delta}\napiernum^{\theta Q_{x,\Delta}}\Xi}.
\end{aligned}$$
The two choices of sign prove \eqref{eq:fock-weighted-exponential-commutator}.
\end{proof}

\begin{lem}[Inserted weighted exponential moment]
Assume \eqref{eq:fock-path-chemical-potential-gap} and pointwise nonnegativity of $v$.
Fix $b>0$, $\beta>0$, $a\in\openinterval{0}{\beta\varepsilon}$, a finite $\Lambda\Subset\Gamma$,
$\Delta\in\setone{\Lambda,\Gamma}$, a finite $\Lambda_0\Subset\Delta$, and $x\in\Gamma$.
Let $M_a$ be defined by \eqref{eq:fock-weighted-moment-constant}.
For every $A\in\fun{\oaresolventalgebra}{\Lambda_0}^{\gamma}$ and
$0<c\leq a/\rbk{S_b+\abscard{\Lambda_0}}$, one has
\begin{equation}
\begin{aligned}
\max\setone{
\fun{\oastate[\overline{\psi}_{\Lambda}]}{
A\napiernum^{cQ_{x,\Delta}}\faadj{A}},
\fun{\oastate[\overline{\psi}_{\Lambda}]}{
\faadj{A}\napiernum^{cQ_{x,\Delta}}A}}
& \leq
M_a\norm{A}^{2}.
\end{aligned}
\label{eq:fock-inserted-exponential-moment-bound}
\end{equation}
\end{lem}

\begin{proof}
Gauge invariance gives $\commutator{A}{N_{\Lambda_0}}=0$.
Write $Q_{x,\Delta}=Q_{x,\Lambda_0}+Q_{x,\Delta\setminus \Lambda_0}$.
For $0<c\leq a/(S_b+\abscard{\Lambda_0})$, the operator inequalities
$Q_{x,\Lambda_0}\leq N_{\Lambda_0}$ and
$$\begin{aligned}
A\napiernum^{cQ_{x,\Lambda_0}}\faadj{A}
\leq
A\napiernum^{cN_{\Lambda_0}}\faadj{A}
\leq
\norm{A}^{2}\napiernum^{cN_{\Lambda_0}}
\end{aligned}$$
hold because $A$ commutes with $N_{\Lambda_0}$.
Since $A$ also commutes with $Q_{x,\Delta\setminus \Lambda_0}$, generalized H\"older and
\eqref{eq:fock-exponential-local-number-bound} can be applied explicitly.
The operator bound is
$$\begin{aligned}
\fun{\oastate[\overline{\psi}_{\Lambda}]}{
A\napiernum^{cQ_{x,\Delta}}\faadj{A}}
& =
\fun{\oastate[\overline{\psi}_{\Lambda}]}{
A\napiernum^{cQ_{x,\Lambda_0}}\faadj{A}
\napiernum^{cQ_{x,\Delta\setminus \Lambda_0}}}
\\
& \leq
\norm{A}^{2}
\fun{\oastate[\overline{\psi}_{\Lambda}]}{
\fnexp{
c\sum_{y\in \Lambda_0}N_y
+
c\sum_{y\in \Delta\setminus \Lambda_0}\fun{q_x}{y}N_y}}.
\end{aligned}$$
The sum of the coefficients of the number operators in the last exponential is bounded by
$$\begin{aligned}
c\abscard{\Lambda_0}
+
c\sum_{y\in \Delta\setminus \Lambda_0}\fun{q_x}{y}
\leq
c\rbk{\abscard{\Lambda_0}+S_b}
\leq
a.
\end{aligned}$$
Generalized H\"older for the joint spectral measure of the commuting number operators gives
$$\begin{aligned}
&
\fun{\oastate[\overline{\psi}_{\Lambda}]}{
\fnexp{
c\sum_{y\in \Lambda_0}N_y
+
c\sum_{y\in \Delta\setminus \Lambda_0}\fun{q_x}{y}N_y}}
\leq
\prod_{y\in \Lambda_0}
\rbk{
\fun{\oastate[\overline{\psi}_{\Lambda}]}{\napiernum^{aN_y}}}^{c/a}
\prod_{y\in \Delta\setminus \Lambda_0}
\rbk{
\fun{\oastate[\overline{\psi}_{\Lambda}]}{\napiernum^{aN_y}}}
^{c\fun{q_x}{y}/a}
\\
& \leq
M_a^{
\frac{1}{a}
\rbk{c\abscard{\Lambda_0}
+
c\sum_{y\in \Delta\setminus \Lambda_0}\fun{q_x}{y}}}
\leq
M_a.
\end{aligned}$$
Combining the operator bound with the generalized Hölder estimate proves the first term in
\eqref{eq:fock-inserted-exponential-moment-bound}.
Replacing $A$ by $\faadj{A}$ proves the second term.
\end{proof}

\begin{lem}[Exponential propagation and cutoff removal]\label{lem:fock-exponential-cutoff-removal}
Assume \eqref{eq:fock-path-chemical-potential-gap}, and suppose that $v$ is pointwise nonnegative and has
finite range.
Fix $\beta>0$, $a \in \openinterval{0}{\beta\varepsilon}$, $T > 0$, and a finite $\Lambda_0 \Subset \Gamma$.
For a finite $\Lambda\Subset\Gamma$, choose $\Delta\in\setone{\Lambda,\Gamma}$ and a finite $\Delta_0$ satisfying
$\Lambda_0\subset \Delta_0\Subset \Delta$.
For $\lambda\geq1$, define the local particle-number cutoff
$$\begin{aligned}
P_{\Delta_0,\lambda}
& =
\prod_{y \in \Delta_0}\fun{\fndef{\closedinterval{0}{\lambda}}}{N_y}.
\end{aligned}$$
For $t\in\fldreal$ and $C\in\opspbddlin{\sphilb{F}_{\Delta}}$, define the corresponding cutoff dynamics by
$$\begin{aligned}
\fun{\widetilde{\alpha}_{\Delta,t}^{\Delta_0,\lambda}}{C}
& =
\napiernum^{\imunit t P_{\Delta_0,\lambda}H_{\Delta}P_{\Delta_0,\lambda}}
C
\napiernum^{-\imunit t P_{\Delta_0,\lambda}H_{\Delta}P_{\Delta_0,\lambda}}.
\end{aligned}$$
There exist constants $a_{\Lambda_0,T}>0$ and $C_{\Lambda_0,T,a}<\infty$, independent of
$\Lambda$, $\Delta$, $\Delta_0$, and $\lambda$, for which the following two conclusions hold.
\begin{enumerate}
\item
The propagated exponential moments satisfy
\begin{equation}
\begin{aligned}
\sup_{x \in \Gamma}
\sup_{\abs{t} \leq T}
\fun{\oastate[\overline{\psi}_{\Lambda}]}{
\fun{\alpha_{\Delta,t}}{\napiernum^{a_{\Lambda_0,T}N_x}}}
& \leq
C_{\Lambda_0,T,a}.
\end{aligned}
\label{eq:fock-exponential-moment-propagation}
\end{equation}
For every $A\in\fun{\oaresolventalgebra}{\Lambda_0}^{\gamma}$, the inserted propagated moments satisfy
\begin{equation}
\begin{aligned}
\sup_{x\in\Gamma}
\sup_{\abs{t}\leq T}
\max\setone{
\fun{\oastate[\overline{\psi}_{\Lambda}]}{
A\fun{\alpha_{\Delta,t}}{\napiernum^{a_{\Lambda_0,T}N_x}}\faadj{A}},
\fun{\oastate[\overline{\psi}_{\Lambda}]}{
\faadj{A}\fun{\alpha_{\Delta,t}}{\napiernum^{a_{\Lambda_0,T}N_x}}A}}
& \leq
C_{\Lambda_0,T,a}\norm{A}^{2}.
\end{aligned}
\label{eq:fock-propagated-inserted-exponential-moment-bound}
\end{equation}

\item
For every
$A \in \fun{\oaresolventalgebra}{\Lambda_0}^{\gamma}$, and every $B \in \oa{B}$, the two cutoff errors satisfy
\begin{equation}
\begin{aligned}
\sup_{\abs{t} \leq T}
\abs{
\fun{\oastate[\overline{\psi}_{\Lambda}]}{
\rbk{
\fun{\alpha_{\Delta,t}}{A}
-
\fun{\widetilde{\alpha}_{\Delta,t}^{\Delta_0,\lambda}}{P_{\Delta_0,\lambda}AP_{\Delta_0,\lambda}}}B}}
& \leq
C_{\Lambda_0,T,a}\abscard{\Delta_0}\lambda
\napiernum^{-a_{\Lambda_0,T}\lambda}\norm{A}\norm{B},
\\
\sup_{\abs{t} \leq T}
\abs{
\fun{\oastate[\overline{\psi}_{\Lambda}]}{
B\rbk{
\fun{\alpha_{\Delta,t}}{A}
-
\fun{\widetilde{\alpha}_{\Delta,t}^{\Delta_0,\lambda}}{P_{\Delta_0,\lambda}AP_{\Delta_0,\lambda}}}}}
& \leq
C_{\Lambda_0,T,a}\abscard{\Delta_0}\lambda
\napiernum^{-a_{\Lambda_0,T}\lambda}\norm{A}\norm{B}.
\end{aligned}
\label{eq:fock-exponential-cutoff-removal}
\end{equation}
\end{enumerate}
\end{lem}

\begin{proof}
Fix $b>0$ and use the weights in Definition \ref{def:fock-exponential-spatial-weights}.
Let $M_a$ and $c_{a,b}$ be defined by
\eqref{eq:fock-weighted-moment-constant} and \eqref{eq:fock-weighted-commutator-constant}, respectively.
The initial estimate is \eqref{eq:fock-initial-weighted-exponential-moment}.
For $t \geq 0$, put $\fun{\theta}{t}=aS_b^{-1}\napiernum^{-c_{a,b}t}$.
The propagated exponential moment is
$$\begin{aligned}
\fun{G_{x,\Delta}}{t}
& =
\fun{\oastate[\overline{\psi}_{\Lambda}] \circ \alpha_{\Delta,t}}
{\napiernum^{\fun{\theta}{t}Q_{x,\Delta}}}.
\end{aligned}$$
For $t\geq0$, one has $0<\fun{\theta}{t}\leq a/S_b$ and
$\fun{\theta'}{t}=-c_{a,b}\fun{\theta}{t}$.
The positive-sign estimate in \eqref{eq:fock-weighted-exponential-commutator} therefore gives the
quadratic-form inequality
$$\begin{aligned}
\fun{\theta'}{t}Q_{x,\Delta}\napiernum^{\fun{\theta}{t}Q_{x,\Delta}}
+
\imunit\commutator{H_{\Delta}}{\napiernum^{\fun{\theta}{t}Q_{x,\Delta}}}
& \leq
0.
\end{aligned}$$
Its derivative is consequently nonpositive:
$$\begin{aligned}
\fun{G_{x,\Delta}'}{t}
& =
\fun{\theta'}{t}
\fun{\oastate[\overline{\psi}_{\Lambda}] \circ \alpha_{\Delta,t}}
{Q_{x,\Delta}\napiernum^{\fun{\theta}{t}Q_{x,\Delta}}}
+
\fun{\oastate[\overline{\psi}_{\Lambda}] \circ \alpha_{\Delta,t}}
{\imunit\commutator{H_{\Delta}}
{\napiernum^{\fun{\theta}{t}Q_{x,\Delta}}}}
\leq
0.
\end{aligned}$$
Integration from $0$ to $t$ gives
$$\begin{aligned}
\fun{G_{x,\Delta}}{t}
\leq
\fun{G_{x,\Delta}}{0}
=
\fun{\oastate[\overline{\psi}_{\Lambda}]}
{\napiernum^{\frac{a}{S_b}Q_{x,\Delta}}}
\leq
M_a,
\quad
0\leq t\leq T.
\end{aligned}$$
Replacing $H_{\Delta}$ by $-H_{\Delta}$ in the commutator estimate gives the same bound for $-T\leq t\leq0$.
Define
$a_{\Lambda_0,T}=a\rbk{S_b+\abscard{\Lambda_0}}^{-1}\napiernum^{-c_{a,b}T}$.
The constant $c_{a,b}$ in \eqref{eq:fock-weighted-commutator-constant} is positive.
For $\abs{t}\leq T$, the definition of $\theta$ gives the coefficient comparison
$$\begin{aligned}
a_{\Lambda_0,T}
& =
\frac{a}{S_b+\abscard{\Lambda_0}}
\napiernum^{-c_{a,b}T}
\leq
\frac{a}{S_b}
\napiernum^{-c_{a,b}\abs{t}}
=
\fun{\theta}{\abs{t}}.
\end{aligned}$$
For $x\in\Delta$, the definition of $Q_{x,\Delta}$ in
\eqref{eq:fock-exponential-spatial-weights} and
$\fun{q_x}{x}=\napiernum^{-b\fun{\topdist}{x,x}}=1$ give
$$\begin{aligned}
Q_{x,\Delta}
& =
N_x+
\sum_{y\in\Delta\setminus\setone{x}}
\fun{q_x}{y}N_y
\geq
N_x.
\end{aligned}$$
Combining the coefficient and operator comparisons gives
$$\begin{aligned}
a_{\Lambda_0,T}N_x
& \leq
\fun{\theta}{\abs{t}}N_x
\leq
\fun{\theta}{\abs{t}}Q_{x,\Delta}.
\end{aligned}$$
The local number operators commute, so the joint spectral calculus gives
$$\begin{aligned}
\napiernum^{a_{\Lambda_0,T}N_x}
& \leq
\napiernum^{\fun{\theta}{\abs{t}}Q_{x,\Delta}}
\end{aligned}$$
for $x\in\Delta$ and $\abs{t}\leq T$.
For $x\notin\Delta$, one necessarily has $\Delta=\Lambda$ and
$x\in\Lambda^{\mathrm{c}}$.
The Hamiltonian $H_{\Lambda}$ acts only on $\sphilb{F}_{\Lambda}$, whereas $N_x$ acts on
$\sphilb{F}_{\Lambda^{\mathrm{c}}}$, and $N_x\Omega_{\Lambda^{\mathrm{c}}}=0$.
Consequently,
$$\begin{aligned}
\fun{\alpha_{\Lambda,t}}{\napiernum^{a_{\Lambda_0,T}N_x}}
& =
\napiernum^{a_{\Lambda_0,T}N_x},
&
\napiernum^{a_{\Lambda_0,T}N_x}\Omega_{\Lambda^{\mathrm{c}}}
& =
\Omega_{\Lambda^{\mathrm{c}}}.
\end{aligned}$$
Using the density matrix in \eqref{eq:fock-vacuum-extended-gibbs-state} therefore gives
$$\begin{aligned}
\fun{\oastate[\overline{\psi}_{\Lambda}]}{
\fun{\alpha_{\Lambda,t}}{\napiernum^{a_{\Lambda_0,T}N_x}}}
=
\sqfun{\trace_{\sphilb{F}_{\Lambda}}}{
\frac{\napiernum^{-\beta K_{\Lambda,\smchemicalpotential_{\Lambda}}}}
{Z_{\Lambda,\beta}}}
\bkt{\Omega_{\Lambda^{\mathrm{c}}}}{
\napiernum^{a_{\Lambda_0,T}N_x}\Omega_{\Lambda^{\mathrm{c}}}}
_{\sphilb{F}_{\Lambda^{\mathrm{c}}}}
=
1.
\end{aligned}$$
The two cases $x\in\Delta$ and $x\notin\Delta$ prove
\eqref{eq:fock-exponential-moment-propagation}.

Define
$$\begin{aligned}
\fun{\theta_{\Lambda_0}}{t}
& =
\frac{a}{S_b+\abscard{\Lambda_0}}\napiernum^{-c_{a,b}t},
\quad
t\geq0.
\end{aligned}$$
For $A\in\fun{\oaresolventalgebra}{\Lambda_0}^{\gamma}$,
Equation \eqref{eq:fock-inserted-exponential-moment-bound} gives the two initial bounds at
$c=\fun{\theta_{\Lambda_0}}{0}$.
Repeating the derivative calculation with $\theta_{\Lambda_0}$ in place of $\theta$ and using
both signs in \eqref{eq:fock-weighted-exponential-commutator} gives, for positive and negative times,
$$\begin{aligned}
\sup_{x\in\Gamma}
\sup_{\abs{t}\leq T}
\max\setone{
\fun{\oastate[\overline{\psi}_{\Lambda}]}{
A\fun{\alpha_{\Delta,t}}{\napiernum^{a_{\Lambda_0,T}N_x}}\faadj{A}},
\fun{\oastate[\overline{\psi}_{\Lambda}]}{
\faadj{A}\fun{\alpha_{\Delta,t}}{\napiernum^{a_{\Lambda_0,T}N_x}}A}}
& \leq
M_a\norm{A}^{2}.
\end{aligned}$$
After enlarging $C_{\Lambda_0,T,a}$ so that $C_{\Lambda_0,T,a}\geq M_a$,
this proves \eqref{eq:fock-propagated-inserted-exponential-moment-bound}.
The spectral inequality
$$\begin{aligned}
1-P_{\Delta_0,\lambda}
& \leq
\sum_{y \in \Delta_0}\fun{\fndef{\openinterval{\lambda}{\infty}}}{N_y}
\leq
\napiernum^{-a_{\Lambda_0,T}\lambda}
\sum_{y \in \Delta_0}\napiernum^{a_{\Lambda_0,T}N_y}
\end{aligned}$$
and \eqref{eq:fock-exponential-moment-propagation} give the propagated cutoff tail,
for $\abs{t}\leq T$,
$$\begin{aligned}
&
\fun{\oastate[\overline{\psi}_{\Lambda}]}{
\fun{\alpha_{\Delta,t}}{1-P_{\Delta_0,\lambda}}}
\leq
\napiernum^{-a_{\Lambda_0,T}\lambda}
\sum_{y\in \Delta_0}
\fun{\oastate[\overline{\psi}_{\Lambda}]}{
\fun{\alpha_{\Delta,t}}{\napiernum^{a_{\Lambda_0,T}N_y}}}
\leq
C_{\Lambda_0,T,a}\abscard{\Delta_0}
\napiernum^{-a_{\Lambda_0,T}\lambda}.
\end{aligned}$$
Let $C=P_{\Delta_0,\lambda}AP_{\Delta_0,\lambda}$.
The initial cutoff error decomposes as
$$\begin{aligned}
A-C
& =
\rbk{1-P_{\Delta_0,\lambda}}A
+
P_{\Delta_0,\lambda}A\rbk{1-P_{\Delta_0,\lambda}}.
\end{aligned}$$
Cauchy--Schwarz for the two positive functionals obtained by inserting $B$ on the right or on the left gives
\begin{equation}
\begin{aligned}
\abs{
\fun{\oastate[\overline{\psi}_{\Lambda}]}{
\fun{\alpha_{\Delta,t}}{
A-C}B}}
& \leq
C_{\Lambda_0,T,a}
\abscard{\Delta_0}
\napiernum^{-\frac{a_{\Lambda_0,T}\lambda}{2}}
\norm{A}\norm{B},
\\
\abs{
\fun{\oastate[\overline{\psi}_{\Lambda}]}{
B\fun{\alpha_{\Delta,t}}{
A-C}}}
& \leq
C_{\Lambda_0,T,a}
\abscard{\Delta_0}
\napiernum^{-\frac{a_{\Lambda_0,T}\lambda}{2}}
\norm{A}\norm{B}.
\end{aligned}
\label{eq:fock-initial-cutoff-error-bound}
\end{equation}

Let $K_{\Delta_0,\lambda}=P_{\Delta_0,\lambda}H_{\Delta}P_{\Delta_0,\lambda}$.
Differentiation of
$\fun{\alpha_{\Delta,t-s}\circ\widetilde{\alpha}_{\Delta,s}^{\Delta_0,\lambda}}{C}$ gives the Duhamel identity
\begin{equation}
\begin{aligned}
\fun{\alpha_{\Delta,t}}{C}
-
\fun{\widetilde{\alpha}_{\Delta,t}^{\Delta_0,\lambda}}{C}
& =
\imunit
\int_{0}^{t}
\fun{\alpha_{\Delta,t-s}}{
\commutator{
H_{\Delta}-K_{\Delta_0,\lambda}}{
\fun{\widetilde{\alpha}_{\Delta,s}^{\Delta_0,\lambda}}{C}}}
\opdmsr{s}
\end{aligned}
\label{eq:fock-cutoff-duhamel-identity}
\end{equation}
for $t\geq0$, with the orientation of the integral reversed for $t<0$.
Since $C=P_{\Delta_0,\lambda}CP_{\Delta_0,\lambda}$, every nonzero term in the commutator contains a factor
$1-P_{\Delta_0,\lambda}$ next to an oriented hopping operator $\faadj{a_u}a_v$ with
$u\sim v$ and $\setone{u,v}\cap \Delta_0\neq\emptyset$.
For every finite $Z\Subset \Delta$, define
$$\begin{aligned}
P_{Z,\lambda}
& =
\prod_{z\in Z}\fun{\fndef{\closedinterval{0}{\lambda}}}{N_z}.
\end{aligned}$$
For these ordered pairs, the inequalities
$$\begin{aligned}
\norm{P_{\setone{u,v},\lambda}\faadj{a_u}a_v}
& \leq
\sqrt{\lambda\rbk{\lambda+1}}
\leq
\sqrt{2}\lambda,
&
\norm{P_{\setone{u},\lambda}\faadj{a_u}a_v\rbk{N_v+1}^{-1/2}}
& \leq
\sqrt{\lambda}
\end{aligned}$$
follow directly from \eqref{eq:fock-occupation-number-basis}.
Define the spectral projection
\begin{equation}
\begin{aligned}
E_{y,\lambda}
& =
\fun{\fndef{\openinterval{\lambda}{\infty}}}{N_y}.
\end{aligned}
\label{eq:fock-local-cutoff-tail-projection}
\end{equation}
Equation \eqref{eq:fock-local-cutoff-tail-projection} and
\eqref{eq:fock-exponential-moment-propagation} give, for $u\in \Delta_0$ and $\abs{t-s}\leq T$,
$$\begin{aligned}
E_{u,\lambda}
\leq
\napiernum^{-a_{\Lambda_0,T}\lambda}\napiernum^{a_{\Lambda_0,T}N_u},
\quad
\fun{\oastate[\overline{\psi}_{\Lambda}]}{
\fun{\alpha_{\Delta,t-s}}{E_{u,\lambda}}}
\leq
C_{\Lambda_0,T,a}\napiernum^{-a_{\Lambda_0,T}\lambda}.
\end{aligned}$$
Expansion over the ordered pairs $u\sim v$ with $\setone{u,v}\cap \Delta_0\neq\emptyset$ produces a factor
$E_{u,\lambda}$ or $E_{v,\lambda}$ at an endpoint belonging to $\Delta_0$.
Fix $u,v\in \Delta_0$ with $u\sim v$.
The occupation-number basis \eqref{eq:fock-occupation-number-basis} gives
$$\begin{aligned}
\norm{\faadj{a_u}a_vP_{\Delta_0,\lambda}}
& \leq
\sqrt{\lambda\rbk{\lambda+1}}
\leq
\sqrt{2}\lambda,
\quad
\lambda\geq1.
\end{aligned}$$
For any projection $R$ and bounded operator $Z$, Cauchy--Schwarz for the state gives
$$\begin{aligned}
\abs{\fun{\oastate[\overline{\psi}_{\Lambda}]}{RZ}}^2
& \leq
\fun{\oastate[\overline{\psi}_{\Lambda}]}{R}
\fun{\oastate[\overline{\psi}_{\Lambda}]}{\faadj{Z}Z}
\leq
\fun{\oastate[\overline{\psi}_{\Lambda}]}{R}\norm{Z}^2.
\end{aligned}$$
Apply this inequality with
$$\begin{aligned}
R
=
\fun{\alpha_{\Delta,t-s}}{E_{u,\lambda}},
\quad
Z
=
\fun{\alpha_{\Delta,t-s}}{
\faadj{a_u}a_vP_{\Delta_0,\lambda}
\fun{\widetilde{\alpha}_{\Delta,s}^{\Delta_0,\lambda}}{C}}B.
\end{aligned}$$
Since $\norm{C}\leq\norm{A}$ and
$\norm{\fun{\widetilde{\alpha}_{\Delta,s}^{\Delta_0,\lambda}}{C}}=\norm{C}$, one has
$\norm{Z}
\leq
\sqrt{2}\lambda\norm{A}\norm{B}$.
Therefore,
\begin{equation}
\begin{aligned}
&
\abs{\fun{\oastate[\overline{\psi}_{\Lambda}]}{
\fun{\alpha_{\Delta,t-s}}{
E_{u,\lambda}
\faadj{a_u}a_vP_{\Delta_0,\lambda}
\fun{\widetilde{\alpha}_{\Delta,s}^{\Delta_0,\lambda}}{C}}B}}
\leq
\sqrt{2}\lambda\norm{A}\norm{B}
\rbk{\fun{\oastate[\overline{\psi}_{\Lambda}]}{
\fun{\alpha_{\Delta,t-s}}{E_{u,\lambda}}}}^{1/2}
\\
& \leq
\sqrt{2C_{\Lambda_0,T,a}}\lambda
\napiernum^{-\frac{a_{\Lambda_0,T}\lambda}{2}}
\norm{A}\norm{B}.
\end{aligned}
\label{eq:fock-right-cutoff-integrand-bound}
\end{equation}
Applying Cauchy--Schwarz to the adjoint gives
\begin{equation}
\begin{aligned}
\abs{\fun{\oastate[\overline{\psi}_{\Lambda}]}{
B\fun{\alpha_{\Delta,t-s}}{
E_{u,\lambda}
\faadj{a_u}a_vP_{\Delta_0,\lambda}
\fun{\widetilde{\alpha}_{\Delta,s}^{\Delta_0,\lambda}}{C}}}}
\leq
\sqrt{2C_{\Lambda_0,T,a}}\lambda
\napiernum^{-\frac{a_{\Lambda_0,T}\lambda}{2}}
\norm{A}\norm{B}.
\end{aligned}
\label{eq:fock-left-cutoff-integrand-bound}
\end{equation}
If Cauchy--Schwarz produces either of the two positive functionals containing $A$ or $\faadj{A}$,
use \eqref{eq:fock-propagated-inserted-exponential-moment-bound} with $A$ or $\faadj{A}$, respectively.
At most $2C_h\abscard{\Delta_0}$ ordered pairs $u\sim v$ satisfy $\setone{u,v}\cap \Delta_0\neq\emptyset$ by
\eqref{eq:fock-uniform-degree-bound}.
Applying \eqref{eq:fock-right-cutoff-integrand-bound} to each commutator term in
\eqref{eq:fock-cutoff-duhamel-identity}, summing over the ordered pairs, integrating in $s$ from $0$ to $t$,
and adding the first initial cutoff error in \eqref{eq:fock-initial-cutoff-error-bound} gives
$$\begin{aligned}
\sup_{\abs{t}\leq T}
\abs{\fun{\oastate[\overline{\psi}_{\Lambda}]}{
\rbk{
\fun{\alpha_{\Delta,t}}{A}
-
\fun{\widetilde{\alpha}_{\Delta,t}^{\Delta_0,\lambda}}{
P_{\Delta_0,\lambda}AP_{\Delta_0,\lambda}}}B}}
& \leq
C_{\Lambda_0,T,a}\abscard{\Delta_0}\lambda
\napiernum^{-\frac{a_{\Lambda_0,T}\lambda}{2}}
\norm{A}\norm{B}.
\end{aligned}$$
Applying \eqref{eq:fock-left-cutoff-integrand-bound} to each commutator term in
\eqref{eq:fock-cutoff-duhamel-identity}, summing over the ordered pairs, integrating in $s$ from $0$ to $t$,
and adding the second initial cutoff error in \eqref{eq:fock-initial-cutoff-error-bound} gives
$$\begin{aligned}
\sup_{\abs{t}\leq T}
\abs{\fun{\oastate[\overline{\psi}_{\Lambda}]}{
B\rbk{
\fun{\alpha_{\Delta,t}}{A}
-
\fun{\widetilde{\alpha}_{\Delta,t}^{\Delta_0,\lambda}}{
P_{\Delta_0,\lambda}AP_{\Delta_0,\lambda}}}}}
& \leq
C_{\Lambda_0,T,a}\abscard{\Delta_0}\lambda
\napiernum^{-\frac{a_{\Lambda_0,T}\lambda}{2}}
\norm{A}\norm{B}.
\end{aligned}$$
Replacing $a_{\Lambda_0,T}/2$ by $a_{\Lambda_0,T}$ in the notation proves
\eqref{eq:fock-exponential-cutoff-removal}.
\end{proof}

\subsection{Dynamics and equilibrium without a local-number hypothesis}\label{dynamics-and-equilibrium-without-a-local-number-hypothesis}

The model assumptions alone now imply both the finite-volume approximation and the existence of an equilibrium state satisfying the KMS boundary relation for \(\alpha_{\Gamma}\) on the concrete Buchholz algebra.

\begin{thm}[All-temperature finite-volume approximation]\label{thm:fock-all-temperature-local-approximation}
Assume that $v$ is pointwise nonnegative and has finite range, and assume
\eqref{eq:fock-path-chemical-potential-gap}.
For every $\beta > 0$, finite $\Lambda_0 \Subset \Gamma$, and $T > 0$, the finite-volume Gibbs family defined by
\eqref{eq:fock-finite-volume-gibbs-state} satisfies
\begin{equation}
\begin{aligned}
\lim_{\Lambda \nearrow \Gamma}
\sup_{\substack{A \in \fun{\oaresolventalgebra}{\Lambda_0}^{\gamma},\ \norm{A} \leq 1\\
B \in \fun{\oaresolventalgebra}{\Lambda}^{\gamma},\ \norm{B} \leq 1}}
\sup_{\abs{t} \leq T}
\abs{
\fun{\oastate[\overline{\psi}_{\Lambda}]}{
\rbk{\fun{\alpha_{\Lambda,t}}{A} - \fun{\alpha_{\Gamma,t}}{A}} B}}
& =
0.
\end{aligned}
\label{eq:fock-all-temperature-thermodynamic-limit}
\end{equation}
Here $\alpha_{\Gamma,t}$ is the pre-existing automorphism of
Fact \ref{fact:dll-buchholz-automorphism}, and the convergence is uniform on compact time intervals.
No local particle-number estimate is assumed in this statement.
\end{thm}

\begin{proof}
Fix $a \in \openinterval{0}{\beta\varepsilon}$.
For $m\in\semigrposint$, put
$$\begin{aligned}
\Delta_m
& =
\Lambda_0[(2m + 1)r],
&
P_{m,\lambda}
& =
P_{\Delta_m,\lambda},
&
\widetilde{\alpha}_{\Delta,t}^{m,\lambda}
& =
\widetilde{\alpha}_{\Delta,t}^{\Delta_m,\lambda}.
\end{aligned}$$
Equation \eqref{eq:fock-polynomial-ball-volume-growth} gives
$$\begin{aligned}
\abscard{\Delta_m}
& \leq
\sum_{x\in \Lambda_0}\abscard{x[(2m+1)r]}
\leq
\sigma\abscard{\Lambda_0}\rbk{(2m+1)r}^d
\leq
\sigma\abscard{\Lambda_0}\rbk{3r}^dm^d.
\end{aligned}$$
Set
$$\begin{aligned}
C_{\Lambda_0}
& =
\max\setone{1,2C_h}\sigma\abscard{\Lambda_0}\rbk{3r}^d.
\end{aligned}$$
Then $C_{\Lambda_0}m^d$ bounds both $\abscard{\Delta_m}$ and the number of ordered pairs $x\sim y$ having
at least one endpoint in $\Delta_m$, by \eqref{eq:fock-uniform-degree-bound}.
Lemma \ref{lem:fock-exponential-cutoff-removal} is the direct exponential substitute for polynomial-moment
propagation.
Apply its cutoff estimate \eqref{eq:fock-exponential-cutoff-removal} with $\Delta_0=\Delta_m$ on both
$\Delta=\Lambda$ and $\Delta=\Gamma$.
Since $\abscard{\Delta_m}\leq C_{\Lambda_0}m^d$, each of the two cutoff-removal errors is bounded by
\begin{equation}
\begin{aligned}
C_{\Lambda_0,T,a}m^d\lambda\napiernum^{-a_{\Lambda_0,T}\lambda}\norm{A}\norm{B}.
\end{aligned}
\label{eq:fock-direct-exponential-cutoff-error}
\end{equation}

For $\Delta\in\setone{\Lambda,\Gamma}$, equations (75)--(76) in the proof of
\cite[Lemma 4.10]{DeuchertLampartLemm001}, with $X=\Lambda_0$ and $P=P_{m,\lambda}$, become
\begin{equation}
\begin{aligned}
&
D_{\Delta,m,\lambda}
=
P_{m,\lambda}H_{\Delta}P_{m,\lambda}
-
P_{m,\lambda}H_{\Delta\setminus\Lambda_0[2 m r]}P_{m,\lambda}
-
P_{m,\lambda}H_{\Lambda_0[2 m r]}P_{m,\lambda},
\\ %%%%%%%%%%%%%%%%
&
\fun{\widetilde{\alpha}_{\Delta,t}^{m,\lambda}}{P_{m,\lambda} A P_{m,\lambda}}
-
\fun{\widetilde{\alpha}_{\Lambda_0[2 m r],t}^{m,\lambda}}{P_{m,\lambda} A P_{m,\lambda}}
\\
& =
\imunit
\int_0^t
\fun{\widetilde{\alpha}_{\Delta,t-s}^{m,\lambda}}{
\commutator{D_{\Delta,m,\lambda}}{
\fun{\widetilde{\alpha}_{\Lambda_0[2 m r],s}^{m,\lambda}}{
P_{m,\lambda} A P_{m,\lambda}}}}
\opdmsr{s}.
\end{aligned}
\label{eq:fock-shell-duhamel-expansion}
\end{equation}
Finite interaction range makes the terms of orders at most $m$ in the iterated form of
\eqref{eq:fock-shell-duhamel-expansion} vanish.
The estimates and the ordered-time integral used for its order-$m+1$ term are
$$\begin{aligned}
\norm{P_{m,\lambda}\faadj{a_x}a_yP_{m,\lambda}}
& \leq
\sqrt{2}\lambda,
\\ %%%%%%%%%%%%%%%%
\abscard{\Delta_m}
& \leq
C_{\Lambda_0}m^d,
\\ %%%%%%%%%%%%%%%%
\int_{\set{(s_1,\ldots,s_{m+1})}{
0\leq s_{m+1}\leq\cdots\leq s_1\leq\abs{t}}}
\opdmsr{s_1}\cdots\opdmsr{s_{m+1}}
& =
\frac{\abs{t}^{m+1}}{\rbk{m+1}!}.
\end{aligned}$$
Equations (84)--(85) of \cite{DeuchertLampartLemm001} supply $\kappa>0$ such that the product of the $m+1$ transported
commutator maps in \eqref{eq:fock-shell-duhamel-expansion} is bounded by
$\rbk{\kappa\lambda}^{m+1}$.
Substitution of these three estimates into its order-$m+1$ Duhamel integral gives, for $\abs{t}\leq T$,
\begin{equation}
\begin{aligned}
&
\norm{
\fun{\widetilde{\alpha}_{\Delta,t}^{m,\lambda}}{P_{m,\lambda} A P_{m,\lambda}}
-
\fun{\widetilde{\alpha}_{\Lambda_0[2 m r],t}^{m,\lambda}}{P_{m,\lambda} A P_{m,\lambda}}}
\\
& \leq
\sqrt{2}C_{\Lambda_0}\lambda m^d\norm{A}
\int_{\set{(s_1,\ldots,s_{m+1})}{
0\leq s_{m+1}\leq\cdots\leq s_1\leq\abs{t}}}
\rbk{\kappa\lambda}^{m+1}
\opdmsr{s_1}\cdots\opdmsr{s_{m+1}}
\\ %%%%%%%%%%%%%%%%
& =
\sqrt{2}C_{\Lambda_0}\lambda m^d\norm{A}
\rbk{\kappa\lambda}^{m+1}
\frac{\abs{t}^{m+1}}{\rbk{m+1}!}
\\
& \leq
\sqrt{2}C_{\Lambda_0}\lambda m^d\norm{A}
\frac{\rbk{\kappa\lambda T}^{m+1}}{\rbk{m+1}!}.
\end{aligned}
\label{eq:fock-cutoff-shell-duhamel-remainder}
\end{equation}
Taking the supremum over $\abs{t}\leq T$ in \eqref{eq:fock-cutoff-shell-duhamel-remainder} gives
\begin{equation}
\begin{aligned}
\sup_{\abs{t} \leq T}
\norm{
\fun{\widetilde{\alpha}_{\Delta,t}^{m,\lambda}}{P_{m,\lambda} A P_{m,\lambda}}
-
\fun{\widetilde{\alpha}_{\Lambda_0[2 m r],t}^{m,\lambda}}{P_{m,\lambda} A P_{m,\lambda}}}
& \leq
\sqrt{2}C_{\Lambda_0} \lambda m^d \norm{A}
\frac{\rbk{\kappa \lambda T}^{m + 1}}{\rbk{m + 1}!}.
\end{aligned}
\label{eq:fock-cutoff-shell-error}
\end{equation}
The estimate \eqref{eq:fock-cutoff-shell-error} is valid with the outer volume equal to either $\Lambda$ or
$\Gamma$.
It contains no state-dependent term.
For all sufficiently large $m$, choose
$$\begin{aligned}
\lambda
& =
\frac{m}{\napiernum^2\kappa T}
\geq
1.
\end{aligned}$$
Substitution into \eqref{eq:fock-direct-exponential-cutoff-error} gives
$$\begin{aligned}
C_{\Lambda_0,T,a}m^d\lambda\napiernum^{-a_{\Lambda_0,T}\lambda}
& =
\frac{C_{\Lambda_0,T,a}}{\napiernum^2\kappa T}
m^{d+1}
\fnexp{-\frac{a_{\Lambda_0,T}m}{\napiernum^2\kappa T}}.
\end{aligned}$$
Use the lower Stirling bound in the explicit form
$$\begin{aligned}
n!
& \geq
\rbk{\frac{n}{\napiernum}}^n,
\quad
n\in\semigrposint.
\end{aligned}$$
Apply this inequality with $n=m+1$ and substitute the selected value of $\lambda$:
$$\begin{aligned}
\rbk{m+1}!
& \geq
\rbk{\frac{m+1}{\napiernum}}^{m+1},
\\
\frac{\rbk{\kappa\lambda T}^{m+1}}{\rbk{m+1}!}
& =
\frac{\rbk{m/\napiernum^2}^{m+1}}{\rbk{m+1}!}
\leq
\rbk{\frac{m}{\napiernum\rbk{m+1}}}^{m+1}
\leq
\napiernum^{-(m+1)}.
\end{aligned}$$
Multiplication by the prefactor in \eqref{eq:fock-cutoff-shell-error} yields
$$\begin{aligned}
\sqrt{2}C_{\Lambda_0}\lambda m^d
\frac{\rbk{\kappa\lambda T}^{m+1}}{\rbk{m+1}!}
& \leq
\frac{\sqrt{2}C_{\Lambda_0}}{\napiernum^3\kappa T}
m^{d+1}\napiernum^{-m}.
\end{aligned}$$
Combining the two computations defines the joint majorant
\begin{equation}
\begin{aligned}
\fun{\delta_m}{T}
& =
C_{\Lambda_0,T,a,\beta}m^{d+1}
\rbk{
\napiernum^{-m}
+
\fnexp{-\frac{a_{\Lambda_0,T}m}{\napiernum^2\kappa T}}}.
\end{aligned}
\label{eq:fock-direct-exponential-shell-majorant}
\end{equation}
Equation \eqref{eq:fock-direct-exponential-shell-majorant} implies
$\fun{\delta_m}{T}\to0$ as $m\to\infty$.
The two cutoff-removal estimates compare the uncut $\Lambda$ and $\Gamma$ dynamics with
their cutoff counterparts, while the shell estimate compares those counterparts through $\Lambda_0[2 m r]$.
Writing $\fun{A_m}{t}=\fun{\alpha_{\Lambda_0[2mr],t}}{A}$, the resulting two-sided estimate is,
for $\Delta \in \setone{\Lambda,\Gamma}$,
\begin{equation}
\begin{aligned}
\sup_{\abs{t}\leq T}
\max\setone{
\abs{\fun{\oastate[\overline{\psi}_{\Lambda}]}{
\rbk{\fun{\alpha_{\Delta,t}}{A}-\fun{A_m}{t}}B}},
\abs{\fun{\oastate[\overline{\psi}_{\Lambda}]}{
B\rbk{\fun{\alpha_{\Delta,t}}{A}-\fun{A_m}{t}}}}}
\leq
\fun{\delta_m}{T}\norm{A}\norm{B}.
\end{aligned}
\label{eq:fock-two-sided-exponential-shell-approximation}
\end{equation}
For $\Delta=\Lambda$ and $\Delta=\Gamma$,
if $\abs{t} \leq T$,
the two estimates in
\eqref{eq:fock-two-sided-exponential-shell-approximation} give
$$\begin{aligned}
&
\abs{
\fun{\oastate[\overline{\psi}_{\Lambda}]}{
\rbk{\fun{\alpha_{\Lambda,t}}{A}-\fun{\alpha_{\Gamma,t}}{A}}B}}
\\
& \leq
\abs{
\fun{\oastate[\overline{\psi}_{\Lambda}]}{
\rbk{\fun{\alpha_{\Lambda,t}}{A}-\fun{A_m}{t}}B}}
+
\abs{
\fun{\oastate[\overline{\psi}_{\Lambda}]}{
\rbk{\fun{A_m}{t}-\fun{\alpha_{\Gamma,t}}{A}}B}}
\\
& \leq
2\fun{\delta_m}{T}\norm{A}\norm{B}.
\end{aligned}$$
The right-hand side tends to zero uniformly for $\norm{A},\norm{B}\leq1$.
Taking the indicated suprema and then $m\to\infty$ proves
\eqref{eq:fock-all-temperature-thermodynamic-limit}.
\end{proof}

\begin{cor}[All-temperature equilibrium states for $\alpha_{\Gamma}$]\label{cor:fock-all-temperature-kms-states}
Assume that $v$ is pointwise nonnegative and has finite range, and assume
\eqref{eq:fock-path-chemical-potential-gap}.
For every $\beta > 0$ and every weak-$\ast$ accumulation point $\oastate[\psi]$ of
$\seq{\oastate[\overline{\psi}_{\Lambda}]}{\Lambda \Subset \Gamma}$ defined in
\eqref{eq:fock-vacuum-extended-gibbs-state}, the state $\oastate[\psi]$ is invariant under
$\alpha_{\Gamma}$ on $\fun{\oaresolventalgebra}{\Gamma}^{\gamma}$ and satisfies the $\beta$-KMS boundary
condition \eqref{eq:fock-infinite-volume-kms-boundary}.
No local particle-number estimate is assumed, and no continuity of the automorphism group on $\oa{B}$ is asserted.
\end{cor}

\begin{proof}
Choose a subnet $\Lambda_j$ along which
$\oastate[\overline{\psi}_{\Lambda_j}] \to \oastate[\psi]$ weak-$\ast$.
Fix a finite $\Lambda_0\Subset\Gamma$ and
$A,B\in\fun{\oaresolventalgebra}{\Lambda_0}^{\gamma}$, and write
$\fun{A_m}{t} = \fun{\alpha_{\Lambda_0[2 m r],t}}{A}$.
The estimate \eqref{eq:fock-two-sided-exponential-shell-approximation} holds for
$\Delta\in\setone{\Lambda_j,\Gamma}$, with $\fun{\delta_m}{T}\to0$ independently of $j$.
For fixed $m$ and $t$, the operators
$\rbk{\fun{\alpha_{\Gamma,t}}{A} - \fun{A_m}{t}}B$ and
$B\rbk{\fun{\alpha_{\Gamma,t}}{A} - \fun{A_m}{t}}$ belong to $\oa{B}$.
Weak-$\ast$ convergence passes the estimates with $\Delta = \Gamma$ to $\oastate[\psi]$.
For every $j$ satisfying $\Lambda_0\subset\Lambda_j$, define
\begin{equation}
\begin{aligned}
\fun{F_j}{z}
& =
\frac{1}{Z_{\Lambda_j,\beta}}
\sqfun{\trace_{\sphilb{F}_{\Lambda_j}}}{
\napiernum^{-\beta K_{\Lambda_j,\smchemicalpotential_{\Lambda_j}}}
\napiernum^{\imunit zH_{\Lambda_j}}
A
\napiernum^{-\imunit zH_{\Lambda_j}}
B},
\quad
-\beta\leq\opimag z\leq0,
\\
\fun{F_j}{t}
& =
\fun{\oastate[\overline{\psi}_{\Lambda_j}]}{
\fun{\alpha_{\Lambda_j,t}}{A}B},
\\
\fun{F_j}{t-\imunit\beta}
& =
\fun{\oastate[\overline{\psi}_{\Lambda_j}]}{
B\fun{\alpha_{\Lambda_j,t}}{A}}.
\end{aligned}
\label{eq:fock-finite-volume-kms-strip-function}
\end{equation}
The trace argument in \cite[Lemma 5.4]{DeuchertLampartLemm001} proves that $F_j$ is holomorphic for
$-\beta<\opimag z<0$, continuous on the closed strip, and bounded by
$\abs{\fun{F_j}{z}}\leq\norm{A}\norm{B}$.
In the standalone derivative calculation of \cite[Lemma 5.5]{DeuchertLampartLemm001}, every number factor on the
fixed shell is bounded directly by the scalar spectral inequality
$\rbk{1+n}^{4}\leq C_c\fnexp{cn}$ and the exponential propagation
\eqref{eq:fock-exponential-moment-propagation}.
More explicitly, differentiation on the fixed finite shell gives
$$\begin{aligned}
\od{\fun{\oastate[\overline{\psi}_{\Lambda_j}]}{
\fun{A_m}{t}B}}{t}
& =
\fun{\oastate[\overline{\psi}_{\Lambda_j}]}{
\imunit\commutator{H_{\Lambda_0[2mr]}}{\fun{A_m}{t}}B},
\\
\od{\fun{\oastate[\overline{\psi}_{\Lambda_j}]}{
B\fun{A_m}{t}}}{t}
& =
\fun{\oastate[\overline{\psi}_{\Lambda_j}]}{
B\imunit\commutator{H_{\Lambda_0[2mr]}}{\fun{A_m}{t}}}.
\end{aligned}$$
The commutator expansion on the fixed shell and Cauchy--Schwarz have the form
$$\begin{aligned}
\abs{
\fun{\oastate[\overline{\psi}_{\Lambda_j}]}{
\imunit\commutator{H_{\Lambda_0[2mr]}}{\fun{A_m}{t}}B}}
\leq
C_{A,B,m}
\rbk{\fun{\oastate[\overline{\psi}_{\Lambda_j}] \circ \alpha_{\Lambda_0[2mr],t}}
{\prod_{y\in \Lambda_0[2mr]}\rbk{1+N_y}^{4}}}^{1/2}.
\end{aligned}$$
For any $c>0$, the scalar bound is obtained from
$$\begin{aligned}
C_c
& =
\sup_{n\in\monnat}\rbk{1+n}^{4}\napiernum^{-cn}
<
\infty,
\quad
\prod_{y\in \Lambda_0[2mr]}\rbk{1+N_y}^{4}
\leq
C_c^{\abscard{\Lambda_0[2mr]}}
\fnexp{c\sum_{y\in \Lambda_0[2mr]}N_y}.
\end{aligned}$$
Choose $c>0$ small enough that generalized H\"older reduces the last exponential to the one-site propagated
moments in \eqref{eq:fock-exponential-moment-propagation}.
The resulting bound is independent of $j$.
The calculation with $B$ on the left is identical.
For fixed $m$ and $T$, we obtain
$$\begin{aligned}
\sup_{\Lambda_j \supset \Lambda_0[2 m r]}\sup_{\abs{t} \leq T}
\rbk{
\abs{\od{\fun{\oastate[\overline{\psi}_{\Lambda_j}]}{\fun{A_m}{t}B}}{t}}
+
\abs{\od{\fun{\oastate[\overline{\psi}_{\Lambda_j}]}{B\fun{A_m}{t}}}{t}}}
& \leq
C_{A,B,m,T}.
\end{aligned}$$
The derivative bound makes the fixed-shell boundary functions equicontinuous.
For fixed $m$, weak-$\ast$ convergence gives pointwise convergence at every $t$.
Given $\eta>0$, choose a finite $\eta/(3C_{A,B,m,T})$-net of $\closedinterval{-T}{T}$.
Pointwise convergence at the finitely many net points and the derivative bound give
$$\begin{aligned}
\sup_{\abs{t}\leq T}
\abs{
\fun{\oastate[\overline{\psi}_{\Lambda_j}]}{\fun{A_m}{t}B}
-
\fun{\oastate[\psi]}{\fun{A_m}{t}B}}
\leq
\eta
\quad
\rbk{j\text{ sufficiently large}}.
\end{aligned}$$
The same argument applies with $B$ on the left.

The upper boundary in \eqref{eq:fock-finite-volume-kms-strip-function} has the exact decomposition
$$\begin{aligned}
\fun{F_j}{t}
-
\fun{\oastate[\psi]}{\fun{\alpha_{\Gamma,t}}{A}B}
& =
\rbk{
\fun{F_j}{t}
-
\fun{\oastate[\overline{\psi}_{\Lambda_j}]}{\fun{A_m}{t}B}}
\\
& \mathrel{\phantom{=}}
+
\rbk{
\fun{\oastate[\overline{\psi}_{\Lambda_j}]}{\fun{A_m}{t}B}
-
\fun{\oastate[\psi]}{\fun{A_m}{t}B}}
+
\rbk{
\fun{\oastate[\psi]}{\fun{A_m}{t}B}
-
\fun{\oastate[\psi]}{\fun{\alpha_{\Gamma,t}}{A}B}}.
\end{aligned}$$
The first and third terms are bounded by
$\fun{\delta_m}{T}\norm{A}\norm{B}$.
The upper boundary satisfies
\begin{equation}
\begin{aligned}
\sup_{\abs{t} \leq T}
\abs{
\fun{F_j}{t}
-
\fun{\oastate[\psi]}{\fun{\alpha_{\Gamma,t}}{A}B}}
& \leq
\sup_{\abs{t} \leq T}
\abs{
\fun{\oastate[\overline{\psi}_{\Lambda_j}]}{\fun{A_m}{t}B}
-
\fun{\oastate[\psi]}{\fun{A_m}{t}B}}
+
2 \fun{\delta_m}{T} \norm{A} \norm{B}.
\end{aligned}
\label{eq:fock-kms-upper-boundary-comparison}
\end{equation}
The lower boundary in \eqref{eq:fock-finite-volume-kms-strip-function} and the analogous decomposition with
$B$ on the left give the lower-boundary estimate.
The lower boundary satisfies
\begin{equation}
\begin{aligned}
\sup_{\abs{t} \leq T}
\abs{
\fun{F_j}{t - \imunit\beta}
-
\fun{\oastate[\psi]}{B\fun{\alpha_{\Gamma,t}}{A}}}
& \leq
\sup_{\abs{t} \leq T}
\abs{
\fun{\oastate[\overline{\psi}_{\Lambda_j}]}{B\fun{A_m}{t}}
-
\fun{\oastate[\psi]}{B\fun{A_m}{t}}}
+
2 \fun{\delta_m}{T} \norm{A} \norm{B}.
\end{aligned}
\label{eq:fock-kms-lower-boundary-comparison}
\end{equation}
For fixed $m$, the first terms on the right-hand sides of
\eqref{eq:fock-kms-upper-boundary-comparison} and
\eqref{eq:fock-kms-lower-boundary-comparison} tend to zero along the subnet.
Letting first $j$ tend along the subnet and then $m \to \infty$ proves compact-uniform convergence of both
boundaries to the two expressions containing $\alpha_{\Gamma,t}$.
The strip functions satisfy
$$\begin{aligned}
\sup_{0\leq s\leq\beta}
\sup_{t\in\fldreal}
\abs{\fun{F_j}{t-\imunit s}}
& \leq
\norm{A}\norm{B}.
\end{aligned}$$
Montel compactness gives a subsequence converging uniformly on compact subsets of the open strip.
The compact-uniform boundary limits identify the two boundary values of every such subsequential limit.
If two bounded holomorphic limits had these boundary values, their difference would have zero values on both
boundaries.
The three-lines theorem gives
$$\begin{aligned}
\abs{\fun{F_{A,B}^{(1)}-F_{A,B}^{(2)}}{t-\imunit s}}
\leq
\rbk{
\sup_{r\in\fldreal}
\abs{\fun{F_{A,B}^{(1)}-F_{A,B}^{(2)}}{r}}}^{1-\frac{s}{\beta}}
\rbk{
\sup_{r\in\fldreal}
\abs{\fun{F_{A,B}^{(1)}-F_{A,B}^{(2)}}{r-\imunit\beta}}}^{\frac{s}{\beta}}
=
0.
\end{aligned}$$
The limit is unique and defines the bounded strip function $F_{A,B}$ with the boundary values in
\eqref{eq:fock-infinite-volume-kms-boundary}.

For invariance, set $B=1$.
Finite-volume invariance and the two shell comparisons give
$$\begin{aligned}
&
\abs{
\fun{\oastate[\psi]}{\fun{\alpha_{\Gamma,t}}{A}}
-
\fun{\oastate[\psi]}{A}}
\leq
\abs{
\fun{\oastate[\psi]}{\fun{\alpha_{\Gamma,t}}{A}}
-
\fun{\oastate[\psi]}{\fun{A_m}{t}}}
+
\abs{
\fun{\oastate[\psi]}{\fun{A_m}{t}}
-
\fun{\oastate[\overline{\psi}_{\Lambda_j}]}{\fun{A_m}{t}}}
\\
&
+
\abs{
\fun{\oastate[\overline{\psi}_{\Lambda_j}]}{
\fun{A_m}{t}-\fun{\alpha_{\Lambda_j,t}}{A}}}
+
\abs{
\fun{\oastate[\overline{\psi}_{\Lambda_j}]}{
\fun{\alpha_{\Lambda_j,t}}{A}-A}}
+
\abs{
\fun{\oastate[\overline{\psi}_{\Lambda_j}]}{A}
-
\fun{\oastate[\psi]}{A}}.
\end{aligned}$$
The fourth term is zero.
For fixed $m$, the second and fifth terms tend to zero along the subnet.
The first and third terms are bounded by $\fun{\delta_m}{T}\norm{A}$.
Taking the subnet limit and then $m\to\infty$ proves invariance on local observables.
The local gauge-invariant resolvent algebras are norm dense in
$\fun{\oaresolventalgebra}{\Gamma}^{\gamma}$.
For local approximating sequences
$\seq{A_n}{n\in\semigrposint}$ and $\seq{B_n}{n\in\semigrposint}$ satisfying
$A_n\to A$ and $B_n\to B$ as $n\to\infty$, the automorphisms are isometric and the state has norm one.
The norm difference is bounded by
$$\begin{aligned}
\abs{
\fun{\oastate[\psi]}{\fun{\alpha_{\Gamma,t}}{A}B}
-
\fun{\oastate[\psi]}{\fun{\alpha_{\Gamma,t}}{A_n}B_n}}
& \leq
\norm{A-A_n}\norm{B}
+
\norm{A_n}\norm{B-B_n}.
\end{aligned}$$
The same bound holds on the lower boundary.
The uniform strip bound extends invariance and
\eqref{eq:fock-infinite-volume-kms-boundary} to the full gauge-invariant resolvent algebra.
This argument does not invoke the conditional equilibrium theorem of \cite{DeuchertLampartLemm001} and does not
transfer a particle-number functional to $\oastate[\psi]$.
\end{proof}

\section{Conclusion and Outlook}\label{conclusion-and-outlook}

Proposition 3.1 of \cite{DeuchertLampartLemm001} already supplies the generally discontinuous automorphism group on the full concrete Buchholz algebra. The obstruction in Theorem 4.1 and Theorem 5.2 of that paper is the uniform local moment condition, not existence of the algebraic automorphism. Theorem \ref{thm:fock-all-temperature-local-number} replaces that condition by a proved exponential estimate under repulsiveness and the low-activity gap. Theorem \ref{thm:fock-all-temperature-local-approximation} gives the finite-volume dynamical approximation, and Corollary \ref{cor:fock-all-temperature-kms-states} gives invariant states satisfying the KMS boundary relation at every temperature.

The discrete functional integral completes the low-activity comparison proposed in \cite[Remark 5.3]{DeuchertLampartLemm001}. The decisive formula is the interacting cancellation \eqref{eq:fock-interacting-loop-cancellation}. It implies reduced-density domination, the factorial estimate \eqref{eq:fock-factorial-local-number-bound}, and the exponential estimate \eqref{eq:fock-exponential-local-number-bound} for every \(\beta > 0\). These bounds remove the local particle-number hypothesis from both conclusions whenever the chemical potential satisfies the uniform gap \eqref{eq:fock-path-chemical-potential-gap}.

The chemical-potential gap keeps the system on the vacuum side of the one-particle threshold. The present comparison does not cover higher-density regimes in which condensation or superfluid behavior may occur. For general superstable interactions outside this gap, an all-temperature local insertion estimate remains open. At high temperature, the volume-uniform estimates of \cite[Theorem 1]{GongKuwaharaTong001} provide a complementary route to uniqueness and boundary-independent convergence of the entire Gibbs family. The equilibrium theorem proved here is an existence result and does not assert that all volume exhaustions have the same limit.

A structural obstruction precedes the equilibrium problem. The free Bose dynamics acts by automorphisms of the resolvent algebra, whereas an interacting dynamics generally sends elements of the original gauge-invariant resolvent algebra into a larger algebra. The observable-algebra extension constructed in \cite{DetlevBuchholz002,DetlevBuchholz003} contains these evolved observables for a large class of pair interactions. In the lattice formulation of \cite[Definition 2.2 and Proposition 3.1]{DeuchertLampartLemm001}, this role is played by the concrete Buchholz algebra on bosonic Fock space.

The resolvent algebra itself has an abstract universal \(\oacstar\)-algebraic presentation and admits faithful regular representations \cite{BuchholzGrundling002}. The present construction of the Buchholz algebra reintroduces a distinguished Fock representation at the level of the enlarged algebra. This dependence is restrictive because a von Neumann algebra obtained as the weak closure of a GNS representation retains the folium normal to that representation. Representations associated with different temperatures in an infinite system are frequently disjoint, and their normal states cannot then be represented simultaneously as normal states of one of the corresponding factors. The purpose of the present program is to free the Buchholz algebra from the Fock representation as a reference representation and to retain, together with the resolvent algebra, as much abstract and universal \(\oacstar\)-algebraic information as possible.

The concrete argument in this paper separates the inputs needed for such an abstraction. The dynamics uses the Buchholz sector consistency, the direct cutoff estimate, and finite-range shell propagation. The present equilibrium proof realizes its local equilibrium states by the Fock trace and derives the required local-number estimates from creation-annihilation insertions and a positive functional integral. The Fock trace itself need not be retained as an axiom of an abstract formulation. Proposition 3.1 of \cite{DeuchertLampartLemm001}, recorded in Fact \ref{fact:dll-buchholz-automorphism}, verifies in the concrete model that every finite \(\oa{B}_X\) is invariant under its local dynamics \(\alpha_X\). An abstract formulation may require this invariance together with the existence of a \(\beta\)-KMS state for \(\alpha_X\) on every finite local algebra. Impose in addition the standard local-normality requirement of algebraic quantum statistical mechanics. Here local normality means that the restriction of a state to every finite local algebra extends to a normal state in a regular representation of that finite subsystem. The local KMS assumption and local normality should make the choice of the Fock representation inessential on each finite local algebra, without building one global Fock representation into the abstract algebra. The remaining task is to formulate the positive functional integral and the resulting local-number estimates in a form compatible with these locally normal KMS states.

Continuous-space bosons form a separate extension problem. Buchholz established stability of enlarged gauge-invariant resolvent algebras under continuous pair-potential dynamics and developed the corresponding sector structures \cite{DetlevBuchholz002,DetlevBuchholz003}. Adapting the present strategy requires spatial localization estimates that control both particle transport and ultraviolet behavior. The canonical-ensemble analysis in \cite{DetlevBuchholz012} provides another test of whether the particle-sector mechanism can be separated from one global Fock folium.

Perturbation theory is a further direction. Standard-form methods control Liouvilleans and equilibrium vectors after a representation has been selected \cite{DerezinskiJaksicPillet001}. The positive-temperature Euclidean construction in \cite[Sections 17.1.5 and 21.4--21.5]{DerezinskiGerard001} identifies the Gaussian periodic path-space representation with an Araki--Woods von Neumann algebra and uses this identification to construct perturbed dynamics and KMS states by functional integration. For bosonic Gaussian path spaces, it gives a particularly sharp von Neumann-algebraic realization of the abstract equivalence between stochastically positive KMS systems, periodic Osterwalder--Schrader positive processes, and positive semigroup structures established in \cite[Sections 1 and 6--8]{KleinLandau001}. The Klein--Landau construction is formally based on a \(\oacstar\)-dynamical system, but its stochastic reconstruction passes to the GNS Hilbert space and the associated von Neumann algebra.

An intrinsic theory should determine which perturbations preserve the abstract resolvent or Buchholz algebra and compare the resulting KMS states with standard-form and Araki--Woods perturbations only after a representation has been selected. This requires a deeper \(\oacstar\)-algebraic theory of probability and stochastic processes in which Euclidean reconstruction and functional-integral perturbation do not begin with a fixed Fock or Araki--Woods reference representation.

\bibliography{myref.bib}


\end{document}
